% !TEX root = ../print.tex
% Chapter 5, continued: section 5.6 of the article, groups smaller than GL(d) (sec:smaller-groups).
% In order: the ray column of the (G, W) scheme for every group (lem:gw-kernels-from-origin, which
% until Q-0278 sat in the first section of 02-affine.tex after lem:parameter-space); the frame
% families and their characterization by the diagonal group (the draft's Proposition 2.11), which the
% diagonal cell of the ray column reads; one linear flow (the anisotropy budget, the one-parameter
% row); the similarities with the whole Loewner order (conjecture A', Q-0264, and
% prop:similarity-cascade-families, Q-0265); and the similarities along the stable order
% (prop:similarity-stable-order-families, Q-0267); and the diagonal group along the stable order
% (prop:diagonal-stable-order-families, Q-0282). The plan's outline lists the ray column before
% the frames; the frames come first here because prop:diagonal-ray-families reads
% prop:fixed-frame-families. The draft's Remark 2.17 read in the nodes' terms, which closed this
% part until Q-0284, is in 05-shape-adapted.tex (5.8), whose material it is. Until Q-0278 (2026-10-04) all of these were
% in 02-affine.tex; no label, statement or proof changed in the move. What the nodes spend is in
% 05-landscape.tex's header.

\section{Smaller groups}
\label{sec:smaller-groups}

For a group smaller than the general linear group the landscape has more cells. Along rays a family
is a continuous covariant field of self-decomposable laws (\ref{lem:gw-kernels-from-origin},
\ref{prop:gw-ray-families}), and for the similarities and the positive diagonal matrices the
parameter set is no longer one orbit, so the isotropic member no longer fixes the family
(\ref{prop:similarity-ray-families}, \ref{prop:diagonal-ray-families}). The affine orbit of a product
family cascades along its own frame's axes, and a family covariant under the diagonal matrices that
cascades along each axis is a product of self-decomposable laws on the line
(\ref{prop:frame-families}, \ref{prop:fixed-frame-families}). The stable law's exponents give a bound
on the anisotropy of the exponent of every rotation-invariant non-Gaussian law
(\ref{prop:anisotropy-budget}), and the families of one linear flow drift in shape along a fixed
curve, locked to scale (\ref{prop:one-parameter-families}). For families covariant under similarities
only, a stable part may carry the shape of $\Sigma$ while the family cascades along the whole Loewner
order, but its eccentricity is bounded when the gauge extends to rank one
(\ref{prop:similarity-stable-family}). Without that extension it need not be
(\ref{prop:similarity-degenerate-family}); for a part with no orbit structure the bound holds when the
part's limit at rank one is comparable to the isotropic member, and it is measured in one of three
anisotropy functionals (\ref{prop:similarity-general-sandwich}, \ref{prop:similarity-general-bounded-shape},
\ref{prop:similarity-line-family}); at index $0$ such a part is shape-blind at the origin only
(\ref{prop:similarity-general-index-zero}). For cascade families the infinite divisibility these nodes
assume is a theorem (\ref{prop:similarity-cascade-families}). The section ends with the similarity
families that cascade along the stable order itself rather than the Loewner order: their kernels are
infinitely divisible, and the anisotropy of a kernel is bounded by a power of that of its parameter
(\ref{prop:similarity-stable-order-kernels}).

% not in the draft: the plan's item D5 (notes/PLAN-article-structure.md, the overview table; Q-0267).
% The bookkeeping every cell of the table needs: a family is determined by its kernels from the
% origin once their transforms have no zeros, those kernels are characterized by four properties,
% and a family of a larger datum restricts to a smaller one. What writing it down found: the
% continuity axiom (K2) on pairs off the origin costs nothing beyond continuity of the kernels from the
% origin, because an increment is a quotient of two of them.
% CHANGED Q-0281 wave 1: blind review W1B-3, \uses gains lem:gw-scheme-claims (the proof reads its (1),
% 0 <=_W Sigma, in (2)(c), and its (3), the transform of an image, in (K3)); W1B-1, the annotation's
% sentence on zeros now points to lem:gw-kernels-self-decomposable (the author's decision, 2026-10-04).
% CHANGED Q-0281 wave 1: blind check KS-6, "Not claimed" now names a field of laws, and "This lemma does
% not use it" replaces an ambiguous "The lemma itself does not use this".
\begin{lemma}[a $(G, W)$ family is its kernels from scale zero]
\label{lem:gw-kernels-from-origin}
\uses{def:gw-cascade-family,lem:gw-scheme-claims}
\lean{AffineHemigroup.gw_family_restrict, AffineHemigroup.gw_kernels_from_origin,
AffineHemigroup.gw_kernels_from_origin_unique, AffineHemigroup.gw_kernels_from_origin_converse}\leanok
Let $(G, W)$ be a wedge datum in dimension $d$ (\ref{def:gw-cascade-family}).
\begin{enumerate}
\item \emph{(Restriction.)} Let $(G', W')$ be a wedge datum in dimension $d$ with $G \subseteq G'$ and
with $\Sigma \le_{W'} \Sigma'$ whenever $\Sigma \le_W \Sigma'$. The restriction to $P_W$ of a
$(G', W')$-cascade family of kernels is a $(G, W)$-cascade family of kernels, with the same kernels from scale zero.
\item \emph{(The kernels from scale zero.)} The kernels from scale zero $\mu_\Sigma = \mu_{0,\Sigma}$,
$\Sigma \in \PSD(d)$, of a $(G, W)$-cascade family of kernels satisfy
\begin{enumerate}
\item[(a)] $\mu_0 = \delta_0$ and $\mu_I \ne \delta_0$;
\item[(b)] $A_*\mu_\Sigma = \mu_{A\cdot\Sigma}$ for every $A \in G$ and $\Sigma \in \PSD(d)$;
\item[(c)] for every $\omega \in \RR^d$ the map $\Sigma \mapsto \widehat\mu_\Sigma(\omega)$ is
continuous on $\PSD(d)$;
\item[(d)] for every pair $\Sigma \le_W \Sigma'$ there is a probability measure $\rho$ on $\RR^d$ with
$\mu_{\Sigma'} = \mu_\Sigma * \rho$, namely $\rho = \mu_{\Sigma,\Sigma'}$.
\end{enumerate}
If no $\widehat\mu_\Sigma$ has a zero, $\mu_{\Sigma,\Sigma'}$ is the only $\rho$ in (d), so the family is
determined by its kernels from scale zero.
\item \emph{(The converse.)} Let $(\mu_\Sigma)_{\Sigma \in \PSD(d)}$ be probability measures on $\RR^d$
whose transforms have no zeros and which satisfy (a)--(d). For $\Sigma \le_W \Sigma'$ let
$\mu_{\Sigma,\Sigma'}$ be the probability measure $\rho$ of (d). Then $(\mu_{\Sigma,\Sigma'})$ is a
$(G, W)$-cascade family of kernels, its kernels from scale zero are the $\mu_\Sigma$, and
$\widehat\mu_{\Sigma,\Sigma'} = \widehat\mu_{\Sigma'}/\widehat\mu_\Sigma$.
\end{enumerate}
\statusT\quad\emph{(Cites nothing beyond the uniqueness of the Fourier transform of a probability
measure on $\RR^d$. \textbf{What it is worth.} It reduces a cell of the plan's table to a statement about
one map on $\PSD(d)$: a family whose kernels from scale zero have transforms without zeros is the same
thing as a continuous $G$-covariant field of laws, trivial at $0$ and not at $I$, in which every
$W$-step has a probability measure as increment. For the kernels from scale zero of a family the
transforms never have zeros: they are self-decomposable, for every wedge datum
(\ref{lem:gw-kernels-self-decomposable}), so the hypothesis of the last sentence of (2) always holds
for a family's own kernels, at no cost in the ledger; it matters in (3), for a field of laws not yet
known to come from a family. This lemma does not use it. Clause (1) is the inclusions of the table: a
family of a larger group or a larger order is a family of every smaller datum. \textbf{Not claimed}:
anything for a field of laws whose transforms vanish somewhere, where an increment need not be unique and
(K2) need not follow from (c).)}
\end{lemma}

\begin{proof}
\leanok
\emph{(1).} $P_W \subseteq P_{W'}$, and a triple $\Sigma \le_W \Sigma' \le_W \Sigma''$ is a triple for
$W'$, so (K1) holds on $P_W$; (K2) is the restriction of a continuous map to the subspace $P_W$; (K3)
holds for every $A \in G \subseteq G'$, and $(A\cdot\Sigma, A\cdot\Sigma') \in P_W$ by the
$G$-invariance of $W$; and $(0, I) \in P_W$, so (K4) is the same condition.

\emph{(2).} (a) is (K1) at the pair $(0,0)$ and (K4). (b) is (K3) at the pair $(0, \Sigma)$, as
$A\cdot 0 = 0$. (c): $0 \le_W \Sigma$ for every $\Sigma$, since $W$ contains the rays, and
$\Sigma \mapsto (0, \Sigma)$ is continuous from $\PSD(d)$ into $P_W$; compose with (K2). (d) is (K1) at
$0 \le_W \Sigma \le_W \Sigma'$. If $\mu_{\Sigma'} = \mu_\Sigma * \rho$, then
$\widehat\rho = \widehat\mu_{\Sigma'}/\widehat\mu_\Sigma$ wherever $\widehat\mu_\Sigma \ne 0$, which is
everywhere when $\widehat\mu_\Sigma$ has no zero; so $\rho$ is determined by its transform.

\emph{(3).} By the last sentence, $\mu_{\Sigma,\Sigma'}$ is well defined and its transform is the
quotient $q_{\Sigma,\Sigma'} := \widehat\mu_{\Sigma'}/\widehat\mu_\Sigma$. At $\Sigma = 0$ the quotient is
$\widehat\mu_{\Sigma'}$, by (a), so $\mu_{0,\Sigma'} = \mu_{\Sigma'}$. (K1): $q_{\Sigma,\Sigma} = 1$, and
$q_{\Sigma',\Sigma''}\,q_{\Sigma,\Sigma'} = q_{\Sigma,\Sigma''}$ whenever $\Sigma \le_W \Sigma' \le_W \Sigma''$;
uniqueness of the transform turns these into $\mu_{\Sigma,\Sigma} = \delta_0$ and the cascade identity.
(K2): by (c) the numerator and the denominator of $q_{\Sigma,\Sigma'}(\omega)$ are continuous on
$\PSD(d) \times \PSD(d)$ and the denominator has no zero, so the quotient is continuous there, and in
particular on $P_W$. (K3): for $A \in G$, $(A\cdot\Sigma, A\cdot\Sigma') \in P_W$ by $G$-invariance, and by
(b) the transform of $A_*\mu_{\Sigma,\Sigma'}$ at $\omega$ is
$q_{\Sigma,\Sigma'}(A^{\tp}\omega) = \widehat\mu_{A\cdot\Sigma'}(\omega)/\widehat\mu_{A\cdot\Sigma}(\omega)
= q_{A\cdot\Sigma,A\cdot\Sigma'}(\omega)$. (K4): $\mu_{0,I} = \mu_I \ne \delta_0$.
\end{proof}

% draft: section 2.3, Proposition 2.11. The draft's factors are admissible exponents of [V]; the
% node needs only that each factor is a cascade on the line, and its last clause -- "the frame
% drops out exactly when every factor is Gaussian", which the draft leaves as a check -- is proved
% by the Herschel-Maxwell argument, which is also what the draft's Proposition 3.3 calls on.
% CHANGED Q-0281 wave 2, group E: proved (AffineHemigroup.FrameFamilies, Lean core, guarded) on
% this proof's route; the Lean reads the equality of the i-th marginals in clause (3) off the
% transforms at t e_i, which is how the marginals are compared here. Statement and proof unchanged.
\begin{proposition}[frame families]
\label{prop:frame-families}
\uses{thm:affine-rigidity}
\lean{AffineHemigroup.frame_family_cascade, AffineHemigroup.frame_family_gaussian,
AffineHemigroup.frame_family_only_gaussian}\leanok
Let $d \ge 1$ and, for $i = 1,\dots,d$, let $(\nu^{(i)}_{s,s'})_{0 \le s \le s'}$ be probability
measures on $\RR$ with $\nu^{(i)}_{s,s} = \delta_0$ and
$\nu^{(i)}_{s,s''} = \nu^{(i)}_{s',s''} * \nu^{(i)}_{s,s'}$ for $s \le s' \le s''$. For
$s \le s'$ in $[0,\infty)^d$, coordinatewise, put
$\mu^{\otimes}_{s,s'} := \nu^{(1)}_{s_1,s'_1} \otimes \dots \otimes \nu^{(d)}_{s_d,s'_d}$, and for
$A \in \GL(d)$ put $\mu_{A;s,s'} := A_*\mu^{\otimes}_{s,s'}$, the \emph{frame family}.
\begin{enumerate}
\item \emph{(Cascade in the frame, and covariance.)}
$\widehat\mu_{A;s,s'}(\omega) = \prod_{i}\widehat\nu^{(i)}_{s_i,s'_i}\bigl((A^{\tp}\omega)_i\bigr)$;
$\mu_{A;s,s} = \delta_0$ and $\mu_{A;s,s''} = \mu_{A;s',s''} * \mu_{A;s,s'}$ for
$s \le s' \le s''$; and $B_*\mu_{A;s,s'} = \mu_{BA;s,s'}$ for $B \in \GL(d)$.
\item \emph{(The Gaussian frame family.)} If $b \ge 0$ and $\nu^{(i)}_{s,s'}$ is the centred
Gaussian of variance $2b(s'^2 - s^2)$ for every $i$, then $\mu_{A;0,s}$ is the centred Gaussian of
covariance $2b\,A\diag(s_1^2,\dots,s_d^2)A^{\tp}$, and so depends on $(A,s)$ only through
$A\diag(s^2)A^{\tp}$.
\item \emph{(Only the Gaussian loses its frame.)} Let $d \ge 2$, and suppose $\mu_{A;0,s}$ depends
on $(A,s) \in \GL(d) \times (0,\infty)^d$ only through $A\diag(s^2)A^{\tp}$. Then there is a
$b \ge 0$ with $\nu^{(i)}_{0,t}$ the centred Gaussian of variance $2bt^2$ for every $i$ and every
$t > 0$ (the point mass $\delta_0$ when $b = 0$).
\end{enumerate}
\statusT\quad\emph{(Cites nothing beyond the uniqueness of the Fourier transform. The draft
asks for admissible factors; clauses (1)--(3) read only the cascade on the line. \textbf{Not
claimed}: the parameter count ($d^2$ parameters, a frame modulo its stabilizer and a shape) and
the stabilizer (signed permutations when the factors coincide), both of which the draft's check
item~7 leaves open; nor any statement about which frame families are cascades along the Loewner
order --- \ref{thm:affine-rigidity} already says only the Gaussian is. \textbf{What it is worth.}
Clause (3) is the draft's ``the frame drops out exactly when every factor is Gaussian'', whose
check item it discharges; the argument is the classical one for a rotation-invariant product
law.)}
\end{proposition}

\begin{proof}
\leanok
\emph{Clause (1).} The transform of a product measure is the product of the transforms, so
$\widehat{\mu^\otimes_{s,s'}}(\zeta) = \prod_i\widehat\nu^{(i)}_{s_i,s'_i}(\zeta_i)$, and
$\widehat{A_*\mu}(\omega) = \widehat\mu(A^{\tp}\omega)$ gives the formula. The cascade is
coordinatewise multiplication of transforms, and $(BA)_* = B_*A_*$ gives the covariance; in each
case the transform determines the law.

\emph{Clause (2).} By clause (1),
$\widehat\mu_{A;0,s}(\omega) = \exp\bigl[-b\sum_i s_i^2(A^{\tp}\omega)_i^2\bigr]
= \exp\bigl[-b\langle\omega, A\diag(s^2)A^{\tp}\omega\rangle\bigr]$.

\emph{Clause (3).} Write $\mathbf 1 = (1,\dots,1)$ and
$\varphi_i := \widehat\nu^{(i)}_{0,1}$. For $Q \in \Orth(d)$, $Q\diag(\mathbf 1)Q^{\tp} = I$, so the
hypothesis gives $Q_*\mu^\otimes_{0,\mathbf 1} = \mu_{Q;0,\mathbf 1} = \mu_{I;0,\mathbf 1}
= \mu^\otimes_{0,\mathbf 1}$: the product law is rotation invariant. So
$\psi(\omega) := \prod_i\varphi_i(\omega_i)$ satisfies $\psi(Q^{\tp}\omega) = \psi(\omega)$, and
$\psi(\omega) = f(|\omega|^2)$ with $f(u) := \psi(\sqrt u\,e_1)$ continuous and $f(0) = 1$;
$Q = -I$ makes $\psi$ real, so $f$ is real, and $|f| \le 1$. Since every $\varphi_j(0) = 1$,
$\varphi_i(t) = \psi(te_i) = f(t^2)$. As $d \ge 2$, $\omega = \sqrt x\,e_1 + \sqrt y\,e_2$ gives
$f(x+y) = \varphi_1(\sqrt x)\varphi_2(\sqrt y) = f(x)f(y)$ for $x, y \ge 0$. Then
$f(x) = f(x/2)^2 \ge 0$, and a zero $x_0$ of $f$ would give $f(x_02^{-n}) = 0$ for all $n$, against
continuity at $0$; so $f > 0$, and $G := -\log f$ is continuous, nonnegative and additive on
$[0,\infty)$. By the last paragraph of the proof of \ref{lem:single-ratio-cauchy} (which reads only
additivity and $G \ge 0$), $G(u) = bu$ with $b = G(1) \ge 0$. So
$\varphi_i(t) = e^{-bt^2}$: each $\nu^{(i)}_{0,1}$ is the centred Gaussian of variance $2b$.
For general $s \in (0,\infty)^d$ take $A = \diag(s)^{-1}$: then $A\diag(s^2)A^{\tp} = I$, and the
hypothesis gives $\diag(s)^{-1}_*\mu^\otimes_{0,s} = \mu^\otimes_{0,\mathbf 1}$. Both sides are
product measures and the map acts coordinatewise, so their $i$-th marginals agree: the image of
$\nu^{(i)}_{0,s_i}$ under $x \mapsto x/s_i$ is the centred Gaussian of variance $2b$, that is
$\nu^{(i)}_{0,s_i}$ is the centred Gaussian of variance $2bs_i^2$.
\end{proof}

% not in the draft: the author's question after Q-0206 (see prop:ray-families). The row "the
% diagonal group, steps along one axis": the converse of prop:frame-families(1) in a fixed frame,
% with covariance, placed after it.
% CHANGED Q-0281 wave 2, group E: proved (AffineHemigroup.FixedFrame, Lean core, guarded) on this
% proof's route; the Lean compares the marginals of the last two paragraphs through their transforms
% (the residual of a factor is the image of rho under x -> c x_i). Statement and proof unchanged.
\begin{proposition}[a fixed frame: the families are products]
\label{prop:fixed-frame-families}
\uses{def:isotropic-kernel-family,lem:charfn-one-on-subspace,prop:frame-families}
\lean{AffineHemigroup.fixed_frame_families, AffineHemigroup.fixed_frame_families_converse}\leanok
Let $d \ge 1$. For $s \in (0,\infty)^d$ write $D_s := \diag(s_1,\dots,s_d)$, and $\mathbf 1 :=
(1,\dots,1)$. Let $(\mu_s)_{s\in(0,\infty)^d}$ be probability measures on $\RR^d$ with
\begin{enumerate}
\item[(D1)] \emph{(covariance under the diagonal group)} $(D_\delta)_*\mu_s = \mu_{D_\delta s}$ for
all $\delta, s \in (0,\infty)^d$;
\item[(D2)] \emph{(a cascade along each axis)} for every $i \le d$ and every $\lambda > 1$ there is
a probability measure $\rho_{i,\lambda}$ on $\RR^d$ with
$\mu_{\mathbf 1 + (\lambda - 1)e_i} = \mu_{\mathbf 1} * \rho_{i,\lambda}$.
\end{enumerate}
Then there are self-decomposable laws $\nu_1, \dots, \nu_d$ on $\RR$
(in the sense of \ref{eq:self-decomposable} at $d = 1$) with
\[
  \mu_s = (s_1)_*\nu_1 \otimes \dots \otimes (s_d)_*\nu_d \qquad \text{for every } s \in (0,\infty)^d ,
\]
where $(s_i)_*$ is the image under $x \mapsto s_ix$ on $\RR$. Conversely, for any self-decomposable
laws $\nu_1, \dots, \nu_d$ on $\RR$ this family satisfies (D1) and (D2), and it cascades along
the coordinatewise order: $\mu_{s'} = \mu_s * (\rho^{(1)}_{s_1,s'_1} \otimes \dots \otimes \rho^{(d)}_{s_d,s'_d})$ for $s \le s'$,
with probability measures $\rho^{(i)}_{s_i,s'_i}$ on $\RR$.
\statusT\quad\emph{(Cites nothing beyond the uniqueness of the Fourier transform. The hypotheses
are weak: covariance under the positive diagonal matrices only, no rotation and no reflection;
steps along one axis at a time, at the one point $\mathbf 1$ (by (D1) that is the same as at every
$s$); no continuity in $s$, and no nonvanishing of the transforms, which is why
\ref{lem:charfn-one-on-subspace} is read in its neighbourhood form. \textbf{What it is worth.} It
is the converse of \ref{prop:frame-families}(1) in a fixed frame: the non-Gaussian anisotropic
families that cascade at a fixed orientation are exactly the products of $d$ independent
self-decomposable laws, one per axis, and they may differ; by \ref{prop:frame-families}(3) the
non-Gaussian factors then fix the frame. So a product family has unbounded anisotropy and
cannot turn; the reading that the obstruction to non-Gaussian kernels is turning alone does not
hold in general, since stable-type isotropic orbits do turn, with bounded anisotropy
(\ref{prop:stable-step-cone}, \ref{prop:anisotropy-budget}).
\textbf{Not claimed}: the case of a one-parameter group $e^{sB}$ with $B$ not diagonal, or with the
cascade only along the curve $s \mapsto e^{sB}$, where the laws are operator-self-decomposable and
need not be products (with rotations added, \ref{prop:stable-exponent-cone} and
\ref{prop:anisotropy-budget}); with reflections added to (D1) the factors
are symmetric, which is not stated.)}
\end{proposition}

\begin{proof}
\leanok
\emph{The form of the family.} (D1) with $\delta = s$ at the point $\mathbf 1$ gives
$\mu_s = (D_s)_*\mu_{\mathbf 1}$. Put $\varphi := \widehat\mu_{\mathbf 1}$, continuous with
$\varphi(0) = 1$; then $\widehat\mu_s(\omega) = \varphi(D_s\omega)$.

\emph{A step along axis $i$ is a blur along axis $i$.} Fix $i$ and $\lambda > 1$, write
$\rho := \rho_{i,\lambda}$, $H := e_i^\perp$, and $\omega = xe_i + \eta$ with $x \in \RR$,
$\eta \in H$. By (D2) and the form, $\varphi(\lambda xe_i + \eta) = \varphi(xe_i + \eta)\,
\widehat\rho(xe_i + \eta)$. At $x = 0$ this is $\varphi(\eta) = \varphi(\eta)\widehat\rho(\eta)$,
so $\widehat\rho = 1$ on $\{\eta \in H : \varphi(\eta) \ne 0\}$, a neighbourhood of $0$ in $H$ by
continuity. By \ref{lem:charfn-one-on-subspace}, $\widehat\rho(xe_i + \eta) = \widehat\rho(xe_i)$.
Hence both $\varphi(\lambda xe_i + \eta) = \varphi(xe_i + \eta)\,\widehat\rho(xe_i)$ and
$\varphi(\lambda xe_i) = \varphi(xe_i)\,\widehat\rho(xe_i)$, and so
\[
  \varphi(\lambda xe_i + \eta)\,\varphi(xe_i) = \varphi(xe_i + \eta)\,\varphi(\lambda xe_i)
  \qquad (x \in \RR,\ \eta \in H,\ \lambda > 1).
\]

\emph{Splitting off axis $i$.} Fix $x_0 \ne 0$ and $\eta \in H$. For $x$ strictly between $0$ and
$x_0$, $\lambda := x_0/x > 1$, and the identity reads
$\varphi(x_0e_i + \eta)\,\varphi(xe_i) = \varphi(xe_i + \eta)\,\varphi(x_0e_i)$. Letting $x \to 0$,
continuity and $\varphi(0) = 1$ give $\varphi(x_0e_i + \eta) = \varphi(\eta)\,\varphi(x_0e_i)$,
which holds trivially at $x_0 = 0$. So $\varphi(xe_i + \eta) = \varphi(xe_i)\,\varphi(\eta)$ for
all $x \in \RR$ and $\eta \in e_i^\perp$.

\emph{The product.} Applying this with $i = 1$, then with $i = 2$ to the remainder
$\omega_2e_2 + \dots + \omega_de_d \in e_1^\perp$, and so on,
$\varphi(\omega) = \prod_{i\le d}\varphi(\omega_ie_i)$. Let $\nu_i$ be the image of
$\mu_{\mathbf 1}$ under the $i$-th coordinate, so that $\widehat\nu_i(t) = \varphi(te_i)$; the
product measure $\nu_1 \otimes \dots \otimes \nu_d$ has transform $\prod_i\widehat\nu_i(\omega_i)
= \varphi(\omega)$, and so equals $\mu_{\mathbf 1}$. As $D_s$ acts coordinatewise,
$\mu_s = (D_s)_*\mu_{\mathbf 1} = (s_1)_*\nu_1 \otimes \dots \otimes (s_d)_*\nu_d$.

\emph{Each factor is self-decomposable.} The image of (D2) under the $i$-th coordinate, a linear
map, is $(\lambda)_*\nu_i = \nu_i * \rho'$ with $\rho'$ the image of $\rho_{i,\lambda}$, since
$\mu_{\mathbf 1 + (\lambda-1)e_i}$ has $i$-th marginal $(\lambda)_*\nu_i$ by the product form.
Applying $(\lambda^{-1})_*$ gives $\nu_i = (c)_*\nu_i * (c)_*\rho'$ with $c := \lambda^{-1}$,
which is \ref{eq:self-decomposable} on $\RR$; every $c \in (0,1)$ arises.

\emph{Converse.} (D1): $D_\delta$ acts coordinatewise, so
$(D_\delta)_*\mu_s = (\delta_1s_1)_*\nu_1 \otimes \dots \otimes (\delta_ds_d)_*\nu_d = \mu_{D_\delta s}$. Let $s \le s'$. For
each $i$ with $s_i < s'_i$, self-decomposability at $c := s_i/s'_i$ gives
$\nu_i = (c)_*\nu_i * \rho_c$, and applying $(s'_i)_*$,
$(s'_i)_*\nu_i = (s_i)_*\nu_i * (s'_i)_*\rho_c =: (s_i)_*\nu_i * \rho^{(i)}_{s_i,s'_i}$; put
$\rho^{(i)}_{s_i,s_i} := \delta_0$. A convolution of product measures is the product of the
convolutions of the factors, so $\mu_{s'} = \mu_s * (\rho^{(1)}_{s_1,s'_1} \otimes \dots \otimes \rho^{(d)}_{s_d,s'_d})$, and (D2)
is the case $s = \mathbf 1$, $s' = \mathbf 1 + (\lambda - 1)e_i$.
\end{proof}

% not in the draft: the plan's item D5 (notes/PLAN-article-structure.md, the overview table, column
% "rays"; Q-0267). The ray column of the (G, W) scheme for every group G at once; the next two nodes
% read it for the similarities and for the positive diagonal group. What writing it down found: a law
% with a self-decomposable transform has no zeros by an elementary argument (the first zero along a ray,
% and the doubling inequality for |phi|^2), so neither A4 nor A6 is spent.
% CHANGED Q-0281 wave 1: \uses prop:ray-families dropped, and the proof of (3) rewritten to the
% Lean route: the orbit on PD(d) is (b) at A = Sigma^{1/2} alone, so the proof no longer rests on
% prop:ray-families (cited for comparison only); the radial form of the transform is the reflection
% carrying z to z', as in lem:affine-reduction. A sentence added to (2), necessity: the argument
% holds for every wedge datum (the Lean helper gw_kernels_selfDecomposable). \uses gains
% lem:gw-scheme-claims, which (3) now cites for (GL(d), rays) being a wedge datum.
% CHANGED Q-0281 wave 1: blind review W1B-2, the annotation's sentence on clause (3) reworded (the
% clause no longer reads prop:ray-families). The general form of (2)'s necessity is now stated as
% lem:gw-kernels-self-decomposable, below (W1B-1, the author's decision).
% CHANGED Q-0281 wave 1: blind check KS-7, "the first conclusion of prop:ray-families" (twice, annotation
% and proof of (3)) is its orbit form, which is that node's second conclusion; KS-8, the last sentence of
% (2), necessity, points to lem:gw-kernels-self-decomposable (a \ref, not a \uses).
% CHANGED Q-0327: notation only (whole-paper review flag 14): the profile of (3) and of its proof is
% written $\vartheta$, as in lem:affine-reduction, no longer $F$, which is an exponent in
% eq:admissible-line and thm:isotropic-classification.
% CHANGED writing-rules pass (2026-10-14; records/affine/PASS-writing-rules-affine.md, owed
% item 1): title only, "the ray column: covariant fields of self-decomposable laws" -> the cell of the overview table said in
% words (writing rule 1: no label of the plan in a title). Label, statement and
% \lean unchanged.
% CHANGED round-2 review (Gemini), item 15 (records/affine/reviews/VERIFY-round2-notation-paper.md):
% the frequency on R^d is omega, as the paper's rule asks from its section 3 on (was xi; and z, z_0
% for the zero in the proof of (1), now omega_1, omega_0). Notation only; no step changes.
% CHANGED blind referee of the response (records/affine/REFEREE-round2-response-paper.md, item 8):
% the radial step of the proof of (3) wrote the two frequencies z, z'; now omega and the second
% frequency xi, as the paper's rule asks. Notation only; no step changes.
\begin{proposition}[every group along rays: covariant fields of self-decomposable laws]
\label{prop:gw-ray-families}
\uses{def:gw-cascade-family,lem:gw-scheme-claims,def:isotropic-kernel-family,lem:gw-kernels-from-origin}
\lean{AffineHemigroup.selfDecomposable_charFun_ne_zero, AffineHemigroup.gw_ray_families,
AffineHemigroup.gw_ray_families_unique, AffineHemigroup.gw_ray_families_gl,
AffineHemigroup.gw_ray_families_gl_converse}\leanok
Let $(G, W)$ be a wedge datum in dimension $d$ whose order is the ray order of
\ref{def:gw-cascade-family}: $\Sigma \le_W \Sigma'$ when $\Sigma = 0$ or $\Sigma' = c\,\Sigma$ for some
$c \ge 1$.
\begin{enumerate}
\item \emph{(Self-decomposable laws.)} If $\mu$ is a self-decomposable probability measure on $\RR^d$
\ref{eq:self-decomposable}, then $\widehat\mu$ has no zeros, and $T_*\mu$ is self-decomposable for
every real $d \times d$ matrix $T$.
\item \emph{(The families.)} A map $\Sigma \mapsto \mu_\Sigma$ from $\PSD(d)$ into the probability
measures on $\RR^d$ is the family of kernels from scale zero of a $(G, W)$-cascade family of kernels if
and only if it satisfies (a), (b) and (c) of \ref{lem:gw-kernels-from-origin}(2) and every $\mu_\Sigma$
is self-decomposable. The family is then unique, and its increments are given by
$\widehat\mu_{\Sigma,c\Sigma} = \widehat\mu_{c\Sigma}/\widehat\mu_\Sigma$.
\item \emph{(The general linear group.)} Let $G = \GL(d)$. The kernels from scale zero of a
$(\GL(d), W)$-cascade family of kernels are $\mu_\Sigma = (\Sigma^{1/2})_*\mu_I$ for every
$\Sigma \in \PSD(d)$, where $\mu_I$ is rotation invariant, self-decomposable and not $\delta_0$; with
$\vartheta(u) := \widehat\mu_I(\sqrt u\,e_1)$, $\widehat\mu_\Sigma(\omega) = \vartheta(\langle\omega,\Sigma\omega\rangle)$.
Conversely every such $\mu_I$ is the isotropic member of exactly one $(\GL(d), W)$-cascade family.
\end{enumerate}
\statusT\quad\emph{(Cites nothing beyond the uniqueness of the Fourier transform. \textbf{What it is
worth.} It is the ray column of the plan's table for every group: along rays the scheme asks for nothing
but self-decomposability, the group only says how the kernels at different shapes are related, and
continuity in $\Sigma$ is the one link between shapes that the order does not supply. Clause (3)
recovers the orbit form $\mu_\Sigma = (\Sigma^{1/2})_*\mu_I$ of \ref{prop:ray-families} on $\PD(d)$ from (b) alone, and
adds the boundary of $\PSD(d)$, by continuity, and the uniqueness of the family. The smaller groups
are \ref{prop:similarity-ray-families} and \ref{prop:diagonal-ray-families}. \textbf{Infinite
divisibility.} The ray order has the path condition (\ref{thm:cascade-infinitely-divisible}(3)(b)), so
every kernel of these families is infinitely divisible by entry A6; the node does not need it, since
clause (1) gives the transforms without zeros directly, and entry A4 (a self-decomposable law is
infinitely divisible), which would also give them, is not spent either. \textbf{Not claimed}: anything
about the kernels at different shapes beyond (b) and (c), which for a group smaller than $\GL(d)$ leave
them largely free (the next two nodes).)}
\end{proposition}

\begin{proof}
\leanok
\emph{(1), images.} If $\widehat\mu(\omega) = \widehat\mu(c\omega)\,\widehat\rho_c(\omega)$, then at $T^{\tp}\omega$,
$\widehat{T_*\mu}(\omega) = \widehat\mu(T^{\tp}\omega) = \widehat{T_*\mu}(c\omega)\,\widehat{T_*\rho_c}(\omega)$, which is
\ref{eq:self-decomposable} for $T_*\mu$ with the residual $T_*\rho_c$.

\emph{(1), no zeros.} First an inequality: for a probability measure $\rho$ on $\RR^d$ and $\omega \in \RR^d$,
\[
  1 - |\widehat\rho(2\omega)|^2 \le 4\bigl(1 - |\widehat\rho(\omega)|^2\bigr) .
\]
Indeed $|\widehat\rho(\omega)|^2 = \widehat\rho(\omega)\widehat\rho(-\omega)$ is the transform of
$\kappa := \rho * (-I)_*\rho$, a symmetric probability measure, so
$1 - |\widehat\rho(\omega)|^2 = \int(1 - \cos\langle\omega,x\rangle)\,\kappa(dx)$. As
$1 - \cos 2\tau = 2(1 - \cos\tau)(1 + \cos\tau) \le 4(1 - \cos\tau)$, integrating at
$\tau = \langle\omega,x\rangle$ against $\kappa$ gives the inequality. Now suppose
$\widehat\mu(\omega_1) = 0$ for some $\omega_1$, and let $t_* := \inf\{t > 0 : \widehat\mu(t\omega_1) = 0\}$. As $\widehat\mu$
is continuous with $\widehat\mu(0) = 1$, $t_* > 0$, $\omega_0 := t_*\omega_1$ has $\widehat\mu(\omega_0) = 0$, and
$\widehat\mu(t\omega_0) \ne 0$ for $0 \le t < 1$. For $c \in (0,1)$, \ref{eq:self-decomposable} at $\omega_0$ reads
$0 = \widehat\mu(c\omega_0)\,\widehat\rho_c(\omega_0)$ with $\widehat\mu(c\omega_0) \ne 0$, so $\widehat\rho_c(\omega_0) = 0$; at
$\omega_0/2$ it gives $\widehat\rho_c(\omega_0/2) = \widehat\mu(\omega_0/2)/\widehat\mu(c\omega_0/2)$, which tends to $1$ as
$c \uparrow 1$ by continuity. The inequality at $\omega = \omega_0/2$ gives
$1 \le 4(1 - |\widehat\rho_c(\omega_0/2)|^2) \to 0$, a contradiction.

\emph{(2), necessity.} By \ref{lem:gw-kernels-from-origin}(2) the kernels from scale zero satisfy (a)--(c).
Let $\Sigma \in \PSD(d)$ and $\lambda > 1$. As $\Sigma \le_W \lambda^2\Sigma$, the cascade gives
$\mu_{\lambda^2\Sigma} = \mu_\Sigma * \rho$ with $\rho := \mu_{\Sigma,\lambda^2\Sigma}$, and (b) with
$A = \lambda I \in G$ gives $\mu_{\lambda^2\Sigma} = (\lambda I)_*\mu_\Sigma$. Applying $(\lambda^{-1}I)_*$,
$\mu_\Sigma = (\lambda^{-1}I)_*\mu_\Sigma * (\lambda^{-1}I)_*\rho$, which is \ref{eq:self-decomposable} with
$c = \lambda^{-1}$; every $c \in (0,1)$ arises. The argument reads only that $W$ contains the rays and
$G$ the dilations, so the kernels from scale zero of a cascade family are self-decomposable for every
wedge datum, not only for the ray order (stated as \ref{lem:gw-kernels-self-decomposable}).

\emph{(2), sufficiency.} By (1) no $\widehat\mu_\Sigma$ has a zero, so by
\ref{lem:gw-kernels-from-origin}(3) it suffices to give each pair $\Sigma \le_W \Sigma'$ a probability
measure $\rho$ with $\mu_{\Sigma'} = \mu_\Sigma * \rho$. For $\Sigma = 0$ take $\rho := \mu_{\Sigma'}$,
using $\mu_0 = \delta_0$. For $\Sigma' = c\,\Sigma$ with $c = 1$ take $\delta_0$; with $c > 1$, the
self-decomposability of $\mu_\Sigma$ at $c^{-1/2}$ gives $\mu_\Sigma = (c^{-1/2}I)_*\mu_\Sigma * \rho'$, and
applying $(c^{1/2}I)_*$ and (b) with $A = c^{1/2}I$, $\mu_{c\Sigma} = (c^{1/2}I)_*\mu_\Sigma = \mu_\Sigma *
(c^{1/2}I)_*\rho'$. Uniqueness and the formula for the increments are
\ref{lem:gw-kernels-from-origin}(2)--(3).

\emph{(3).} First, the transform of a rotation-invariant law $\nu$ is radial: if $|\omega| = |\xi|$, the
reflection $Q$ in the hyperplane $(\omega - \xi)^\perp$ is orthogonal and symmetric and carries $\omega$ to $\xi$,
so $\widehat\nu(\xi) = \widehat\nu(Q^{\tp}\omega) = \widehat{Q_*\nu}(\omega) = \widehat\nu(\omega)$. With
$\vartheta(u) := \widehat\nu(\sqrt u\,e_1)$, continuous on $[0,\infty)$, this is $\widehat\nu(\omega) = \vartheta(|\omega|^2)$; and
for $\Sigma \in \PSD(d)$, $\Sigma^{1/2}$ being symmetric,
$\widehat{(\Sigma^{1/2})_*\nu}(\omega) = \widehat\nu(\Sigma^{1/2}\omega) = \vartheta(|\Sigma^{1/2}\omega|^2)
= \vartheta(\langle\omega,\Sigma\omega\rangle)$.

Now let the family be given; $(\GL(d), W)$ is a wedge datum by \ref{lem:gw-scheme-claims}(2). By (2),
$\mu_I$ is self-decomposable, by (a) it is not $\delta_0$, and by (b) with $A = Q \in \Orth(d)$ at
$\Sigma = I$, where $Q \cdot I = I$, it is rotation invariant. Let $\vartheta$ be its function as above. On
$\PD(d)$, $\Sigma^{1/2} \in \GL(d)$ and $\Sigma^{1/2}\cdot I = \Sigma$, so (b) with $A = \Sigma^{1/2}$ gives
$\mu_\Sigma = (\Sigma^{1/2})_*\mu_I$ and $\widehat\mu_\Sigma(\omega) = \vartheta(\langle\omega,\Sigma\omega\rangle)$;
this is the orbit form $\mu_\Sigma = (\Sigma^{1/2})_*\mu_I$ of \ref{prop:ray-families}, obtained here from
(b) alone. For
$\Sigma \in \PSD(d)$, $\Sigma + \varepsilon I \in \PD(d)$ tends to $\Sigma$ as $\varepsilon \downarrow 0$,
and (c) gives $\widehat\mu_\Sigma(\omega) = \lim_{\varepsilon\downarrow0}\vartheta(\langle\omega,(\Sigma +
\varepsilon I)\omega\rangle) = \vartheta(\langle\omega,\Sigma\omega\rangle)$, the transform of
$(\Sigma^{1/2})_*\mu_I$; so $\mu_\Sigma = (\Sigma^{1/2})_*\mu_I$ there too.

Conversely, let $\mu_I$ be rotation invariant, self-decomposable and not $\delta_0$, with function $\vartheta$,
and put $\mu_\Sigma := (\Sigma^{1/2})_*\mu_I$, so that $\widehat\mu_\Sigma(\omega) =
\vartheta(\langle\omega,\Sigma\omega\rangle)$. Then $\mu_I$ is its own member at $\Sigma = I$, as
$\vartheta(|\omega|^2) = \widehat\mu_I(\omega)$, and the conditions of (2) hold.
\begin{itemize}
\item (a), as $\vartheta(0) = 1$.
\item (b) for every $A \in \GL(d)$, as $\widehat{A_*\mu_\Sigma}(\omega) =
\vartheta(\langle A^{\tp}\omega,\Sigma A^{\tp}\omega\rangle) = \vartheta(\langle\omega,A\Sigma A^{\tp}\omega\rangle) =
\widehat\mu_{A\Sigma A^{\tp}}(\omega)$.
\item (c), as $\vartheta$ and $\Sigma \mapsto \langle\omega,\Sigma\omega\rangle$ are continuous.
\item Each $\mu_\Sigma$ is self-decomposable by (1), as an image of $\mu_I$.
\end{itemize} By (2) these are the kernels from scale zero of a family, whose isotropic
member is $\mu_I$. Two families with the isotropic member $\mu_I$ have the kernels from scale zero
$(\Sigma^{1/2})_*\mu_I$, by the first half, and so agree on the pairs by the uniqueness in (2).
\end{proof}

% not in the draft. CHANGED Q-0281 wave 1: a new node, the author's decision (2026-10-04) on the blind
% review's finding W1B-1. The necessity half of prop:gw-ray-families(2) reads only that W contains the
% rays and G the dilations; its Lean (the helper gw_kernels_selfDecomposable, proved and guarded with
% wave 1B, Lean core) states it for every wedge datum, and this node states what the Lean proves.
% CHANGED Q-0281 wave 1: blind check KS-1, the \lean tag gains gw_kernels_hasNoZeros (the clause "its
% transform has no zeros" for a family's kernels, now a declaration of its own) and
% gw_family_determined_by_kernels (the annotation's "determined by its kernels from scale zero, for every
% wedge datum"); KS-4, \uses gains def:isotropic-kernel-family, where eq:self-decomposable is; KS-2, the
% route through infinite divisibility attributed to its entries; KS-3, the example of an increment that is
% not self-decomposable. Statement unchanged.
% CHANGED Q-0281 wave 1: the statement gains its last sentence, the author's decision (2026-10-04) on the
% blind review's finding W1C-1 (a widening; proved, gw_family_determined_by_kernels, Lean core).
\begin{lemma}[the kernels from scale zero are self-decomposable]
\label{lem:gw-kernels-self-decomposable}
\uses{def:gw-cascade-family,def:isotropic-kernel-family,lem:gw-kernels-from-origin,prop:gw-ray-families}
\lean{AffineHemigroup.gw_kernels_selfDecomposable, AffineHemigroup.selfDecomposable_charFun_ne_zero,
AffineHemigroup.gw_kernels_hasNoZeros, AffineHemigroup.gw_family_determined_by_kernels}\leanok
Let $(G, W)$ be a wedge datum in dimension $d$ (\ref{def:gw-cascade-family}). Every kernel from scale zero $\mu_\Sigma = \mu_{0,\Sigma}$, $\Sigma \in \PSD(d)$, of a $(G, W)$-cascade family of kernels is
self-decomposable \ref{eq:self-decomposable}, and its transform has no zeros. In particular two
$(G, W)$-cascade families of kernels with the same kernels from scale zero agree on $P_W$: a family is
determined by its kernels from scale zero.
\statusT\quad\emph{(Cites nothing beyond the uniqueness of the Fourier transform. \textbf{What it is
worth.} The hypothesis of the last sentence of \ref{lem:gw-kernels-from-origin}(2) always holds for a
family's own kernels, so the determination holds for every wedge datum, at no cost in the ledger: neither the route through infinite divisibility (the kernels
are infinitely divisible by \ref{thm:cascade-infinitely-divisible}(2) where the order has the path
condition, entry A6, and an infinitely divisible transform has no zeros by entry A3(i)) nor entry A4 is
needed for it. \textbf{Not claimed}: anything about the increments $\mu_{\Sigma,\Sigma'}$ off the origin,
which need not be self-decomposable --- at $d = 1$, for the $(\GL(1), \text{rays})$ family whose
isotropic member is the Laplace law, $\widehat\mu_\sigma(\xi) = (1 + \sigma\xi^2)^{-1}$, the increment from
$\sigma$ to $\sigma' > \sigma > 0$ has transform $\psi(\xi) = a + (1 - a)(1 + \sigma'\xi^2)^{-1}$,
$a := \sigma/\sigma' \in (0,1)$; were it self-decomposable, a residual $\rho_c$ would have
$\widehat\rho_c(\xi) = \psi(\xi)/\psi(c\xi) \to 1$ as $\xi \to \infty$, averaging $1 - \cos\xi x$ over
$\xi \in [0,T]$ would give $\rho_c = \delta_0$, and then $\psi(\xi) = \psi(c^n\xi) \to 1$, which it is not
---; nor anything about a field of laws not known to come from a family, for which
\ref{lem:gw-kernels-from-origin}(3) still asks that the transforms have no zeros.)}
\end{lemma}

\begin{proof}
\leanok
Let $\Sigma \in \PSD(d)$ and $c \in (0,1)$, and put $\lambda := c^{-1} > 1$. As $W$ contains the rays,
$\Sigma \le_W \lambda^2\Sigma$, and the cascade gives $\mu_{\lambda^2\Sigma} = \mu_\Sigma * \rho$ with
$\rho := \mu_{\Sigma,\lambda^2\Sigma}$; as $G$ contains the dilations, (K3) with $A = \lambda I$ at the
pair $(0, \Sigma)$ gives $\mu_{\lambda^2\Sigma} = (\lambda I)_*\mu_\Sigma$
(\ref{lem:gw-kernels-from-origin}(2)(b)). Applying $(cI)_*$,
$\mu_\Sigma = (cI)_*\mu_\Sigma * (cI)_*\rho$, which is \ref{eq:self-decomposable} at $c$. So $\mu_\Sigma$
% CHANGED Q-0281 wave 1: blind check KS-5 (the parenthesis), KS-1 (the last sentence).
is self-decomposable, and by \ref{prop:gw-ray-families}(1) (clause (1) reads nothing of $(G, W)$) its
transform has no zeros. With the uniqueness in \ref{lem:gw-kernels-from-origin}(2), two $(G, W)$-cascade
families with the same kernels from scale zero therefore agree on $P_W$.
\end{proof}

% not in the draft: the plan's item D5, the cell (similarities, rays) (Q-0267). The parameter set of a
% group smaller than GL(d) is not one orbit; the plan's section 4.2 discusses the stabilizer.
% CHANGED Q-0281: \uses re-pointed to the split child it reads (lem:parameter-space(5) only).
% CHANGED Q-0281 wave 2, group E: proved (AffineHemigroup.SimilarityRays, with the toolbox
% AffineHemigroup.SimilaritySpectrum; Lean core, guarded), and the proof rewritten to the Lean route
% in three places: the closedness of S(d) (through the closed condition AA^T = (AA^T)_11 I, as
% simSet_isClosedSubgroupGL), the simplex Delta_d (its vertices and their independence, which the
% proof had asserted), and the two families of (3) differing (through Stab(mu_I) = O(d),
% lem:parameter-space-examples(5), in place of the eigenvector argument). Statement unchanged.
% CHANGED writing-rules pass (2026-10-14; records/affine/PASS-writing-rules-affine.md, owed
% item 1): title only, "(similarities, rays): the families, and the parameter set" -> the cell of the overview table said in
% words (writing rule 1: no label of the plan in a title). Label, statement and
% \lean unchanged.
\begin{proposition}[similarity covariance along rays: the families, and the parameter set]
\label{prop:similarity-ray-families}
% CHANGED Q-0281 wave 2: blind review W2E-4: the proof of (3) reads lem:parameter-space(2) (the
% Lean's isFull_of_isRotationInvariant), which needs only that mu_I is rotation invariant and not
% delta_0, not a cascade family: \uses lem:parameter-space in place of lem:parameter-space-examples.
\uses{lem:gw-kernels-from-origin,prop:gw-ray-families,lem:affine-gauge,lem:parameter-space}
\lean{AffineHemigroup.similarity_ray_parameter, AffineHemigroup.similarity_ray_families,
AffineHemigroup.similarity_ray_kernels, AffineHemigroup.similarity_ray_isotropic,
AffineHemigroup.similarity_ray_two_families, AffineHemigroup.similarity_row_isotropic_members}\leanok
Let $d \ge 1$, let $\mathrm S(d) := \{cQ : c > 0,\ Q \in \Orth(d)\}$ be the similarities, and let $W$ be
the ray order. For $\Sigma \in \PSD(d) \setminus \{0\}$ write $\sigma(\Sigma) \in \RR^d$ for its
eigenvalues in nonincreasing order divided by $\tr\Sigma$, and put
$\Delta_d := \{\sigma \in \RR^d : \sigma_1 \ge \dots \ge \sigma_d \ge 0,\ \sigma_1 + \dots + \sigma_d = 1\}$.
\begin{enumerate}
\item \emph{(The parameter set.)} $(\mathrm S(d), W)$ is a wedge datum. Two matrices
$\Sigma, \Sigma' \in \PSD(d) \setminus \{0\}$ lie in one orbit of $\mathrm S(d)$, acting by
$A\cdot\Sigma = A\Sigma A^{\tp}$, exactly when $\sigma(\Sigma) = \sigma(\Sigma')$, and each orbit contains
exactly one matrix $\diag(\sigma)$ with $\sigma \in \Delta_d$; so the orbits in $\PSD(d) \setminus \{0\}$
are indexed by $\Delta_d$, a simplex of dimension $d - 1$. The stabilizer of $\Sigma \ne 0$ in
$\mathrm S(d)$ is $\{Q \in \Orth(d) : Q\Sigma = \Sigma Q\}$, and it is $\Orth(d)$ exactly when $\Sigma$ is
a multiple of $I$.
\item \emph{(The families.)} The kernels from scale zero of the $(\mathrm S(d), W)$-cascade families of
kernels are the maps of \ref{prop:gw-ray-families}(2) with $G = \mathrm S(d)$. In such a family every
$\mu_\Sigma$ is symmetric, $(-I)_*\mu_\Sigma = \mu_\Sigma$, and invariant under the stabilizer of
$\Sigma$, and $\mu_{cQ\diag(\sigma)Q^{\tp}} = (\sqrt c\,Q)_*\mu_{\diag(\sigma)}$ for $c > 0$, $Q \in \Orth(d)$
and $\sigma \in \Delta_d$: the family is determined by its kernels at the matrices $\diag(\sigma)$,
$\sigma \in \Delta_d$.
\item \emph{(The isotropic member does not determine the family.)} In such a family
$\mu_{sI} = (\sqrt s\,I)_*\mu_I$ for $s \ge 0$, and $\mu_I$ is rotation invariant, self-decomposable and
not $\delta_0$. Conversely, let $\mu_I$ be such a law. Then (i) $\Sigma \mapsto (\Sigma^{1/2})_*\mu_I$ and
(ii) $\Sigma \mapsto \nu_\Sigma := ((\det\Sigma)^{1/(2d)}I)_*\mu_I$ are the kernels from scale zero of
$(\mathrm S(d), W)$-cascade families with isotropic member $\mu_I$, and (ii) is also the family of kernels
from scale zero of an $(\mathrm S(d), \text{Loewner})$-cascade family. For $d \ge 2$ the two differ at
every $\Sigma \in \PD(d)$ that is not a multiple of $I$.
\item \emph{(The isotropic members of the similarity row.)} Let $W'$ be an $\mathrm S(d)$-invariant
partial order on $\PSD(d)$ that contains the rays and is contained in the Loewner order. The isotropic
members $\mu_I$ of the $(\mathrm S(d), W')$-cascade families of kernels are exactly the rotation-invariant
self-decomposable laws other than $\delta_0$.
\end{enumerate}
\statusT\quad\emph{(Cites nothing beyond the uniqueness of the Fourier transform and the spectral
theorem. \textbf{The parameter set.} \ref{lem:parameter-space}(5) gives, for a group containing
$\Orth(d)$, a full isotropic member with stabilizer $\Orth(d)$; with $\GL(d)$ that member fixes the whole
family (\ref{prop:gw-ray-families}(3)), and the parameter space is one orbit, $\PD(d)$. With the
similarities the orbits are the rays through the shapes $\diag(\sigma)$, $\sigma \in \Delta_d$, so
$\PSD(d) \setminus \{0\}$ is a $(d-1)$-parameter set of orbits, and the kernel at a shape is
constrained pointwise only by its stabilizer (besides self-decomposability; continuity constrains it
where eigenvalues meet): for distinct eigenvalues that is the group of $2^d$ sign changes in an
eigenbasis, and the kernel need not be elliptical (at $\Sigma = \diag(1,2)$, $d = 2$, the kernel
$N(0, 2b(\Sigma + \kappa\tr(\Sigma)I)) * S_\alpha(\Sigma)$ of \ref{prop:similarity-stable-order-kernels}(3),
$b, \kappa > 0$, is not: ellipticity would force the shape $\Sigma$ from the stable factor and then
$\kappa = 0$ from the Gaussian one). \textbf{What it is worth.} The cell
(similarities, rays) of the plan's table is not empty and is large: a continuous field of symmetric
self-decomposable laws over the shapes, covariant under similarities. The isotropic member is any
rotation-invariant self-decomposable law other than $\delta_0$, as in the $\GL(d)$ row, but for
$d \ge 2$ it no longer determines the family; clause (3) gives two families for each member, the $\GL(d)$ orbit and the shape-blind family, and
clause (4) shows that the isotropic members are the same in every column of the similarity row, the
Loewner column included, unlike the $\GL(d)$ row, where the Loewner order leaves the Gaussian alone
(\ref{thm:affine-rigidity}). The constraints a larger order puts on the shapes are
\ref{prop:similarity-stable-order-kernels} and the A$'$ nodes. \textbf{Not claimed}: that a choice of
kernels at the matrices $\diag(\sigma)$, invariant under their stabilizers and continuous in $\sigma$,
always defines a family: continuity in $\Sigma$ where eigenvalues meet is a condition on the choice that
is not worked out.)}
\end{proposition}

\begin{proof}
\leanok
\emph{(1).} $\mathrm S(d)$ contains the dilations, and it is closed in $\GL(d)$: every $A \in \mathrm S(d)$
satisfies $AA^{\tp} = (AA^{\tp})_{11}\,I$, a closed condition, and an invertible $A$ satisfying it has
$k := (AA^{\tp})_{11} > 0$ and $A = \sqrt k\,(k^{-1/2}A)$ with $k^{-1/2}A$ orthogonal. The ray order is invariant under every subgroup of $\GL(d)$
and contains the rays (\ref{def:gw-cascade-family}). For $A = \sqrt c\,Q$, $A\cdot\Sigma = cQ\Sigma Q^{\tp}$
has the eigenvalues of $\Sigma$ multiplied by $c$ and the trace $c\tr\Sigma$, so $\sigma$ is constant on
orbits. If $\sigma(\Sigma) = \sigma(\Sigma') = \sigma$, the spectral theorem gives $Q, Q' \in \Orth(d)$ with
$\Sigma = (\tr\Sigma)\,Q\diag(\sigma)Q^{\tp}$ and $\Sigma' = (\tr\Sigma')\,Q'\diag(\sigma)Q'^{\tp}$, and then
$\Sigma' = A\cdot\Sigma$ with $A := (\tr\Sigma'/\tr\Sigma)^{1/2}\,Q'Q^{\tp}$. For $\sigma \in \Delta_d$,
$\sigma(\diag(\sigma)) = \sigma$, so the orbit of $\Sigma$ contains $\diag(\sigma(\Sigma))$ and no other
$\diag(\sigma')$ with $\sigma' \in \Delta_d$. The set $\Delta_d$ is the convex hull of the $d$ points
$v_k := k^{-1}(e_1 + \dots + e_k)$, $1 \le k \le d$: they lie in the convex set $\Delta_d$, and every
$\sigma \in \Delta_d$ is $\sum_kw_kv_k$ with $w_k := k(\sigma_k - \sigma_{k+1}) \ge 0$ ($\sigma_{d+1} := 0$),
whose sum is $\sum_k\sigma_k = 1$ by summation by parts. The coordinates of $\sum_kg_kv_k$ are
$\sum_{k \ge i}g_k/k$, whose successive differences are the $g_i/i$, so the $v_k$ are linearly, hence
affinely, independent, and $\Delta_d$ is a simplex of dimension $d - 1$. If $(cQ)\cdot\Sigma = c^2Q\Sigma Q^{\tp} = \Sigma$, the traces
give $c = 1$, and $Q\Sigma Q^{\tp} = \Sigma$ is $Q\Sigma = \Sigma Q$. A multiple of $I$ commutes with every
$Q$; conversely, if every $Q \in \Orth(d)$ commutes with $\Sigma$, the reflections and permutations of the
proof of \ref{lem:affine-gauge} make $\Sigma$ a multiple of $I$.

\emph{(2).} The first sentence is \ref{prop:gw-ray-families}(2). Property (b) of
\ref{lem:gw-kernels-from-origin}(2) with $A = -I$, which has $(-I)\cdot\Sigma = \Sigma$, gives the
symmetry; with $A$ in the stabilizer it gives the invariance; with $A = \sqrt c\,Q$ and $\diag(\sigma)$ in
place of $\Sigma$ it gives the formula, and by (1) every $\Sigma \ne 0$ is of the form
$cQ\diag(\sigma)Q^{\tp}$, while $\mu_0 = \delta_0$.

\emph{(3).} (b) with $A = \sqrt s\,I$ gives $\mu_{sI} = (\sqrt s\,I)_*\mu_I$ for $s > 0$, and (a) gives it at
$s = 0$; $\mu_I$ is invariant under the stabilizer $\Orth(d)$ of $I$ by (2), self-decomposable by
\ref{prop:gw-ray-families}(2) and not $\delta_0$ by (a). (i): \ref{prop:gw-ray-families}(3) gives a
$(\GL(d), W)$-cascade family with these kernels from scale zero, and \ref{lem:gw-kernels-from-origin}(1)
restricts it to $(\mathrm S(d), W)$. (ii): the Loewner order is a partial order on $\PSD(d)$, invariant under
$\GL(d)$ and containing the rays (\ref{def:gw-cascade-family}), so $(\mathrm S(d), \text{Loewner})$ is a
wedge datum, and the ray order is contained in it, since $\Sigma' - 0$ and $c\Sigma - \Sigma$, $c \ge 1$,
are positive semidefinite. So by \ref{lem:gw-kernels-from-origin}(3) and (1) it suffices to check (a)--(d)
of \ref{lem:gw-kernels-from-origin}(2) for the Loewner order, with transforms without zeros. (a):
$\det 0 = 0$, so $\nu_0 = \delta_0$, and $\nu_I = \mu_I \ne \delta_0$. (b): for $A = \sqrt c\,Q$,
$\det(cQ\Sigma Q^{\tp})^{1/(2d)} = \sqrt c\,(\det\Sigma)^{1/(2d)}$, and $Q_*\mu_I = \mu_I$ gives
$A_*\nu_\Sigma = (\sqrt c\,(\det\Sigma)^{1/(2d)}I)_*Q_*\mu_I = \nu_{A\cdot\Sigma}$. (c):
$\widehat\nu_\Sigma(\omega) = \widehat\mu_I((\det\Sigma)^{1/(2d)}\omega)$ is continuous in $\Sigma$. No
zeros: $\nu_\Sigma$ is self-decomposable by \ref{prop:gw-ray-families}(1). (d): let $\Sigma \le \Sigma'$.
Then $\det\Sigma \le \det\Sigma'$: trivially if $\Sigma$ is singular, and otherwise
$\det\Sigma' = \det\Sigma\det M$ with $M := \Sigma^{-1/2}\Sigma'\Sigma^{-1/2} \ge I$, whose eigenvalues are
at least $1$. Put $\eta := (\det\Sigma)^{1/(2d)} \le \eta' := (\det\Sigma')^{1/(2d)}$. If
$\eta = 0$, take $\rho := \nu_{\Sigma'}$; if $\eta = \eta'$, take $\delta_0$; if
$0 < \eta < \eta'$, self-decomposability at $c := \eta/\eta'$ gives
$\mu_I = (cI)_*\mu_I * \rho_c$, and applying $(\eta'I)_*$, $\nu_{\Sigma'} = \nu_\Sigma * (\eta'I)_*\rho_c$.

For $d \ge 2$ let $\Sigma \in \PD(d)$ be no multiple of $I$, put $t := (\det\Sigma)^{1/(2d)} > 0$, and
suppose the two kernels agree at $\Sigma$, $(\Sigma^{1/2})_*\mu_I = (tI)_*\mu_I$. Then
$P := t^{-1}\Sigma^{1/2}$ is invertible and $P_*\mu_I = \mu_I$. As $\mu_I$ is rotation invariant and not
$\delta_0$, \ref{lem:parameter-space}(2) gives $\operatorname{Stab}(\mu_I) = \Orth(d)$, so $P$ is orthogonal:
$t^{-2}\Sigma = PP^{\tp} = I$, and $\Sigma = t^2I$, which is excluded.

\emph{(4).} $W'$ contains the rays, so every pair of the ray order is a pair of $W'$ ($0 \le_{W'} \Sigma'$
and $\Sigma \le_{W'} c\Sigma$ for $c \ge 1$), and \ref{lem:gw-kernels-from-origin}(1) restricts an
$(\mathrm S(d), W')$-family to an $(\mathrm S(d), W)$-family with the same isotropic member; (3) applies.
Conversely, for a rotation-invariant self-decomposable $\mu_I \ne \delta_0$, the
$(\mathrm S(d), \text{Loewner})$-family of (3)(ii) restricts to $(\mathrm S(d), W')$, as $W'$ is contained
in the Loewner order.
\end{proof}

% not in the draft: the plan's item D5, the cell (diagonal, rays) (Q-0267). The diagonal group is the
% positive one, as in prop:fixed-frame-families. What writing it down found: the product form of that
% node comes from the axis steps of the Loewner order, not from the group; with rays only, the frame's
% own shapes carry any self-decomposable law.
% CHANGED Q-0281 wave 2, group E: proved (AffineHemigroup.DiagonalRays, Lean core, guarded). Proof
% rewritten to the Lean route in (1) (the openness of the set U through the compact unit sphere,
% PD(d) not being open among all matrices; the dimension d(d-1)/2, which the proof had asserted) and
% in (5) (through Stab(mu_I) = O(d), lem:parameter-space-examples(5), in place of the argument on F
% that cited the old proof of prop:similarity-ray-families(3)); its uses list trades
% prop:similarity-ray-families for lem:parameter-space-examples accordingly. Statement unchanged.
% CHANGED Q-0327: notation only (whole-paper review flag 15): the correlation matrix is $\Theta$, no
% longer $R$, as in prop:diagonal-stable-order-families(4).
% CHANGED writing-rules pass (2026-10-14; records/affine/PASS-writing-rules-affine.md, owed
% item 1): title only, "(diagonal, rays): the families, and the parameter set" -> the cell of the overview table said in
% words (writing rule 1: no label of the plan in a title). Label, statement and
% \lean unchanged.
\begin{proposition}[diagonal covariance along rays: the families, and the parameter set]
\label{prop:diagonal-ray-families}
% CHANGED Q-0281 wave 2: blind review W2E-4: the proof of (5) reads lem:parameter-space(2) (the
% Lean's isFull_of_isRotationInvariant), which needs only that mu_I is rotation invariant and not
% delta_0: \uses lem:parameter-space in place of lem:parameter-space-examples.
\uses{lem:gw-kernels-from-origin,prop:gw-ray-families,prop:fixed-frame-families,lem:parameter-space}
\lean{AffineHemigroup.diagonal_ray_parameter, AffineHemigroup.diagonal_ray_families,
AffineHemigroup.diagonal_ray_frame, AffineHemigroup.diagonal_loewner_products,
AffineHemigroup.diagonal_rays_no_product, AffineHemigroup.diagonal_rays_off_frame}\leanok
Let $d \ge 1$, let $\mathrm D_+(d)$ be the group of diagonal matrices with positive diagonal entries, and
let $W$ be the ray order. For $\Sigma \in \PSD(d)$ write $\Delta(\Sigma) := \diag(\Sigma_{11}, \dots,
\Sigma_{dd})$, for $s \in (0,\infty)^d$ write $D_s := \diag(s_1, \dots, s_d)$, and put
$\mathrm{Corr}(d) := \{\Theta \in \PD(d) : \Theta_{11} = \dots = \Theta_{dd} = 1\}$.
\begin{enumerate}
\item \emph{(The parameter set.)} $(\mathrm D_+(d), W)$ is a wedge datum. Every $\Sigma \in \PD(d)$ has
the trivial stabilizer in $\mathrm D_+(d)$, and is $D\cdot \Theta = D\Theta D$ for exactly one pair
$D \in \mathrm D_+(d)$, $\Theta \in \mathrm{Corr}(d)$, namely $D = \Delta(\Sigma)^{1/2}$. So the orbits in
$\PD(d)$ are indexed by $\mathrm{Corr}(d)$, a nonempty open subset of the affine space of symmetric
matrices with unit diagonal, of dimension $d(d-1)/2$; the orbit of $I$ is the set of positive definite
diagonal matrices.
\item \emph{(The families on the frame's shapes.)} The kernels from scale zero of the
$(\mathrm D_+(d), W)$-cascade families of kernels are the maps of \ref{prop:gw-ray-families}(2) with
$G = \mathrm D_+(d)$. On the orbit of $I$ they are $\mu_{D_s^2} = (D_s)_*\mu_I$, $s \in (0,\infty)^d$, with
$\mu_I$ self-decomposable and not $\delta_0$. Conversely, for every self-decomposable $\mu_I \ne \delta_0$
the map $\Sigma \mapsto (\Delta(\Sigma)^{1/2})_*\mu_I$ is the family of kernels from scale zero of a
$(\mathrm D_+(d), W)$-cascade family.
\item \emph{(The Loewner order: products.)} The kernels from scale zero of a
$(\mathrm D_+(d), \text{Loewner})$-cascade family of kernels are, on the orbit of $I$,
\[
  \mu_{D_s^2} = (s_1)_*\nu_1 \otimes \dots \otimes (s_d)_*\nu_d , \qquad s \in (0,\infty)^d ,
\]
with self-decomposable laws $\nu_1, \dots, \nu_d$ on $\RR$, not all $\delta_0$. Conversely, for such laws
the map $\Sigma \mapsto (\Delta(\Sigma)^{1/2})_*(\nu_1 \otimes \dots \otimes \nu_d)$ is the family of kernels
from scale zero of a $(\mathrm D_+(d), \text{Loewner})$-cascade family.
\item \emph{(Rays only: no product.)} Let $d \ge 2$ and $\Theta \in \mathrm{Corr}(d)$ not diagonal. The centred
Gaussian law $N(0,\Theta)$ is self-decomposable and is not of the form $\nu_1 \otimes \dots \otimes \nu_d$. So
the $(\mathrm D_+(d), W)$-family of (2) with $\mu_I = N(0,\Theta)$ is the restriction of no
$(\mathrm D_+(d), \text{Loewner})$-family.
\item \emph{(Off the frame's shapes.)} Let $d \ge 2$ and $\mu_I$ rotation invariant, self-decomposable and
not $\delta_0$. The maps $\Sigma \mapsto (\Delta(\Sigma)^{1/2})_*\mu_I$ and
$\Sigma \mapsto (\Sigma^{1/2})_*\mu_I$ are the kernels from scale zero of two
$(\mathrm D_+(d), W)$-cascade families with the same isotropic member, and they differ at every
$\Sigma \in \PD(d)$ that is not diagonal.
\end{enumerate}
\statusT\quad\emph{(Cites nothing beyond the uniqueness of the Fourier transform. \textbf{The parameter
set.} The group has dimension $d$ and $\PD(d)$ has dimension $d(d+1)/2$, so for $d \ge 2$ the parameter
set is not one orbit: an orbit is a frame's own shapes, the axis-aligned ellipsoids of one correlation
matrix, and the stabilizer of $I$ in $\mathrm D_+(d)$ is trivial, so $\mu_I$ need carry no symmetry.
\textbf{What it is worth.} The cell (diagonal, rays) of the plan's table is not empty. On the orbit of
$I$ it is \ref{prop:ray-families} without rotations: the frame images of one self-decomposable law, any
such law. Clauses (3) and (4) locate the product form of \ref{prop:fixed-frame-families}: it is forced by
the steps along single axes, which the Loewner order has and the ray order lacks, and not by the group;
on the orbit of $I$, the Loewner cell of the diagonal row is that node read in the scheme. Off the orbit
of $I$ the ray order links nothing across the correlation matrices, and the isotropic member need not
determine the family there (for a rotation-invariant member, clause (5)). \textbf{Not claimed}: a description of the continuous covariant fields over
$\mathrm{Corr}(d)$; the group of all invertible diagonal matrices, which would add invariance of $\mu_I$
under the coordinate reflections; and anything about $(\mathrm D_+(d), \text{Loewner})$-families off the
orbit of $I$ beyond the example in (3).)}
\end{proposition}

\begin{proof}
\leanok
\emph{(1).} A limit in $\GL(d)$ of matrices of $\mathrm D_+(d)$ is diagonal with nonnegative entries and
invertible, so its entries are positive: $\mathrm D_+(d)$ is closed, and it contains the dilations; the ray
order is invariant under every subgroup and contains the rays. For $D = D_s$, $(D\Sigma D)_{ii} =
s_i^2\Sigma_{ii}$, and $\Sigma_{ii} > 0$ on $\PD(d)$, so $D\cdot\Sigma = \Sigma$ forces $D = I$. The matrix
$\Theta := \Delta(\Sigma)^{-1/2}\Sigma\Delta(\Sigma)^{-1/2}$ is positive definite with unit diagonal and
$\Delta(\Sigma)^{1/2}\cdot \Theta = \Sigma$; if $D\Theta D = D'\Theta'D'$ with $\Theta, \Theta' \in \mathrm{Corr}(d)$, the diagonals
give $D^2 = D'^2$, so $D = D'$ and $\Theta = \Theta'$. $\mathrm{Corr}(d)$ is the intersection of that affine space
with the set $U$ of matrices whose quadratic form is positive on the unit sphere, and contains $I$; $U$ is
open, since the unit sphere is compact and $(M, x) \mapsto \langle x, Mx\rangle$ is continuous. The
direction of the affine space, the symmetric matrices with zero diagonal, is parametrized linearly by the
$d(d-1)/2$ entries above the diagonal. The orbit of $I$ is $\{D^2 : D \in \mathrm D_+(d)\}$.

\emph{(2).} The first sentence is \ref{prop:gw-ray-families}(2); (b) of
\ref{lem:gw-kernels-from-origin}(2) with $A = D_s$ at $\Sigma = I$ gives $\mu_{D_s^2} = (D_s)_*\mu_I$, and
$\mu_I$ is self-decomposable and not $\delta_0$. Conversely, put $\mu_\Sigma :=
(\Delta(\Sigma)^{1/2})_*\mu_I$. (a): $\Delta(0) = 0$ and $\Delta(I) = I$. (b): $\Delta(D\Sigma D) =
D^2\Delta(\Sigma)$, and diagonal matrices commute, so $\Delta(D\cdot\Sigma)^{1/2} = D\,\Delta(\Sigma)^{1/2}$ and
$\mu_{D\cdot\Sigma} = D_*\mu_\Sigma$. (c): $\widehat\mu_\Sigma(\omega) = \widehat\mu_I(\Delta(\Sigma)^{1/2}\omega)$
is continuous in $\Sigma$. Each $\mu_\Sigma$ is self-decomposable by \ref{prop:gw-ray-families}(1), and
\ref{prop:gw-ray-families}(2) applies.

\emph{(3), necessity.} By \ref{lem:gw-kernels-from-origin}(2), $\mu_s := \mu_{D_s^2}$,
$s \in (0,\infty)^d$, satisfies (D1) of \ref{prop:fixed-frame-families}, as
$D_\delta\cdot D_s^2 = D_{\delta s}^2$, and (D2), as $D^2_{\mathbf 1 + (\lambda-1)e_i} - I$ is diagonal with
the one nonzero entry $\lambda^2 - 1 > 0$, so that $I \le D^2_{\mathbf 1 + (\lambda-1)e_i}$ is a Loewner
step with the increment $\mu_{I, D^2_{\mathbf 1 + (\lambda-1)e_i}}$. That node gives the product form with
self-decomposable $\nu_i$, and they are not all $\delta_0$ since $\mu_I \ne \delta_0$.

\emph{(3), sufficiency.} $(\mathrm D_+(d), \text{Loewner})$ is a wedge datum: the Loewner order is
invariant under every invertible matrix and contains the rays. Put $\nu := \nu_1 \otimes \dots \otimes \nu_d$ and $\mu_\Sigma :=
(\Delta(\Sigma)^{1/2})_*\nu$; as $\Delta(\Sigma)^{1/2}$ acts coordinatewise,
$\mu_\Sigma = (\Sigma_{11}^{1/2})_*\nu_1 \otimes \dots \otimes (\Sigma_{dd}^{1/2})_*\nu_d$. (a)--(c) hold as in
(2), with $\mu_I = \nu \ne \delta_0$, and the transform $\prod_i\widehat\nu_i(\Sigma_{ii}^{1/2}\omega_i)$ has
no zeros by \ref{prop:gw-ray-families}(1) at $d = 1$. Let $\Sigma \le \Sigma'$ in the Loewner order; then
$a_i := \Sigma_{ii}^{1/2} \le a'_i := \Sigma'^{1/2}_{ii}$ for each $i$. Put $\rho_i := \delta_0$ if
$a_i = a'_i$, $\rho_i := (a'_i)_*\nu_i$ if $0 = a_i < a'_i$, and, if $0 < a_i < a'_i$,
$\rho_i := (a'_i)_*\rho^{(i)}$ with $\nu_i = (a_i/a'_i)_*\nu_i * \rho^{(i)}$ by self-decomposability. Then
$(a'_i)_*\nu_i = (a_i)_*\nu_i * \rho_i$ in each case, and a convolution of product measures is the product of
the convolutions, so $\mu_{\Sigma'} = \mu_\Sigma * (\rho_1 \otimes \dots \otimes \rho_d)$.
\ref{lem:gw-kernels-from-origin}(3) gives the family.

\emph{(4).} $\widehat{N(0,\Theta)}(\xi) = e^{-\langle\xi,\Theta\xi\rangle/2}$ is
$e^{-c^2\langle\xi,\Theta\xi\rangle/2}\,e^{-(1-c^2)\langle\xi,\Theta\xi\rangle/2}$, the second factor the transform of
$N(0,(1-c^2)\Theta)$, which is \ref{eq:self-decomposable}. If $N(0,\Theta) = \nu_1 \otimes \dots \otimes \nu_d$, then
$\widehat\nu_i(t) = \widehat{N(0,\Theta)}(te_i) = e^{-t^2/2}$ as $\Theta_{ii} = 1$, so
$e^{-\langle\xi,\Theta\xi\rangle/2} = e^{-|\xi|^2/2}$ for every $\xi$, and $\Theta = I$, which is diagonal. If the
family of (2) with $\mu_I = N(0,\Theta)$ were the restriction of a $(\mathrm D_+(d), \text{Loewner})$-family, the
two would have the same kernels from scale zero (\ref{lem:gw-kernels-from-origin}(1)), and (3) at $s =
\mathbf 1$ would make $N(0,\Theta)$ a product.

\emph{(5).} The first map is (2); the second is \ref{prop:gw-ray-families}(3) restricted to
$\mathrm D_+(d) \subseteq \GL(d)$ by \ref{lem:gw-kernels-from-origin}(1). Both have $\mu_I$ at $I$. Let
$\Sigma \in \PD(d)$ be not diagonal, put $D := \Delta(\Sigma)^{1/2}$, invertible as every $\Sigma_{ii} > 0$,
and suppose the two kernels agree at $\Sigma$, $D_*\mu_I = (\Sigma^{1/2})_*\mu_I$. Then
$P := D^{-1}\Sigma^{1/2}$ is invertible and $P_*\mu_I = \mu_I$. As $\mu_I$ is rotation invariant and not
$\delta_0$, \ref{lem:parameter-space}(2) gives $\operatorname{Stab}(\mu_I) = \Orth(d)$, so $P$ is orthogonal: $D^{-1}\Sigma D^{-1} = PP^{\tp} = I$, and
$\Sigma = D^2 = \Delta(\Sigma)$ is diagonal, which is excluded.
\end{proof}

% not in the draft: P-0019 (ii), Q-0241. The anisotropy budget. Sufficiency is prop:stable-exponent-cone. Necessity
% spends A5 (Jurek-Mason Thm. 3.3.5 and section 3.4's opening; Urbanik 1972 Prop. 5.2) for infinite
% divisibility of the law and of its residuals, A3 for the Levy measures, and smooths the Levy
% measure radially so that the infinitesimal argument of the proposal applies to every
% rotation-invariant Levy measure, with or without a density. Jurek-Mason Prop. 3.4.1(b), p. 90,
% states the inequality of Levy measures; it is derived here from A3 rather than cited.
% CHANGED Q-0281 wave 2, group E: the proof of clause (2) is rewritten to the Lean route of the child
% prop:anisotropy-budget-converse (proved): the contraction by the derivative of the squared norm, not
% Gaussian through a one-dimensional marginal. Clause (1) is not proved; the node stays notready.
% CHANGED Q-0281: typed (the statement skeleton). Clause (2), sufficiency, is split to
% prop:anisotropy-budget-converse, placed after this node (ADR-0025; records/formalization/PRICING.md:
% it spends A3(iii) through prop:stable-exponent-cone, where clause (1) spends A5), and \uses gains
% it; the text is unchanged and the \lean tag names the child's declaration too.
% CHANGED Q-0281 wave 2, group G: clause (1) proved (AffineHemigroup.AnisotropyBudget; A5 and
% A3(i)-(ii), the smoothing and the budget Lean core in budget_of_contracting; guarded). Its proof is
% rewritten to the Lean route: the symmetric form of the exponent in place of the compensator
% change, the smoothing by s in (1/2, 3/2) with ds, the density identified by the radial image
% (no Haar measure on O(d)), and the budget read on M(r) = K(r)/r^{d-1}. Statement unchanged.
% CHANGED round-2 review (Gemini), item 15 (records/affine/reviews/VERIFY-round2-notation-paper.md):
% the frequency on R^d is omega, as the paper's rule asks from its section 3 on (was xi in the proofs
% of clause (2) and of the radial smoothing). The blind referee of the response
% (records/affine/REFEREE-round2-response-paper.md, item 8): the symmetric form of the transform in
% the proof of (1) wrote it z, now omega as well. Notation only; no step changes.
\begin{proposition}[the anisotropy budget of a rotation-invariant law]
\label{prop:anisotropy-budget}
\uses{prop:stable-exponent-cone,lem:elliptical-stable,prop:anisotropy-budget-converse,thm:cited-A3,thm:cited-A5}
\lean{AffineHemigroup.anisotropy_budget, AffineHemigroup.anisotropy_budget_converse}\leanok
Let $d \ge 1$, let $B$ be a real $d \times d$ matrix, and write
$\lambda := \lambda_{\min}(\mathrm{sym}\,B)$ (\ref{prop:stable-exponent-cone}).
\begin{enumerate}
\item If $e^{-tB} \to 0$ as $t \to \infty$ and $\mu$ is a rotation-invariant $B$-self-decomposable
law on $\RR^d$ that is not Gaussian (a point mass counting as Gaussian), then $\lambda > 0$ and
\[
  \tr B < (d+2)\,\lambda .
\]
\item Conversely, if $\tr B < (d+2)\lambda$, then $\lambda > 0$, $e^{-tB} \to 0$, and $S_\alpha(I)$ is
rotation invariant, not Gaussian, and $B$-self-decomposable for every $\alpha \in (0,2)$ with
$\alpha \ge \tr B/\lambda - d$.
\end{enumerate}
\statusA\quad\ledger{A5}\quad\ledger{A3}\quad\emph{(\textbf{Assignment.} \ledger{A5},
\sscite{jurek1993operator} Thm.~3.3.5, p.~83, with the opening of its \S 3.4, pp.~89--90 (the
original is \sscite{urbanik1972levy} Prop.~5.2, p.~134), is spent once, in clause (1): a
$B$-self-decomposable law with $e^{-tB} \to 0$ and its residuals are infinitely divisible.
\ledger{A3} then gives the law a \Levy{} measure, by (i), and by uniqueness (ii) makes it rotation
invariant and gives the inequality $(e^{-tB})_*\nu \le \nu$; clause (2) spends \ledger{A3}(iii)
through \ref{prop:stable-exponent-cone}. The radial smoothing, the infinitesimal inequality and the
integrability argument are this module's. \textbf{What it is worth.} With rotation invariance and
one linear flow, a non-Gaussian kernel exists exactly when the ratio $\tr B/\lambda$, which is at
least $d$ and equals $d$ for $B$ a multiple of $I$ plus a skew part, stays below $d + 2$: an
anisotropy budget of $2$, which is the ceiling of the stable index, and only the stable laws with
$\alpha$ close to $2$ come near it. The bound is strict, unlike \ref{eq:stable-exponent-cone}, because
the \Levy{} measure must integrate $|x|^2$ near the origin. The skew part of $B$ is free. With
\ref{prop:fixed-frame-families} this separates two obstructions: frame families take the
Gaussian's anisotropy per step and cannot turn; rotation-invariant non-Gaussian laws turn freely and have a bounded anisotropy
budget; only the Gaussian has neither limit. \textbf{Not claimed}: the case where $e^{-tB}$ does not
tend to $0$ (for $B$ skew every rotation-invariant law is $B$-self-decomposable with
$\rho_t = \delta_0$, and in general \ledger{A5} is not available); laws that are not rotation
invariant; the exponent cone of a law other than $S_\alpha(I)$.)}
\end{proposition}

\begin{proof}
\leanok
\emph{Clause (2).} $\tr B = \tr\mathrm{sym}\,B \ge d\lambda$, so $\tr B < (d+2)\lambda$ forces
$\lambda > 0$ (if $\lambda \le 0$ then $(d+2)\lambda \le d\lambda \le \tr B$); then, with
$y := e^{-tB}x$, $\tfrac{d}{dt}|y|^2 = -2\langle y,(\mathrm{sym}\,B)y\rangle \le -2\lambda|y|^2$, so
$|e^{-tB}x| \le e^{-t\lambda}|x|$ and $e^{-tB} \to 0$. The interval of $\alpha$ is nonempty, as
$\tr B/\lambda - d < 2$, and for such $\alpha$, $(d+\alpha)\lambda \ge \tr B$, so
\ref{prop:stable-exponent-cone} applies. $S_\alpha(I)$ has the transform $e^{-|\omega|^\alpha}$
(\ref{lem:elliptical-stable}(2)), which depends on $|\omega|$ only, so it is rotation invariant. It is
not Gaussian: its marginal along $e_1$ has the transform $e^{-|t|^\alpha}$, while a Gaussian marginal
has a transform of modulus $e^{-vt^2/2}$ for some $v \ge 0$; equality at $t = 1$ and $t = 2$ would give
$v = 2$ and $2^\alpha = 4$, so $\alpha = 2$.

\emph{Clause (1): infinite divisibility, and the \Levy{} measure.} Let $L_t := e^{-tB}$ and
$\mu = (L_t)_*\mu * \rho_t$. By \ledger{A5}, $\mu$ and every $\rho_t$ are infinitely divisible.
Let $(A_\mu, \nu, \gamma)$ be the generating triplet of $\mu$ (\ledger{A3}(i)). As $-I \in \Orth(d)$,
$\mu$ is symmetric, and uniqueness (\ledger{A3}(ii)) applied to the reflected triplet
$(A_\mu, (-I)_*\nu, -\gamma)$ gives $\gamma = 0$ and $(-I)_*\nu = \nu$, so that
$\widehat\mu(\omega) = \exp[-\frac12\langle \omega,A_\mu \omega\rangle - \int(1 - \cos\langle \omega,x\rangle)\,\nu(dx)]$.
Then $\nu \ne 0$, since otherwise $\mu$ is the Gaussian law $N(0, A_\mu)$. For $L \in \GL(d)$ this form,
read at $L^{\tp}\omega$, shows that $L_*\mu$ has the triplet $(LA_\mu L^{\tp}, L_*\nu, 0)$. The exponents of
a convolution add, so uniqueness applied to $\mu = (L_t)_*\mu * \rho_t$ gives $\nu = (L_t)_*\nu + \nu_t$
with $\nu_t$ the \Levy{} measure of $\rho_t$: hence $(L_t)_*\nu \le \nu$ for every $t > 0$. Applied to
$Q_*\mu = \mu$, $Q \in \Orth(d)$, it gives $Q_*\nu = \nu$.

\emph{Radial smoothing.} Fix a smooth $\eta$ on $\RR$ with $0 \le \eta \le 1$, $\eta = 1$ on
$[3/4, 5/4]$ and $\eta = 0$ off $(1/2, 3/2)$, and put $\nu_\eta(U) := \int\eta(s)\,\nu(s^{-1}U)\,ds$ for
Borel $U$, the average of the dilates $(sI)_*\nu$. It is invariant under $\Orth(d)$, null at the
origin and not $0$, and $(L_t)_*\nu_\eta \le \nu_\eta$, because $L_t$ commutes with $sI$. Where
$\eta(s) \ne 0$, $s < 3/2$ and $s^2|x|^2 \wedge 1 \le \frac94(|x|^2 \wedge 1)$, so
$\int(|y|^2 \wedge 1)\,\nu_\eta(dy) \le \frac94\int\eta\,\int(|x|^2 \wedge 1)\,\nu(dx) < \infty$, and
$\nu_\eta$ is finite off every ball about the origin. Substituting $r = s|x|$ and exchanging the
integrals, its image under $y \mapsto |y|$ is $K(r)\,dr$ on $(0,\infty)$, with
\[
  K(r) := \int\eta(r/|x|)\,|x|^{-1}\,\nu(dx) .
\]
For $r$ in a compact subset of $(0,\infty)$ only $x$ with $|x|$ in a compact subset of $(0,\infty)$
contribute, a set of finite $\nu$-measure, so $K$ is finite there and may be differentiated under the
integral. By polar coordinates the measure with density $h(y) := K(|y|)/(|S^{d-1}|\,|y|^{d-1})$ has the
same radial image $K(r)\,dr$, and it is invariant under $\Orth(d)$. A measure invariant under
$\Orth(d)$, null at the origin and finite off every ball about it is determined by its radial image:
on each shell $\{|x| > \varepsilon\}$ its restriction $m$ is finite and invariant, and averaging
$\widehat m(|\omega|\theta)$ over unit vectors $\theta$ against $\sigma_{d-1}$ and exchanging the integrals
gives $\widehat m(\omega) = |S^{d-1}|^{-1}\int\Phi(|\omega|,r)\,(|\cdot|_*m)(dr)$ with
$\Phi(s,r) := \int_{S^{d-1}}e^{isr\langle e_1,\theta\rangle}\,\sigma_{d-1}(d\theta)$. So $\nu_\eta$ has the
density $h$; no Haar measure on $\Orth(d)$ is needed.

\emph{The infinitesimal inequality.} Write $M(r) := K(r)/r^{d-1}$, so that $h(y) = M(|y|)/|S^{d-1}|$.
The linear change of variables $z = L_ty$, with $|\det L_t|^{-1} = e^{t\tr B}$, turns
$(L_t)_*\nu_\eta \le \nu_\eta$ into $e^{t\tr B}M(|e^{tB}z|) \le M(|z|)$ for almost every $z$, hence for
every $z \ne 0$: the set where it fails is open, both sides being continuous off the origin, and
null. Take $z = rv$ with $v$ a unit eigenvector of $\mathrm{sym}\,B$ for $\lambda$. Then
$\tfrac{d}{dt}|e^{tB}rv|$ at $t = 0$ is $r\langle v,Bv\rangle = r\lambda$, and a function of $t \ge 0$
that is largest at $t = 0$ has a nonpositive right derivative there:
\[
  \tr B\;M(r) + \lambda\,r\,M'(r) \le 0 \qquad (r > 0).
\]
Since $\det e^{-tB} = e^{-t\tr B} \to 0$, $\tr B > 0$. As $\nu_\eta \ne 0$ there is $r_0$ with
$M(r_0) > 0$.

\emph{The budget.} If $\lambda = 0$ the inequality gives $M(r_0) \le 0$, which is excluded. Otherwise
put $\kappa := \tr B/\lambda$; then $(r^\kappa M(r))' = r^{\kappa-1}(\tr B\,M(r) + \lambda rM'(r))/\lambda$,
which is $\le 0$ if $\lambda > 0$ and $\ge 0$ if $\lambda < 0$. If $\lambda < 0$, then $\kappa < 0$ and
$M(r) \ge M(r_0)r_0^\kappa r^{-\kappa} \ge M(r_0)r_0^\kappa$ for $r \ge \max(r_0, 1)$, so that
$K(r) = r^{d-1}M(r)$ is bounded below there by a positive constant and
$\int_0^\infty(r^2 \wedge 1)K(r)\,dr = \int(|y|^2 \wedge 1)\,\nu_\eta(dy)$ is infinite, which is excluded.
So $\lambda > 0$, and $M(r) \ge M(r_0)r_0^\kappa r^{-\kappa}$ for $0 < r \le r_0$, so that with
$r_1 := \min(r_0, 1)$
\[
  \int_0^{r_1}r^2K(r)\,dr \ge M(r_0)\,r_0^\kappa\int_0^{r_1}r^{d+1-\kappa}\,dr ,
\]
which is infinite if $\kappa \ge d+2$. Hence $\kappa < d + 2$, that is $\tr B < (d+2)\lambda$.
\end{proof}

% CHANGED Q-0281: split from prop:anisotropy-budget (ADR-0025; records/formalization/PRICING.md):
% clause (2), sufficiency, copied verbatim with the parent's clause number. It spends A3 through
% prop:stable-exponent-cone, where the parent's clause (1) spends A5, and it is a few hundred lines of
% Lean against the necessity's thousand; the parent \uses this node, and its text and numbering are
% unchanged.
% CHANGED Q-0281 wave 2, group E: proved (AffineHemigroup.AnisotropyConverse; prints A3(iii) through
% stable_exponent_cone, guarded). The parent proof clause (2) is rewritten to the Lean route: the
% contraction by the derivative of the squared norm, and not Gaussian by the transform of a one-
% dimensional marginal at t = 1 and t = 2, in place of the uniqueness A3(ii) of the Levy measure,
% which the assignment no longer names. Statement unchanged.
\begin{proposition}[the anisotropy budget is attained]
\label{prop:anisotropy-budget-converse}
\uses{prop:stable-exponent-cone,lem:elliptical-stable,thm:cited-A3}
\lean{AffineHemigroup.anisotropy_budget_converse}\leanok
Let $d \ge 1$, let $B$ be a real $d \times d$ matrix, and write
$\lambda := \lambda_{\min}(\mathrm{sym}\,B)$ (\ref{prop:stable-exponent-cone}).
\begin{enumerate}
\item[(2)] If $\tr B < (d+2)\lambda$, then $\lambda > 0$, $e^{-tB} \to 0$, and $S_\alpha(I)$ is
rotation invariant, not Gaussian, and $B$-self-decomposable for every $\alpha \in (0,2)$ with
$\alpha \ge \tr B/\lambda - d$.
\end{enumerate}
\statusA\quad\ledger{A3}\quad\emph{(\textbf{Assignment.} \ledger{A3}(iii) through
\ref{prop:stable-exponent-cone}, for the self-decomposability; ``not Gaussian'' reads the transform of
$S_\alpha(I)$ alone; \ledger{A5} is not spent. Clause (2) of \ref{prop:anisotropy-budget},
split off on 2026-10-04 (Q-0281, ADR-0025); the numbering stays the parent's.)}
\end{proposition}

\begin{proof}
\leanok
This is the part of the proof of \ref{prop:anisotropy-budget} headed clause (2).
\end{proof}

% not in the draft: the author's question after Q-0206, the row "one linear flow" (hub RESEARCH.md,
% Paper VII, the weakenings item), Q-0246 item 1. The converse in the manner of prop:ray-families,
% with lambda I replaced by e^{sB}; the class is Urbanik's (operator-self-decomposable laws), named
% in the annotation, not cited by the proof.
% CHANGED Q-0281: typed (the statement skeleton). The second sentence of clause (5), the stable
% member, is split to prop:one-parameter-stable-member, placed after this node (ADR-0025;
% records/formalization/PRICING.md: the one sentence of the node that spends A3), and \uses gains it;
% the text is unchanged and the \lean tag names the child's declaration too.
% CHANGED Q-0281: blind review R4-2: the tag names isExpContracting_iff, the equivalence clause (2)
% asserts with "that is" (every eigenvalue of B with positive real part iff e^{tB} -> 0 as t -> -oo),
% which its proof proves.
% CHANGED Q-0281 wave 2, group E: proved (AffineHemigroup.OneParameter, with isExpContracting_iff in
% AffineHemigroup.ExpContracting; Lean core, guarded; the child prints A3(iii)). The proof of clause (2)'s
% equivalence is written out on the Lean route (the eigenvector identity, an eigenvector of B inside an
% eigenspace of e^{-B}, Gelfand's formula) in place of the Jordan form. The proof
% of clause (3) is rewritten to the Lean route (the derivative of the quadratic form along the curve,
% in place of the matrix integral and the closed cone; the origin through the entries of
% e^{tB}Sigma_0^{1/2}), and the second half of clause (5) too (the transform at e_1 + e_2, in place of
% Hamel's equation), so the uses list drops prop:frame-families. Statement unchanged.
% CHANGED round-3 review (report E, required change 9; VERIFY-round3-smaller-groups-paper.md section 9):
% the annotation gives the whole range 1/(1+alpha) <= b <= 1 + alpha of the planar curves that S_alpha(I)
% rides, by prop:stable-exponent-cone, where it gave the half 1 <= b <= 1 + alpha. Annotation only.
\begin{proposition}[the families of one linear flow]
\label{prop:one-parameter-families}
\uses{prop:stable-exponent-cone,lem:elliptical-stable,prop:one-parameter-stable-member,thm:cited-A3}
\lean{AffineHemigroup.one_parameter_families, AffineHemigroup.one_parameter_families_converse,
AffineHemigroup.isExpContracting_iff, AffineHemigroup.one_parameter_origin, AffineHemigroup.one_parameter_curve,
AffineHemigroup.one_parameter_gaussian, AffineHemigroup.one_parameter_products,
AffineHemigroup.one_parameter_stable_member}\leanok
Let $d \ge 1$ and let $B$ be a real $d \times d$ matrix. Let $(\mu_t)_{t\in\RR}$ be probability
measures on $\RR^d$ with
\begin{enumerate}
\item[(F1)] \emph{(covariance under the flow)} $(e^{sB})_*\mu_t = \mu_{t+s}$ for all $s, t \in \RR$;
\item[(F2)] \emph{(a cascade along the flow)} for every $u > 0$ there is a probability measure
$\rho_u$ on $\RR^d$ with $\mu_u = \mu_0 * \rho_u$.
\end{enumerate}
\begin{enumerate}
\item \emph{(The characterization.)} Then $\mu_0$ is $B$-self-decomposable
(\ref{prop:stable-exponent-cone}), $\mu_t = (e^{tB})_*\mu_0$ for every $t$, and the family cascades
between any two of its members: $\mu_t = \mu_s * (e^{sB})_*\rho_{t-s}$ for $s < t$. Conversely, for
every $B$-self-decomposable law $\mu_0$ the family $\mu_t := (e^{tB})_*\mu_0$ satisfies (F1) and
(F2).
\item \emph{(The origin.)} If every eigenvalue of $B$ has positive real part, that is if
$e^{tB} \to 0$ as $t \to -\infty$, then $\mu_t \to \delta_0$ weakly as $t \to -\infty$.
\item \emph{(The curve of apertures.)} For $\Sigma_0 \in \PD(d)$ put
$\Sigma(t) := e^{tB}\Sigma_0e^{tB^{\tp}}$. Then $\Sigma(s) \le \Sigma(t)$ for all $s \le t$ in the
Loewner order if and only if $B\Sigma_0 + \Sigma_0B^{\tp} \in \PSD(d)$; and $\Sigma(t) \to 0$ as
$t \to -\infty$ if and only if $e^{tB} \to 0$.
\item \emph{(The Gaussian rides exactly the monotone curves.)} For $b > 0$ the family
$\mu_t := N(0, 2b\,\Sigma(t))$ satisfies (F1), and it satisfies (F2) if and only if
$B\Sigma_0 + \Sigma_0B^{\tp} \in \PSD(d)$.
\item \emph{(Two kinds of member.)} If $B = \diag(b_1,\dots,b_d)$ with every $b_i > 0$, then every
product $\nu_1 \otimes \dots \otimes \nu_d$ of self-decomposable laws on $\RR$ is
$B$-self-decomposable. If $d \ge 2$, $0 < \alpha < 2$ and
$(d+\alpha)\lambda_{\min}(\mathrm{sym}\,B) \ge \tr B$, then $S_\alpha(I)$ is $B$-self-decomposable
and is not a product law in any orthonormal frame.
\end{enumerate}
\statusA\quad\ledger{A3}\quad\emph{(\textbf{Assignment.} Clauses (1)--(4) cite nothing beyond the
uniqueness of the Fourier transform: pushforward and convolution, as in \ref{prop:ray-families}, and
for (4) the transform of a centred Gaussian law. Clause (5) spends \ledger{A3}(iii) through
\ref{prop:stable-exponent-cone}, for its second sentence only. \textbf{The class.} For $B$ with
every eigenvalue of positive real part these are Urbanik's operator-self-decomposable laws
(\sscite{urbanik1972levy} Thm.~5.1, p.~136, with his $Q = -B$; \sscite{sato1984operator} (2.1),
p.~76, ``$Q$-selfdecomposable'' with $Q = B$), named, not cited: no proof here reads them.
\textbf{What it is worth.} It is the row of the weakened axioms between the rays and a fixed frame.
With $B = I$ it is the ray cascade of \ref{prop:ray-families} without the rest of the affine group,
and so without rotation invariance; with $B$ diagonal it contains the products of
\ref{prop:fixed-frame-families} along every curve of the frame at once; and clause (5) shows it is
larger than both, since $S_\alpha(I)$ rides the curves of $B = \diag(1, b)$ in the plane exactly for
$1/(1+\alpha) \le b \le 1 + \alpha$ (\ref{prop:stable-exponent-cone}: $(2+\alpha)\min(1, b) \ge 1 + b$). For vision it is shape drift locked to scale: the aperture's shape is a fixed function
of its scale, set by $B$, and does not adapt to the image. Clause (4) is why the Loewner order
matters on the curve at all: only for the Gaussian is the label $\Sigma(t)$ the kernel's covariance,
and then a cascade along the curve is exactly a Loewner-monotone curve (the Gaussian's exponent cone
$\{B\Sigma_0 + \Sigma_0B^{\tp} \in \PSD(d)\}$, cf. \sscite{jurek1993operator} Prop.~3.3.10, p.~88,
named, not cited). \textbf{Not claimed}: anything for $B$ whose flow does not contract beyond
clause (1), which holds for every $B$ (a skew $B$ gives rotations only, under which a
rotation-invariant law is $B$-self-decomposable with $\rho_u = \delta_0$, and a $B$ with an
eigenvalue of negative real part never reaches the origin); the exponent cone of a law other than
the Gaussian and $S_\alpha(I)$; a representation of the class, which is Sato--Yamazato's (2.6),
p.~77.)}
\end{proposition}

\begin{proof}
\leanok
\emph{Clause (1).} (F1) with $t = 0$ gives $\mu_s = (e^{sB})_*\mu_0$. With $u > 0$, (F2) reads
$(e^{uB})_*\mu_0 = \mu_0 * \rho_u$; the image of a convolution under a linear map is the
convolution of the images, so applying $(e^{-uB})_*$ gives
$\mu_0 = (e^{-uB})_*\mu_0 * (e^{-uB})_*\rho_u$, which is $B$-self-decomposability with the residual
$(e^{-uB})_*\rho_u$. For $s < t$, $\mu_t = (e^{sB})_*\mu_{t-s} = (e^{sB})_*(\mu_0 * \rho_{t-s})
= \mu_s * (e^{sB})_*\rho_{t-s}$. Conversely, $(e^{sB})_*(e^{tB})_*\mu_0 = (e^{(s+t)B})_*\mu_0$ gives
(F1), and applying $(e^{uB})_*$ to $\mu_0 = (e^{-uB})_*\mu_0 * \rho'_u$ gives
$\mu_u = \mu_0 * (e^{uB})_*\rho'_u$, which is (F2).

\emph{Clause (2).} For $f$ bounded and continuous, $\int f\,d\mu_t = \int f(e^{tB}x)\,\mu_0(dx)$,
and $f(e^{tB}x) \to f(0)$ for every $x$ as $t \to -\infty$; dominated convergence gives
$\int f\,d\mu_t \to f(0)$. That $e^{tB} \to 0$ as $t \to -\infty$ exactly when every eigenvalue of
$B$ has positive real part rests on one identity: if $Au = zu$ for a complex matrix $A$ and
$u \ne 0$, then $e^Au = e^zu$, by the exponential series. If $Bu = zu$ and $e^{-tB} \to 0$ as
$t \to \infty$, then $e^{-tB}u = e^{-tz}u \to 0$, so $|e^{-tz}| = e^{-t\,\mathrm{Re}\,z} \to 0$ and
$\mathrm{Re}\,z > 0$. Conversely, let every eigenvalue of $B$ have positive real part and put
$E := e^{-B}$. Every eigenvalue $\mu$ of $E$ is $e^{-z}$ for an eigenvalue $z$ of $B$: $B$ commutes
with $E$, so it maps the eigenspace $\ker(E - \mu)$ into itself and has an eigenvector $u$ there,
and then $\mu u = Eu = e^{-z}u$. So $|\mu| = e^{-\mathrm{Re}\,z} < 1$, the spectral radius of $E$ is
below $1$, and by Gelfand's formula, $\|E^n\|^{1/n}$ tending to the spectral radius, $E^n \to 0$.
Finally $e^{-tB} = E^ne^{-(t - n)B}$ for $t \ge 0$ with $n$ the integer part of $t$, and the second
factor stays bounded, so $e^{-tB} \to 0$.

\emph{Clause (3).} $C := B\Sigma_0 + \Sigma_0B^{\tp}$ is symmetric. For $\omega \in \RR^d$ put
$v(t) := e^{tB^{\tp}}\omega$, so that $\langle\omega,\Sigma(t)\omega\rangle = \langle v(t),\Sigma_0v(t)\rangle$
and $v'(t) = B^{\tp}v(t)$; then
$\tfrac{d}{dt}\langle\omega,\Sigma(t)\omega\rangle = \langle v(t), Cv(t)\rangle$. If $C \in \PSD(d)$, this
derivative is nonnegative, so $t \mapsto \langle\omega,\Sigma(t)\omega\rangle$ is nondecreasing for every
$\omega$, which is $\Sigma(s) \le \Sigma(t)$ for $s \le t$. Conversely, if the curve is Loewner-monotone,
each of these functions is nondecreasing, so its derivative at $0$, $\langle\omega,C\omega\rangle$, is
nonnegative, and $C \in \PSD(d)$. For the origin, put $M(t) := e^{tB}\Sigma_0^{1/2}$, so that
$\Sigma(t) = M(t)M(t)^{\tp}$ and $\Sigma(t)_{ii} = \sum_kM(t)_{ik}^2 \ge M(t)_{ij}^2$. If $\Sigma(t) \to 0$,
then $M(t) \to 0$ entrywise and $e^{tB} = M(t)\Sigma_0^{-1/2} \to 0$; conversely, $e^{tB} \to 0$ gives
$e^{tB^{\tp}} = (e^{tB})^{\tp} \to 0$, and $\Sigma(t) \to 0$.

\emph{Clause (4).} The image of $N(0, C')$ under a linear map $L$ is $N(0, LC'L^{\tp})$, so
$(e^{sB})_*N(0, 2b\Sigma(t)) = N(0, 2b\Sigma(t+s))$, which is (F1). For (F2) at $u > 0$, a
probability measure $\rho_u$ with $N(0,2b\Sigma(u)) = N(0,2b\Sigma_0) * \rho_u$ must have transform
$\widehat\rho_u(\omega) = e^{-b\langle\omega,(\Sigma(u) - \Sigma_0)\omega\rangle}$, the Gaussian
transforms having no zeros. If $\Sigma(u) - \Sigma_0$ has a negative eigenvalue with eigenvector $v$,
then $|\widehat\rho_u(v)| > 1$, which no characteristic function allows; if it is positive
semidefinite, $N(0, 2b(\Sigma(u) - \Sigma_0))$ is such a $\rho_u$. So (F2) holds exactly when
$\Sigma_0 \le \Sigma(u)$ for every $u > 0$, which by (F1)'s transport $\Sigma(s+u) =
e^{sB}\Sigma(u)e^{sB^{\tp}}$ is the monotonicity of the whole curve, and clause (3) ends the proof.

\emph{Clause (5).} For $B = \diag(b_i)$ and $u > 0$ put $c_i := e^{-ub_i} \in (0,1)$; by
self-decomposability on $\RR$, $\nu_i = (c_i)_*\nu_i * \rho^{(i)}$ for probability measures
$\rho^{(i)}$, and since $e^{-uB}$ acts coordinatewise and a convolution of product measures is the
product of the convolutions of the factors,
$\nu_1 \otimes \dots \otimes \nu_d = (e^{-uB})_*(\nu_1 \otimes \dots \otimes \nu_d) * (\rho^{(1)} \otimes \dots \otimes \rho^{(d)})$. The second sentence's
first half is \ref{prop:stable-exponent-cone}. For the second half: $S_\alpha(I)$ has the transform
$e^{-|\omega|^\alpha}$ (\ref{lem:elliptical-stable}(2)), which depends on $|\omega|$ only, so it is
rotation invariant, and if it were $Q_*(\nu_1 \otimes \dots \otimes \nu_d)$ with $Q \in \Orth(d)$, applying
$(Q^{\tp})_*$ would make it the product $\nu_1 \otimes \dots \otimes \nu_d$ itself. Reading the transform at
$te_i$ gives $\widehat\nu_i(t) = e^{-|t|^\alpha}$, and at $\omega = e_1 + e_2$ the product gives
$e^{-2} = e^{-|e_1 + e_2|^\alpha} = e^{-2^{\alpha/2}}$, so $\alpha = 2$, which is excluded.
\end{proof}

% CHANGED Q-0281: split from prop:one-parameter-families (ADR-0025; records/formalization/PRICING.md):
% the second sentence of clause (5), copied verbatim with the parent's clause number. It is the one
% sentence of the node that spends A3 (through prop:stable-exponent-cone); the parent \uses this node,
% and its text and numbering are unchanged.
% CHANGED Q-0281 wave 2, group E: proved (AffineHemigroup.OneParameter; prints A3(iii) through
% stable_exponent_cone, guarded). The second half's proof (in the parent's clause (5)) is rewritten
% to the Lean route: the transform read at e_1 + e_2 forces alpha = 2, in place of Hamel's equation of
% prop:frame-families(3), so the uses list drops prop:frame-families. Statement unchanged.
\begin{proposition}[the stable member of one linear flow]
\label{prop:one-parameter-stable-member}
\uses{prop:stable-exponent-cone,lem:elliptical-stable,thm:cited-A3}
\lean{AffineHemigroup.one_parameter_stable_member}\leanok
Let $B$ be a real $d \times d$ matrix, and use $B$-self-decomposability as in
\ref{prop:stable-exponent-cone}.
\begin{enumerate}
\item[(5)] If $d \ge 2$, $0 < \alpha < 2$ and $(d+\alpha)\lambda_{\min}(\mathrm{sym}\,B) \ge \tr B$,
then $S_\alpha(I)$ is $B$-self-decomposable and is not a product law in any orthonormal frame.
\end{enumerate}
\statusA\quad\ledger{A3}\quad\emph{(\textbf{Assignment.} \ledger{A3}(iii) through
\ref{prop:stable-exponent-cone}, for the first half; the second half reads the transform of
$S_\alpha(I)$ (\ref{lem:elliptical-stable}(2)) at $e_i$ and at $e_1 + e_2$.
The second sentence of \ref{prop:one-parameter-families}(5), split off on 2026-10-04 (Q-0281,
ADR-0025); the numbering stays the parent's.)}
\end{proposition}

\begin{proof}
\leanok
This is the part of the proof of \ref{prop:one-parameter-families} headed clause (5) that treats
the clause's second sentence (its first half and its second half).
\end{proof}

% not in the draft: P-0019 (i) and the elliptical-stable case of conjecture A' (notes/HANDOFF.md),
% Q-0241. Conjecture A (hub RESEARCH.md, 2026-09-30): under similarity covariance with the whole
% Loewner order, every family is an anisotropic Gaussian convolved with an isotropic, shape-blind
% part. Clause (3) refutes it; clause (2) is A' for elliptical stable parts, by the proposal's
% sandwich. The sufficient bound on kappa is the proposal's crude one; the sampled threshold
% (about 1 + sqrt 2 at alpha = 1, d = 2) is not claimed.
% CHANGED Q-0281: \uses re-pointed to the split child it reads ((1)(b)-(c), (2) and (c) of prop:stable-step-cone beside its (1)(a)).
% CHANGED Q-0281 wave 2, group G: proved (AffineHemigroup.SimilarityStable; clauses (1)-(2) Lean
% core, clause (3) A3(iii); guarded). The proof of clause (3)'s stable increments is rewritten to the
% Lean route: condition (c) for the whole step at once, by the arithmetic-geometric mean inequality
% and Bernoulli's inequality, in place of the derivative of log n_{H_s} along the segment.
% Statement unchanged.
% CHANGED Q-0469 (wording only; records/affine/GATE-affine.md, Q-0317 row "App. H 11"): clause (3)'s
% heading "(A family that is not shape-blind.)" -> "(A stable part that is not rotation invariant.)",
% the term "shape-blind" being undefined in the node; the heading now says what the clause's last
% sentence states. Label, clause text and \lean unchanged.
\begin{proposition}[similarity-covariant families with a stable part]
\label{prop:similarity-stable-family}
\uses{lem:elliptical-stable,lem:stable-tail,prop:stable-step-cone,lem:affine-gauge,lem:stable-step-order,thm:cited-A3}
\lean{AffineHemigroup.sim_stable_parts_ordered, AffineHemigroup.sim_stable_bounded_shape,
AffineHemigroup.sim_stable_not_shape_blind}\leanok
Let $d \ge 2$ and $0 < \alpha < 2$. For maps $h : \PD(d) \to \PD(d)$ and $C : \PD(d) \to \PSD(d)$
put $\mu_\Sigma := N(0,C(\Sigma)) * S_\alpha(h(\Sigma))$, with $N(0,C)$ the centred Gaussian law of
covariance $C$.
\begin{enumerate}
\item \emph{(The stable parts are ordered.)} If $\Sigma, \Sigma' \in \PD(d)$ and
$\mu_{\Sigma'} = \mu_\Sigma * \rho$ for a positive finite measure $\rho$, then
$h(\Sigma) \preceq_\alpha h(\Sigma')$.
\item \emph{(Bounded shape.)} Suppose that $h(cQ\Sigma Q^{\tp}) = c\,Q\,h(\Sigma)\,Q^{\tp}$ for $c > 0$
and $Q \in \Orth(d)$, that $h$ extends continuously to $\PSD(d) \setminus \{0\}$ with values in
$\PD(d)$, and that every Loewner step $\Sigma \le \Sigma'$ in $\PD(d)$ has an increment as in (1).
Then $h(I) = h_0I$ with $h_0 > 0$, and with $c_1 := \lambda_{\max}(h(e_1e_1^{\tp})) > 0$,
\[
  \frac{\lambda_{\max}(h(\Sigma))}{\lambda_{\min}(h(\Sigma))} \le \Bigl(\frac{h_0}{c_1}\Bigr)^{\alpha}
  \qquad \text{for every } \Sigma \in \PD(d) .
\]
\item \emph{(A stable part that is not rotation invariant.)} Let $b \ge 0$ and $\kappa > 0$ with
$\alpha\kappa^2 - (d+1)\kappa - 1 \ge 0$, and put $h(\Sigma) := \Sigma + \kappa\tr(\Sigma)I$ and
$C(\Sigma) := 2b\Sigma$ for $\Sigma \in \PSD(d)$, with $S_\alpha(0) := \delta_0$. For every
$\Sigma \le \Sigma'$ in $\PSD(d)$ there is a probability measure $\rho_{\Sigma,\Sigma'}$ with
$\mu_{\Sigma'} = \mu_\Sigma * \rho_{\Sigma,\Sigma'}$, and the family $(\rho_{\Sigma,\Sigma'})$ satisfies
(K1), (K2) and (K4) of \ref{def:affine-cascade-family}, and (K3) for every similarity $A = cQ$,
$c > 0$, $Q \in \Orth(d)$. Its stable part $S_\alpha(h(\Sigma))$ is rotation invariant only when
$\Sigma$ is a multiple of $I$.
\end{enumerate}
\statusA\quad\ledger{A3}\quad\emph{(\textbf{Assignment.} \ledger{A3}(iii), through
\ref{lem:elliptical-stable}(2) and \ref{prop:stable-step-cone}(1), for the laws and the stable
increments of clause (3); clauses (1) and (2) spend nothing beyond the existence of the laws.
\textbf{What it is worth.} Clause (3) refutes the hub's conjecture A (2026-09-30), that under
similarity covariance with the whole Loewner order every family is an anisotropic Gaussian convolved
with an isotropic, shape-blind part: here the non-Gaussian part carries the shape of $\Sigma$, and
$b = 0$ is allowed, so no Gaussian part is needed. Clause (2) is the elliptical-stable case of the
replacement, conjecture A$'$ (P-0019): the stable part's eccentricity is bounded,
by a bound that tends to $1$ as $\alpha \to 0$; for clause (3)'s family it is
$((1+d\kappa)/(1+\kappa))^\alpha$. \textbf{Not claimed here}: conjecture A$'$ for non-Gaussian parts
that are not elliptical stable, and its index-$0$ case, exact shape-blindness, which are
\ref{prop:similarity-orbit-family} for elliptical parts of every self-decomposable profile (under
infinitely divisible increments, or without a Gaussian part), and
\ref{prop:similarity-general-bounded-shape} and \ref{prop:similarity-general-index-zero} for a general
part; the least $\kappa$ for which clause (3) holds (the bound used is sufficient, not sharp);
clause (2) without the extension of $h$ to rank one with values in $\PD(d)$, which supplies the lower
end of the sandwich: without it the clause is false (\ref{prop:similarity-degenerate-family}, Q-0264),
and only the trade-off $\lambda_{\max}(h(\Sigma))/\lambda_{\min}(h(\Sigma)) \le (h_0/\lambda_{\max}(h(\Sigma)))^\alpha$
for $\lambda_{\max}(\Sigma) = 1$, which the proof reads off the step $\Sigma \le I$ alone, survives.)}
\end{proposition}

% CHANGED writing-rules pass (2026-10-14; the blind check of records/affine/PASS-writing-rules-affine.md):
% notation only in clause (1), "with G = (C(Sigma') - C(Sigma))/2" -> "with (C(Sigma') - C(Sigma))/2 for
% its matrix C", the matrix of the Gaussian factor of lem:stable-tail being C since Q-0327.
\begin{proof}
\leanok
\emph{Clause (1).} The transforms of $\mu_\Sigma$ have no zero, so
$\widehat\rho(\omega) = \exp[-\tfrac12\langle\omega,(C(\Sigma') - C(\Sigma))\omega\rangle
- \psi_{h(\Sigma')}(\omega) + \psi_{h(\Sigma)}(\omega)]$, which is \ref{eq:stable-tail-transform} with
$\tfrac12(C(\Sigma') - C(\Sigma))$ for its matrix $C$. By \ref{lem:stable-tail}, $n_{h(\Sigma)} \le n_{h(\Sigma')}$,
which is $h(\Sigma) \preceq_\alpha h(\Sigma')$ by \ref{prop:stable-step-cone}(1).

\emph{Clause (2).} With $c = 1$ the hypothesis is the equivariance of \ref{lem:affine-gauge}'s proof
under $\Orth(d)$, which gives $h(I) = h_0I$, $h_0 > 0$; by continuity the equivariance extends to
$\PSD(d) \setminus \{0\}$, and $c_1 > 0$ because $h(e_1e_1^{\tp}) \in \PD(d)$. Let $\Sigma \in \PD(d)$;
by the equivariance under $c > 0$, which multiplies $h(\Sigma)$ by $c$, we may take
$\lambda_{\max}(\Sigma) = 1$. Let $Q \in \Orth(d)$ carry $e_1$ to a unit eigenvector for that
eigenvalue, and for $0 < \varepsilon \le \lambda_{\min}(\Sigma)$ put
$\Sigma_\varepsilon := Q\diag(1,\varepsilon,\dots,\varepsilon)Q^{\tp}$, so that
$\Sigma_\varepsilon \le \Sigma \le I$. By clause (1), $h(\Sigma_\varepsilon) \preceq_\alpha h(\Sigma)
\preceq_\alpha h_0I$, and $\preceq_\alpha$ is contained in the Loewner order
(\ref{prop:stable-step-cone}(2)). Letting $\varepsilon \to 0$,
$Q\,h(e_1e_1^{\tp})\,Q^{\tp} \le h(\Sigma)$, so the largest eigenvalue $\eta_d$ of $h(\Sigma)$ is at
least $c_1$. Let $\eta_1 \le \dots \le \eta_d$ be the eigenvalues of $h(\Sigma)$. For
$h(\Sigma) \preceq_\alpha h_0I$ the matrix $M = h_0h(\Sigma)^{-1}$ has $\lambda_{\min}(M) = h_0/\eta_d$ and
$\det M = h_0^d/(\eta_1\cdots\eta_d)$, and condition (c) of \ref{prop:stable-step-cone} reads
$\prod_i(\eta_i/\eta_d) \ge (\eta_d/h_0)^\alpha$. The left-hand side is at most $\eta_1/\eta_d$, so
$\eta_d/\eta_1 \le (h_0/\eta_d)^\alpha \le (h_0/c_1)^\alpha$.

% CHANGED Q-0281 wave 2, group G: proof rewritten to the Lean route (one step, no derivative).
\emph{Clause (3), the stable increments.} Let $\Sigma \le \Sigma'$ in $\PSD(d)$ and $P := \Sigma' -
\Sigma$. If $\Sigma = 0$ take $\rho^{\mathrm s} := S_\alpha(h(\Sigma'))$. Otherwise $\tau := \tr\Sigma > 0$,
so $H := h(\Sigma)$ and $h(\Sigma') = H + V$, $V := P + \kappa\tr(P)I$, are in $\PD(d)$; we check
condition (c) of \ref{prop:stable-step-cone}(1) for the relative step
\[
  M := H^{-1/2}h(\Sigma')H^{-1/2} = I + H^{-1/2}PH^{-1/2} + \kappa\tr(P)H^{-1} .
\]
Since $0 \le \Sigma \le \tau I$, $\kappa\tau I \le H \le (1+\kappa)\tau I$. Put $u := \tr(P)/\tau$,
$A := (1+d\kappa)/\kappa$ and $B := \kappa/(1+\kappa)$. As $H^{-1/2}PH^{-1/2} \ge 0$ and
$H^{-1} \ge ((1+\kappa)\tau)^{-1}I$, $\lambda_{\min}(M) \ge 1 + Bu$. As
$\tr(H^{-1/2}PH^{-1/2}) = \tr(P^{1/2}H^{-1}P^{1/2}) \le \tr(P)/\lambda_{\min}(H)$ and
$\tr H^{-1} \le d/\lambda_{\min}(H)$, with $\lambda_{\min}(H) \ge \kappa\tau$, $\tr M \le d + Au$. The
hypothesis on $\kappa$ is $(d+\alpha)\kappa^2 \ge (1+d\kappa)(1+\kappa)$, that is $A \le (d+\alpha)B$.
By the inequality of arithmetic and geometric means for the eigenvalues of $M$, and Bernoulli's
inequality $1 + qx \le (1+x)^q$ for $q = (d+\alpha)/d \ge 1$ and $x \ge 0$,
\[
  \det M \le \Bigl(\frac{\tr M}{d}\Bigr)^d \le \Bigl(1 + \frac{Au}{d}\Bigr)^d
  \le \Bigl(1 + \frac{d+\alpha}{d}\,Bu\Bigr)^d \le (1 + Bu)^{d+\alpha} \le \lambda_{\min}(M)^{d+\alpha} .
\]
Hence $n_H \le n_{h(\Sigma')}$, and \ref{prop:stable-step-cone}(1) gives an infinitely divisible
probability measure $\rho^{\mathrm s}$ with $S_\alpha(h(\Sigma')) = S_\alpha(H) * \rho^{\mathrm s}$. Put
$\rho_{\Sigma,\Sigma'} := N(0,2bP) * \rho^{\mathrm s}$.

\emph{Clause (3), the axioms.} In every case $\widehat\rho_{\Sigma,\Sigma'}(\omega) =
\exp[-b\langle\omega,P\omega\rangle - \psi_{h(\Sigma')}(\omega) + \psi_{h(\Sigma)}(\omega)]$, with
$\psi_0 := 0$. These exponents telescope, which is (K1); they are continuous in $(\Sigma,\Sigma')$,
which is (K2); $\mu_I \ne \delta_0$, which is (K4). For $A = cQ$,
$\widehat\rho_{\Sigma,\Sigma'}(A^{\tp}\omega)$ is the same expression with $A\Sigma A^{\tp}$,
$A\Sigma'A^{\tp}$ and $h(A\Sigma A^{\tp}) = c^2Q\,h(\Sigma)\,Q^{\tp} = A\,h(\Sigma)A^{\tp}$, which is (K3)
for $A$; each time the transform determines the law. Finally, $S_\alpha(H)$ is rotation invariant
exactly when $\psi_H(Q^{\tp}\omega) = \psi_H(\omega)$ for all $Q$, that is when $H$ commutes with
$\Orth(d)$, which forces $H$ to be a multiple of $I$ as in the proof of \ref{lem:affine-gauge}; and
$h(\Sigma) = \Sigma + \kappa\tr(\Sigma)I$ is a multiple of $I$ only when $\Sigma$ is.
\end{proof}

% not in the draft: Q-0246 item 3, conjecture A' (notes/HANDOFF.md) for non-Gaussian parts that are
% elliptical orbits of any self-decomposable profile, under infinitely divisible increments; and its
% index-0 clause, exact shape-blindness, two ways. The record check-gaussian-matern-step.py shows the
% hypothesis cannot be dropped when a Gaussian part is present.
% Title until Q-0284: "conjecture A' for elliptical parts". The article drops the working-note names
% "conjecture A" and "conjecture A'" and states the results directly (the author, 2026-10-04).
% CHANGED Q-0281 wave 2, group G: proved (AffineHemigroup.SimilarityOrbit, on the triplet calculus of
% AffineHemigroup.TripletSum; clauses (1)-(3) print A3(i)-(iii), clause (4) Lean core; guarded). The
% last step of clause (1) is rewritten to the Lean route: D m(r) <= m(r/lambda) is the ordering at
% every point directly, since u(y) <= |y|/lambda. Statement unchanged.
% CHANGED round-3 review (report D, item 5): clause (1) and the proof of (2) say "the step set of
% prop:orbit-step-cone(2)" where they said "the cone", that set being a cone only for some profiles.
% Wording only; no claim, \uses or \lean changes. prop:wedge-axis-profiles reads clause (1) with h = id.
\begin{proposition}[similarity-covariant families with an elliptical part: bounded shape, and shape-blind at index $0$]
\label{prop:similarity-orbit-family}
\uses{prop:orbit-step-cone,prop:log-orbit-rays,prop:similarity-stable-family,lem:affine-gauge,thm:cited-A3}
\lean{AffineHemigroup.sim_orbit_id_increment, AffineHemigroup.sim_orbit_bounded_shape,
AffineHemigroup.sim_orbit_index_zero, AffineHemigroup.sim_orbit_log}\leanok
Let $d \ge 2$, let $k$ be as in \ref{prop:orbit-step-cone} and right-continuous, with dilation ratio
$R$, and write $\nu_H := (H^{1/2})_*\nu_I$, $H \in \PD(d)$, for the orbit of the law $\nu_I$ of
\ref{prop:orbit-step-cone}(1), with \Levy{} densities $n_H$. For maps $h : \PD(d) \to \PD(d)$ and
$C : \PD(d) \to \PSD(d)$ put $\mu_\Sigma := N(0,C(\Sigma)) * \nu_{h(\Sigma)}$.
\begin{enumerate}
\item \emph{(Infinitely divisible increments order the parts.)} If $\Sigma, \Sigma' \in \PD(d)$ and
$\mu_{\Sigma'} = \mu_\Sigma * \rho$ with $\rho$ infinitely divisible, then $C(\Sigma) \le C(\Sigma')$
and $n_{h(\Sigma)} \le n_{h(\Sigma')}$ everywhere: the step $h(\Sigma) \to h(\Sigma')$ is in the step
set of \ref{prop:orbit-step-cone}(2).
\item \emph{(Bounded shape.)} Suppose $h(cQ\Sigma Q^{\tp}) = c\,Q\,h(\Sigma)\,Q^{\tp}$ for $c > 0$ and
$Q \in \Orth(d)$, that $h$ extends continuously to $\PSD(d) \setminus \{0\}$ with values in $\PD(d)$,
and that every Loewner step $\Sigma \le \Sigma'$ in $\PD(d)$ has an infinitely divisible increment.
Then $h(I) = h_0I$ with $h_0 > 0$, and with $c_1 := \lambda_{\max}(h(e_1e_1^{\tp}))$,
\[
  \frac{\lambda_{\max}(h(\Sigma))}{\lambda_{\min}(h(\Sigma))} \le R\bigl((h_0/c_1)^{1/2}\bigr)^2
  \qquad \text{for every } \Sigma \in \PD(d) .
\]
\item \emph{(Index $0$: shape-blind.)} If $k(0{+})$ is finite, $h$ is equivariant as in (2) (no
extension needed), and every Loewner step has an infinitely divisible increment, then $h(\Sigma)$
is a multiple of $I$ for every $\Sigma$.
\item \emph{(Logarithmic exponent, no Gaussian part.)} Let instead the non-Gaussian parts be the
orbit $(\nu_H)$ of a law satisfying the hypothesis of \ref{prop:log-orbit-rays}, let $C \equiv 0$
and $h$ be equivariant as in (2), and suppose every Loewner step has a positive finite increment.
Then $h(\Sigma)$ is a multiple of $I$ for every $\Sigma$.
\end{enumerate}
\statusA\quad\ledger{A3}\quad\emph{(\textbf{Assignment.} \ledger{A3}(ii), the uniqueness of the
generating triplet, in clause (1); \ledger{A3}(iii) through \ref{prop:orbit-step-cone} for the
orbit itself. Clause (4) cites nothing beyond \ref{prop:log-orbit-rays}. \textbf{What it is worth.}
It is conjecture A$'$ for every non-Gaussian part that is an elliptical orbit of a self-decomposable
profile: under similarity covariance with infinitely divisible increments along the whole Loewner
order, the part's eccentricity is bounded by the profile's dilation ratio, $(h_0/c_1)^\alpha$ for the
stable profile (\ref{prop:similarity-stable-family}(2), there without the infinite divisibility, by
\ref{lem:stable-tail}), and the part is exactly shape-blind at index $0$, a profile bounded at the
origin. Without a Gaussian part, clause (4) gets shape-blindness from positive increments alone.
\textbf{Why the infinite divisibility.} With a Gaussian part a positive increment need not order the
parts. In the plane, with $C(\Sigma) = 2b\Sigma$, $b = 2$, the \Matern{} orbit of index $1$ and
$h(\Sigma) = \Sigma + \kappa\tr(\Sigma)I$, $\kappa = 2$, the step from $\diag(1, 0.3)$ by
$\diag(0.05, 0.55)$ has a positive increment (its density computed by FFT: minimum $-10^{-17}$ of its
maximum, as for the ray steps that calibrate it, at two resolutions), although $h(\Sigma')$ is not a
multiple of $h(\Sigma)$; for $b = 0.1$ the same step is negative with mass $10^{-2}$
(\texttt{records/affine/check-gaussian-matern-step.py}, output in \texttt{check-gaussian-matern-step-output.txt}). That increment is
positive and not infinitely divisible, since its \Levy{} measure would be the signed
$n_{h(\Sigma')} - n_{h(\Sigma)}$. The Gaussian part does real work: it smooths the obstruction of
\ref{prop:log-orbit-rays} away on full-rank steps. Infinitely divisible increments are what the
draft's Lemma~2.6 obtains for a continuous cascade by the null-array argument along a Loewner segment,
and the module now takes that route: \ref{thm:cascade-infinitely-divisible} proves it (ledger entry
A6), and \ref{prop:similarity-cascade-families}(3) discharges the hypothesis of clauses (2) and (3) for
every family of kernels from scale zero of a (similarities, Loewner)-cascade family. The step above is a
single step, and by that node its kernels are those of no such cascade.
\textbf{Not claimed here}: A$'$ for non-Gaussian parts that are not elliptical orbits of one profile,
which is \ref{prop:similarity-general-bounded-shape} under a rank-one limit comparable to the isotropic
member and \ref{prop:similarity-general-index-zero} at index $0$, where exact shape-blindness fails for a
general part (\ref{prop:similarity-line-family}(3)); clause (2) without the extension of $h$ to rank
one with values in $\PD(d)$, which is false (\ref{prop:similarity-degenerate-family}, Q-0264);
clause (1), and so A$'$ with a Gaussian part, under positive increments only, for kernel families
that are not cascades (open: the near-rank-one steps, where the Gaussian part is thin, were not
resolved numerically; for cascade families the question does not arise,
\ref{prop:similarity-cascade-families}).)}
\end{proposition}

\begin{proof}
\leanok
\emph{Clause (1).} $\mu_\Sigma$ is infinitely divisible with triplet
$(C(\Sigma), n_{h(\Sigma)}(x)\,dx, 0)$, a Gaussian law and the law of \ref{prop:orbit-step-cone}(1)
convolved, and so is $\mu_{\Sigma'}$. Let $(A_\rho, \nu_\rho, \gamma_\rho)$ be the triplet of $\rho$
(\ledger{A3}(i)). The exponents of a convolution add, so $\mu_\Sigma * \rho$ has the triplet
$(C(\Sigma) + A_\rho, n_{h(\Sigma)}(x)\,dx + \nu_\rho, \gamma_\rho)$, and by uniqueness
(\ledger{A3}(ii)) $C(\Sigma') = C(\Sigma) + A_\rho \ge C(\Sigma)$ and $n_{h(\Sigma')}(x)\,dx =
n_{h(\Sigma)}(x)\,dx + \nu_\rho$, so $n_{h(\Sigma)} \le n_{h(\Sigma')}$ almost everywhere. To make it
everywhere, substitute as in the proof of \ref{prop:orbit-step-cone}(2): with
$M := h(\Sigma)^{-1/2}h(\Sigma')h(\Sigma)^{-1/2}$, $D := (\det M)^{1/2}$,
$\lambda := \lambda_{\min}(M)^{1/2}$ and $u(y) := \langle y,M^{-1}y\rangle^{1/2} \le |y|/\lambda$, the ordering
says $D\,m(|y|) \le m(u(y))$ for almost every $y$. Fix $r > 0$ and $\delta > 0$. The open set
$U := \{y : r < |y| < r + \delta,\ u(y) > r/\lambda\}$ is not empty (it contains $sv$ for $v$ a unit
eigenvector for $\lambda_{\min}(M)$ and $s$ slightly above $r$), so it has positive measure, and at almost
every $y \in U$, $m$ being nonincreasing,
$D\,m(r + \delta) \le D\,m(|y|) \le m(u(y)) \le m(r/\lambda)$. As $\delta \downarrow 0$, the
right-continuity of $m$ gives $D\,m(r) \le m(r/\lambda)$ for every $r > 0$.
% CHANGED Q-0281 wave 2, group G: proof rewritten to the Lean route (the ordering everywhere read off directly).
That is the ordering everywhere: for $y \ne 0$, $0 < u(y) \le |y|/\lambda$, and $m$ is
nonincreasing, so $D\,m(|y|) \le m(|y|/\lambda) \le m(u(y))$. The cone's inequalities are then
\ref{prop:orbit-step-cone}(2).

\emph{Clause (2).} As in the proof of \ref{prop:similarity-stable-family}(2): $\Orth(d)$-equivariance
gives $h(I) = h_0I$; after scaling, $\lambda_{\max}(\Sigma) = 1$, and $\Sigma_\varepsilon \le \Sigma \le I$
with $\Sigma_\varepsilon := Q\diag(1,\varepsilon,\dots,\varepsilon)Q^{\tp}$. By clause (1) both steps of
$h$ are in the step set of \ref{prop:orbit-step-cone}(2), which is contained in the Loewner order; letting
$\varepsilon \to 0$, $\eta_d := \lambda_{\max}(h(\Sigma)) \ge c_1$. Let $\eta_1 \le \dots \le \eta_d$ be the
eigenvalues of $h(\Sigma)$. For the step $h(\Sigma) \to h_0I$, $M = h_0h(\Sigma)^{-1}$,
$\lambda_{\min}(M) = h_0/\eta_d$ and $\det M = h_0^d/(\eta_1\cdots\eta_d)$, so the inequality of that step set reads
$\prod_i(\eta_d/\eta_i) \le R((h_0/\eta_d)^{1/2})^2$. The left-hand side is at least $\eta_d/\eta_1$,
and $R$ is nondecreasing ($k(r/\lambda)$ is nondecreasing in $\lambda$) with $h_0/\eta_d \le h_0/c_1$.

\emph{Clause (3).} $R \equiv 1$ (\ref{prop:orbit-step-cone}(3)), so by clause (1) every Loewner step
moves $h$ along a ray. Given $\Sigma$, take $t \ge \lambda_{\max}(\Sigma)$, so $\Sigma \le tI$; then
$h(tI) = t\,h(I) = th_0I$ is a multiple of $h(\Sigma)$, and $h(\Sigma)$ is a multiple of $I$.

\emph{Clause (4).} With $C \equiv 0$, $\mu_\Sigma = \nu_{h(\Sigma)}$ and the increment of
$\Sigma \le tI$ is a positive increment of the orbit $(\nu_H)$ from $h(\Sigma)$ to $h(tI) = th_0I$;
\ref{prop:log-orbit-rays} makes $th_0I$ a multiple of $h(\Sigma)$.
\end{proof}

% not in the draft: Q-0264 item 1 (the author's request of 2026-10-03; conjecture A' is the
% postdoc's, hub offices/postdoc/notes/2026-09-30-affine-weakenings.md, P-0019). The sandwich of
% prop:similarity-orbit-family(2) for a non-Gaussian part with no orbit structure. What writing it
% down found: the rank-one limit needs no continuity hypothesis, monotonicity alone makes it a
% decreasing limit; the question is its size, which the next three nodes settle.
% CHANGED Q-0281 wave 2, group G: proved (AffineHemigroup.SimilarityGeneral; every clause prints
% A3(i)-(ii) only, clause (3) through clause (1)'s symmetric collapse; guarded). The proof of record
% is the Lean's route and is unchanged. Statement unchanged.
\begin{proposition}[the sandwich for a general similarity-covariant part]
\label{prop:similarity-general-sandwich}
\uses{def:affine-cascade-family,prop:similarity-orbit-family,lem:elliptical-stable,thm:cited-A3}
\lean{AffineHemigroup.sim_general_covariance, AffineHemigroup.sim_general_id_increment,
AffineHemigroup.sim_general_finite_increment, AffineHemigroup.sim_general_rank_one_limit,
AffineHemigroup.sim_general_rank_one_exp_limit, AffineHemigroup.sim_general_sandwich}\leanok
Let $d \ge 2$. For each $\Sigma \in \PD(d)$ let $\mu_\Sigma$ be a symmetric infinitely divisible
probability measure on $\RR^d$ with generating triplet $(C(\Sigma), L_\Sigma, 0)$, and put
$\psi_\Sigma(\omega) := \int(1 - \cos\langle\omega,x\rangle)\,L_\Sigma(dx)$; so
$\mu_\Sigma = N(0,C(\Sigma)) * \nu_\Sigma$, where the \emph{non-Gaussian part} $\nu_\Sigma$ has
$\widehat\nu_\Sigma = e^{-\psi_\Sigma}$. Suppose the family is covariant under similarities,
$(\sqrt c\,Q)_*\mu_\Sigma = \mu_{cQ\Sigma Q^{\tp}}$ for $c > 0$, $Q \in \Orth(d)$ and
$\Sigma \in \PD(d)$, which is (K3) of \ref{def:affine-cascade-family} for $A = \sqrt c\,Q$. Put
$\Sigma_\varepsilon := \diag(1,\varepsilon,\dots,\varepsilon)$ for $0 < \varepsilon \le 1$.
\begin{enumerate}
\item \emph{(Covariance.)} $C(cQ\Sigma Q^{\tp}) = c\,Q\,C(\Sigma)\,Q^{\tp}$,
$L_{cQ\Sigma Q^{\tp}} = (\sqrt c\,Q)_*L_\Sigma$ and
$\psi_{cQ\Sigma Q^{\tp}}(\omega) = \psi_\Sigma(\sqrt c\,Q^{\tp}\omega)$. In particular $L_I$ and
$\psi_I$ are rotation invariant.
\item \emph{(Infinitely divisible increments.)} If $\mu_{\Sigma'} = \mu_\Sigma * \rho$ with $\rho$
infinitely divisible, then $C(\Sigma) \le C(\Sigma')$, $L_\Sigma \le L_{\Sigma'}$ and
$\psi_\Sigma \le \psi_{\Sigma'}$.
\item \emph{(Positive increments without a Gaussian part.)} If $C \equiv 0$ and
$\mu_{\Sigma'} = \mu_\Sigma * \rho$ with $\rho$ a positive finite measure, then $\rho$ is a
probability measure, $\widehat\rho = e^{-(\psi_{\Sigma'} - \psi_\Sigma)}$, and
$\psi_\Sigma \le \psi_{\Sigma'}$.
\item \emph{(The rank-one limit.)} Suppose every Loewner step $\Sigma \le \Sigma'$ in $\PD(d)$ has an
infinitely divisible increment. Then $L_*(B) := \lim_{\varepsilon\downarrow0}L_{\Sigma_\varepsilon}(B)$
exists for every Borel set $B$ bounded away from the origin, and $L_*$ extends to a \Levy{} measure
with $L_* \le L_{\Sigma_\varepsilon}$ for every $\varepsilon$; and
$\psi_*(\omega) := \lim_{\varepsilon\downarrow0}\psi_{\Sigma_\varepsilon}(\omega)$ exists for every
$\omega$. Both are invariant under every $R \in \Orth(d)$ with $Re_1 = \pm e_1$. If instead
$C \equiv 0$ and every Loewner step has a positive finite increment, $\psi_*$ exists and has the same
invariance.
\item \emph{(The sandwich.)} Let $\Sigma \in \PD(d)$ with $\lambda_{\max}(\Sigma) = 1$, and let
$Q \in \Orth(d)$ have as its first column a unit eigenvector of $\Sigma$ for the eigenvalue $1$. Under
the first hypothesis of (4),
\[
  Q_*L_* \le L_\Sigma \le L_I
  \qquad\text{and}\qquad
  \psi_*(Q^{\tp}\omega) \le \psi_\Sigma(\omega) \le \psi_I(\omega) \quad (\omega \in \RR^d) ;
\]
under the second, the inequalities between the exponents hold.
\end{enumerate}
\statusA\quad\ledger{A3}\quad\emph{(\textbf{Assignment.} \ledger{A3}(i) and the uniqueness (ii) in
clauses (1) and (2), as in \ref{prop:similarity-orbit-family}(1); clauses (3)--(5) are order
% CHANGED Q-0281 wave 2: blind review W2G-3
arguments and cite nothing of their own: (3) reads \ledger{A3}(ii) through the symmetric form of (1),
and (4)--(5) through (2). The hypothesis that the triplet has drift $0$ and a symmetric \Levy{}
measure is what symmetry of $\mu_\Sigma$ gives through the uniqueness. \textbf{What it is worth.} It is
item~1 of Q-0264: the sandwich of \ref{prop:similarity-orbit-family}(2) for a non-Gaussian part
that is not assumed to be an image of one law, in the \Levy{} order under infinitely divisible
increments and in the order of exponents under positive increments when there is no Gaussian part.
Writing it down found two things. The rank-one limit needs no continuity hypothesis: the \Levy{}
measures along $\Sigma_\varepsilon$ decrease as $\varepsilon \downarrow 0$, so the limit exists by
monotonicity alone. The analogue of the elliptical nodes' extension of $h$ to
$\PSD(d) \setminus \{0\}$ with values in $\PD(d)$ is therefore a hypothesis on the limit's
\emph{size}, and the limit may be $0$ (\ref{prop:similarity-degenerate-family}). And the
self-decomposability of $\nu_\Sigma$ is not used. \textbf{Superpositions.} A finite sum of elliptical
parts, $L_\Sigma = \sum_jL^{(j)}_{h_j(\Sigma)}$, falls under \ref{prop:similarity-orbit-family}(2)
term by term only when each term is itself Loewner-monotone and has $h_j(e_1e_1^{\tp}) \in \PD(d)$;
then the exponent spread of the sum is bounded
(\ref{prop:similarity-general-bounded-shape}(1)). Infinitely divisible increments of the sum order
only the sum, and \ref{prop:similarity-degenerate-family}(3) is a two-term sum, with a nonzero
rank-one limit, whose shape is unbounded. So nothing here assumes a structure for $L_\Sigma$, and
the elliptical nodes do not cover sums. \textbf{The axiomatic families.} For the kernels from scale zero of a (similarities, Loewner)-cascade family the standing hypotheses (a symmetric infinitely
divisible law with drift $0$, similarity covariance) and the first hypothesis of (4) are theorems,
\ref{prop:similarity-cascade-families}(1). \textbf{Not claimed}: any order between the parts under
positive increments when there is a Gaussian part, where neither order is available
(\ref{prop:similarity-orbit-family}'s annotation), for kernel families that are not cascades.)}
\end{proposition}

\begin{proof}
\leanok
\emph{Clause (1).} Let $A := \sqrt c\,Q$. The reflection $(-I)_*$ carries a triplet $(C, L, 0)$ to
$(C, (-I)_*L, 0)$, the truncation $\mathbf 1_D$ being even; as $(-I)_*\mu_\Sigma = \mu_\Sigma$, the uniqueness
\ledger{A3}(ii) makes $L_\Sigma$ symmetric. Since $L_\Sigma$ is symmetric and the drift is $0$, the
imaginary part of the integrand of \ledger{A3}(i) integrates to $0$, as in the proof of
\ref{lem:elliptical-stable}(2), and
$\widehat\mu_\Sigma(\omega) = \exp[-\tfrac12\langle\omega,C(\Sigma)\omega\rangle - \psi_\Sigma(\omega)]$.
So $\widehat{A_*\mu_\Sigma}(\omega) = \widehat\mu_\Sigma(A^{\tp}\omega) =
\exp[-\tfrac12\langle\omega,AC(\Sigma)A^{\tp}\omega\rangle - \int(1 - \cos\langle\omega,y\rangle)\,
(A_*L_\Sigma)(dy)]$. The measure $A_*L_\Sigma$ is symmetric, and
$\int(|y|^2 \wedge 1)\,A_*L_\Sigma(dy) = \int(c|x|^2 \wedge 1)\,L_\Sigma(dx) \le \max(1,c)\int(|x|^2 \wedge
1)\,L_\Sigma(dx)$, so the same collapse read backwards shows that this is the transform of the law
with triplet $(AC(\Sigma)A^{\tp}, A_*L_\Sigma, 0)$. That law is $\mu_{cQ\Sigma Q^{\tp}}$, and the
uniqueness \ledger{A3}(ii) gives the first two identities; the third is the substitution
$y = Ax$. For $c = 1$ and $\Sigma = I$, $QIQ^{\tp} = I$.

\emph{Clause (2).} Let $(A_\rho, \nu_\rho, \gamma_\rho)$ be the triplet of $\rho$
(\ledger{A3}(i)). The exponents of a convolution add, so $\mu_\Sigma * \rho$ has the triplet
$(C(\Sigma) + A_\rho, L_\Sigma + \nu_\rho, \gamma_\rho)$, and by \ledger{A3}(ii)
$C(\Sigma') = C(\Sigma) + A_\rho$ and $L_{\Sigma'} = L_\Sigma + \nu_\rho$. Integrating
$1 - \cos\langle\omega,\cdot\rangle \ge 0$ gives $\psi_\Sigma \le \psi_{\Sigma'}$.

\emph{Clause (3).} With $C \equiv 0$, $\widehat\mu_\Sigma = e^{-\psi_\Sigma}$ is positive, so
$\widehat\rho = \widehat\mu_{\Sigma'}/\widehat\mu_\Sigma = e^{-(\psi_{\Sigma'} - \psi_\Sigma)}$;
$\rho(\RR^d) = \widehat\rho(0) = 1$, and $|\widehat\rho| \le 1$ gives $\psi_{\Sigma'} - \psi_\Sigma \ge 0$.

\emph{Clause (4).} For $0 < \varepsilon \le \varepsilon' \le 1$, $\Sigma_\varepsilon \le
\Sigma_{\varepsilon'}$, so by (2) $\varepsilon \mapsto L_{\Sigma_\varepsilon}(B)$ is nondecreasing, and it
is bounded by $L_{\Sigma_1}(B) = L_I(B)$, which is finite when $B \subseteq \{|x| \ge \delta\}$ for some
$\delta > 0$. So the limit exists. On the Borel sets of $\{|x| \ge \delta\}$ it is countably
additive: for disjoint $B_n$ with union $B$, $L_{\Sigma_\varepsilon}(B) = \sum_nL_{\Sigma_\varepsilon}(B_n)$,
each term converges and is at most $L_I(B_n)$, and $\sum_nL_I(B_n) < \infty$, so the limit of the sum
is the sum of the limits. For an arbitrary Borel set $B \subseteq \RR^d \setminus \{0\}$ and $\delta > 0$
put $B_\delta := \{x \in B : |x| \ge \delta\}$ and
$L_*(B) := \lim_{\delta\downarrow0}L_*(B_\delta)$, an increasing limit of measures and
so a measure, which agrees with the limit above on sets bounded away from the origin. Since
$L_*(B_\delta) \le L_{\Sigma_\varepsilon}(B)$ for every $\delta$, $L_* \le
L_{\Sigma_\varepsilon}$, and in particular $\int(|x|^2 \wedge 1)\,dL_* \le \int(|x|^2 \wedge 1)\,dL_I <
\infty$. The exponents $\psi_{\Sigma_\varepsilon}(\omega) \ge 0$ are nondecreasing in $\varepsilon$, by (2)
under the first hypothesis and by (3) under the second, so $\psi_*$ exists. If $Re_1 = \pm e_1$, then
$R\Sigma_\varepsilon R^{\tp} = \varepsilon I + (1-\varepsilon)Re_1(Re_1)^{\tp} = \Sigma_\varepsilon$, and (1)
with $c = 1$ gives $R_*L_{\Sigma_\varepsilon} = L_{\Sigma_\varepsilon}$ and
$\psi_{\Sigma_\varepsilon}(R^{\tp}\omega) = \psi_{\Sigma_\varepsilon}(\omega)$; as $R$ maps sets bounded
away from the origin to such sets, both identities pass to the limit.

\emph{Clause (5).} Write $v := Qe_1$ and $\Sigma = vv^{\tp} + \sum_{i\ge2}\sigma_iw_iw_i^{\tp}$ with
$(v, w_2, \dots, w_d)$ orthonormal and $\sigma_i \ge \lambda_{\min}(\Sigma)$. For
$0 < \varepsilon \le \lambda_{\min}(\Sigma)$, $Q\Sigma_\varepsilon Q^{\tp} = \varepsilon I + (1-\varepsilon)vv^{\tp}$,
so $\Sigma - Q\Sigma_\varepsilon Q^{\tp} = \sum_{i\ge2}(\sigma_i - \varepsilon)w_iw_i^{\tp} \ge 0$, and
$\Sigma \le I$. Under the first hypothesis, (2) and (1) give
$Q_*L_{\Sigma_\varepsilon} = L_{Q\Sigma_\varepsilon Q^{\tp}} \le L_\Sigma \le L_I$ and
$\psi_{\Sigma_\varepsilon}(Q^{\tp}\omega) \le \psi_\Sigma(\omega) \le \psi_I(\omega)$, and by (4)
$Q_*L_* \le Q_*L_{\Sigma_\varepsilon}$ and $\psi_* \le \psi_{\Sigma_\varepsilon}$. Under the second, (3)
replaces (2) for the exponents.
\end{proof}

% not in the draft: Q-0264 item 2, the positive half. Three anisotropy functionals; the exponent
% spread is the one chosen for A' (the author, 2026-10-03) (it is the eccentricity's analogue for a
% general part, and the only one positive increments see). Each bound is the sandwich read in one
% functional, under a hypothesis that the rank-one limit is comparable to the isotropic member.
% Title until Q-0284: "conjecture A' under a comparable rank-one limit" (the names dropped, the
% author, 2026-10-04).
% CHANGED Q-0281: typed (the statement skeleton). The identity
% M_s((cR)_*L) = c^2 R M_{s/c}(L) R^T of clause (1) is typed for a Levy measure L, where the entries
% of the moment matrix are finite; for a measure with int (|x|^2 ^ s^2) dL infinite the entries are
% infinite or undefined, and the Lean integral is junk there. The text is unchanged (a narrowing
% recorded in records/formalization/SKELETON.md, dispatch 4).
% CHANGED Q-0281 wave 2, group G: proved (AffineHemigroup.SimilarityBounded; clause (1) Lean core,
% the orbit's bound read on the exponent of the profile itself, so the law of
% prop:orbit-step-cone(1) is not needed; clauses (2)-(4) A3(i)-(ii) through
% prop:similarity-general-sandwich; guarded). The proof of record follows the Lean and is unchanged.
% CHANGED Q-0327: notation only (whole-paper review flag 15): the rotation is $Q$ (and $Q'$ for a second
% one), no longer $R$, which is the dilation ratio of prop:orbit-step-cone.
% CHANGED round-1 review (Gemini), required change 9: notation only, completing Q-0327's rename. The
% statement's clause (1) had (cR) for (cQ) three times; the proof's clause (1) had the garbled
% conjugation Q -> Q'^{-1}RR' (now Q -> Q'^{-1}QQ' for Q' in O(d)) and R x x^T Q^T (now Q x x^T Q^T);
% its clause (4) had Re_1 = +-e_1 for Q'e_1 = +-e_1. Read with R = Q every identity was already true;
% nothing claimed changes, and the Lean statement (which binds R throughout) is untouched.
\begin{proposition}[bounded shape under a comparable rank-one limit]
\label{prop:similarity-general-bounded-shape}
\uses{prop:similarity-general-sandwich,prop:orbit-step-cone,lem:elliptical-stable,lem:affine-gauge,thm:cited-A3}
\lean{AffineHemigroup.anisotropy_functionals, AffineHemigroup.sim_bounded_shape_levy,
AffineHemigroup.sim_bounded_shape_exp, AffineHemigroup.sim_bounded_shape_moments}\leanok
For a measure $L$ on $\RR^d \setminus \{0\}$, a function $\psi \ge 0$ on $\RR^d$, $s > 0$ and
$H \in \PD(d)$, put
\[
  \operatorname{rs}(L) := \inf\{K \ge 1 : Q_*L \le K\,L \text{ for every } Q \in \Orth(d)\},
\]
\[
  \operatorname{es}(\psi) := \inf\{K \ge 1 : \psi(Q\omega) \le K\,\psi(\omega) \text{ for every } Q \in \Orth(d),\ \omega \in \RR^d\},
\]
\[
  \mathcal M_s(L) := \int(|x|^2 \wedge s^2)\,\frac{xx^{\tp}}{|x|^2}\,L(dx),
  \qquad T_s(L) := \tr\mathcal M_s(L) = \int(|x|^2 \wedge s^2)\,L(dx),
\]
\[
  \operatorname{ecc}(H) := \frac{\lambda_{\max}(H)}{\lambda_{\min}(H)} ,
\]
the \emph{rotational spread}, the \emph{exponent spread}, the \emph{moment matrix at scale $s$} and
the eccentricity, the infimum of the empty set being $\infty$, and $\operatorname{cond}M := \lambda_{\max}(M)/\lambda_{\min}(M)
\in [1,\infty]$ for $M \in \PSD(d) \setminus \{0\}$. Let the family be as in
\ref{prop:similarity-general-sandwich}.
\begin{enumerate}
\item \emph{(The functionals.)} For $c > 0$ and $Q \in \Orth(d)$,
$\operatorname{rs}((cQ)_*L) = \operatorname{rs}(L)$, $\operatorname{es}(\psi(cQ\,\cdot)) = \operatorname{es}(\psi)$
and $\mathcal M_s((cQ)_*L) = c^2Q\,\mathcal M_{s/c}(L)\,Q^{\tp}$. $\operatorname{rs}(L) = 1$ exactly when $L$ is rotation
invariant, $\operatorname{es}(\psi) = 1$ exactly when $\psi$ is, and $\operatorname{es}(\psi_L) \le \operatorname{rs}(L)$ for the
exponent $\psi_L(\omega) := \int(1 - \cos\langle\omega,x\rangle)\,L(dx)$; the exponent spread of a sum is at
most the largest exponent spread of its terms. For the orbit $(\nu_H)$ of a profile as in
\ref{prop:orbit-step-cone}(1), $\operatorname{es}(\psi_{\nu_H}) \le 4\operatorname{ecc}(H)$; for the stable orbit,
$\operatorname{es}(\psi_{S_\alpha(H)}) = \operatorname{ecc}(H)^{\alpha/2}$.
\item \emph{(\Levy{} level.)} Suppose every Loewner step in $\PD(d)$ has an infinitely divisible
increment and $L_* \ge \theta L_I$ for some $\theta \in (0,1]$. Then for every $\Sigma \in \PD(d)$, with
$t := \lambda_{\max}(\Sigma)$,
\[
  \theta\,L_{tI} \le L_\Sigma \le L_{tI} ,
  \qquad
  \operatorname{rs}(L_\Sigma) \le 1/\theta ,
  \qquad
  \operatorname{es}(\psi_\Sigma) \le 1/\theta .
\]
\item \emph{(Exponent level.)} Suppose every Loewner step in $\PD(d)$ has an infinitely divisible
increment, or that $C \equiv 0$ and every Loewner step has a positive finite increment; and that
$\psi_* \ge \theta\psi_I$ for some $\theta \in (0,1]$. Then $\theta\psi_{tI} \le \psi_\Sigma \le \psi_{tI}$ and
$\operatorname{es}(\psi_\Sigma) \le 1/\theta$ for every $\Sigma \in \PD(d)$, $t := \lambda_{\max}(\Sigma)$.
\item \emph{(Moments, scale by scale.)} Suppose every Loewner step in $\PD(d)$ has an infinitely
divisible increment, and $L_I \ne 0$. For $s > 0$ the limit $\mathcal M_s^* := \lim_{\varepsilon\downarrow0}\mathcal M_s(L_{\Sigma_\varepsilon})$
exists and is $\diag(a_s, b_s, \dots, b_s)$; put $\theta_s := d\min(a_s,b_s)/T_s(L_I) \in [0,1]$. Then for every
$\Sigma \in \PD(d)$ and $s > 0$, with $t := \lambda_{\max}(\Sigma)$,
$\operatorname{cond}\mathcal M_s(L_\Sigma) \le 1/\theta_{s/\sqrt t}$; so
$\operatorname{cond}\mathcal M_s(L_\Sigma) \le 1/\inf_{r>0}\theta_r$ for all $\Sigma$ and $s$.
\end{enumerate}
\statusA\quad\ledger{A3}\quad\emph{(\textbf{Assignment.} \ledger{A3} through
\ref{prop:similarity-general-sandwich}(1)--(2) in clauses (2)--(4), and through
\ref{prop:orbit-step-cone}(1) for the orbit named in clause (1); the rest is the order of measures,
of functions and of matrices. \textbf{The functional.} Q-0264 asks for one in which to decide
conjecture A$'$; the exponent spread is the one chosen, and is A$'$'s functional of record (the
author, 2026-10-03). It is defined for every symmetric
infinitely divisible law, it is the condition number for a Gaussian ($\operatorname{es}(\langle\omega,C\omega\rangle)
= \operatorname{cond}C$), it is $\operatorname{ecc}(H)^{\alpha/2}$ on the stable orbit and at most $4\operatorname{ecc}(H)$ on every
elliptical orbit of a self-decomposable profile as in \ref{prop:orbit-step-cone}(1) (the constant
uses that $k$ is nonincreasing), so it extends the eccentricity that \ref{prop:similarity-stable-family}(2) and
\ref{prop:similarity-orbit-family}(2) bound, and it equals $1$ exactly for the rotation-invariant parts
(the exponent determines the law). It is also the only one of the three that positive increments see
when there is no Gaussian part (clause (3)), which is Q-0264's item~4. The rotational spread is the
two-sided ratio of the polar profiles across directions, uniformly in the radius, that the brief
names; it dominates the exponent spread, $\operatorname{es}(\psi_L) \le \operatorname{rs}(L)$. The moment
matrices are read at each scale separately. No comparison between the exponent spread and the moment
matrices is claimed.
\textbf{What it is worth.} Conjecture A$'$ holds, in each functional, as soon as the rank-one limit is
comparable to the isotropic member in that functional; the hypotheses are sufficient, not necessary,
since the shape-blind family $\nu_\Sigma := ((\det\Sigma)^{1/(2d)})_*\nu_I$, with $\nu_I$ rotation invariant and
self-decomposable, has $L_* = 0$ and spread $1$. For the elliptical stable parts of
\ref{prop:similarity-stable-family}(2), whose gauge is continuous at $e_1e_1^{\tp}$ with
$h(e_1e_1^{\tp}) \in \PD(d)$, the hypothesis of (2) holds, $L_*$ and $L_I$ having densities homogeneous of
the same degree, positive and continuous on the sphere, whose ratio is bounded below; so under
infinitely divisible increments (2) recovers that clause with another constant. Without such a
hypothesis A$'$ fails in all three functionals (\ref{prop:similarity-degenerate-family}), and the
rotational spread is not equivalent to the other two (\ref{prop:similarity-line-family}). Conjecture
A$'$ as restated (the author, 2026-10-03) is what this node and \ref{prop:similarity-general-index-zero}
prove: the shape is bounded, in the exponent spread, whenever the rank-one limit is comparable to the
isotropic member, and at index $0$ the part is shape-blind at the origin. \textbf{Not claimed}: a converse, and a uniform moment
bound without $\inf_r\theta_r > 0$ (\ref{prop:similarity-degenerate-family}(3) has $\theta_r > 0$ for every
$r$ and no uniform bound).)}
\end{proposition}

\begin{proof}
\leanok
\emph{Clause (1).} For $Q' \in \Orth(d)$, conjugation $Q \mapsto Q'^{-1}QQ'$ is a bijection of $\Orth(d)$, so
$Q_*(Q'_*L) \le K\,Q'_*L$ for every $Q$ exactly when $Q_*L \le K\,L$ for every $Q$; dilations commute with
rotations; the same holds for $\psi$; and
$\mathcal M_s((cQ)_*L) = \int(c^2|x|^2 \wedge s^2)\,\frac{Qxx^{\tp}Q^{\tp}}{|x|^2}\,L(dx) = c^2Q\,\mathcal M_{s/c}(L)\,Q^{\tp}$.
If $Q_*L \le L$ for every $Q$, then $L = Q_*(Q^{-1})_*L \le Q_*L$, so $Q_*L = L$; the same argument
applies to $\psi$; conversely invariance gives $K = 1$. If $Q_*L \le K\,L$ for every $Q$, then
$\psi_L(Q^{\tp}\omega) = \int(1 - \cos\langle\omega,y\rangle)\,(Q_*L)(dy) \le K\,\psi_L(\omega)$. If
$\psi_j(Q\omega) \le K_j\psi_j(\omega)$, then $\sum_j\psi_j(Q\omega) \le \max_jK_j\sum_j\psi_j(\omega)$.
For the orbit, put $g(u) := \psi_{\nu_I}(\sqrt u\,e)$ for a unit vector $e$; by rotation invariance and
$\widehat\nu_H(\omega) = \widehat\nu_I(H^{1/2}\omega)$, $\psi_{\nu_H}(\omega) = g(\langle\omega,H\omega\rangle)$. In
the notation of \ref{prop:orbit-step-cone}(1), the substitution $x = y/a$ gives
$g(a^2u) = \int(1 - \cos\langle\sqrt u\,e,y\rangle)\,k(|y|/a)\,|y|^{-d}\,dy$, nondecreasing in $a > 0$ as $k$
is nonincreasing; and $1 - \cos 2\tau = 2(1 - \cos\tau)(1 + \cos\tau) \le 4(1 - \cos\tau)$ gives
$g(4u) \le 4g(u)$. For $c \ge 1$ take the least $n \ge 0$ with $4^n \ge c$; then
$g(cu) \le g(4^nu) \le 4^ng(u) \le 4c\,g(u)$. For $|\omega| = |\omega'|$, $\langle\omega,H\omega\rangle \le
\operatorname{ecc}(H)\langle\omega',H\omega'\rangle$, so $\psi_{\nu_H}(\omega) \le 4\operatorname{ecc}(H)\,\psi_{\nu_H}(\omega')$. For the stable
orbit $\psi_{S_\alpha(H)}(\omega) = \langle\omega,H\omega\rangle^{\alpha/2}$ (\ref{lem:elliptical-stable}(2)), whose
largest ratio on a sphere is $\operatorname{ecc}(H)^{\alpha/2}$, attained at unit eigenvectors.

\emph{Clause (2).} By \ref{prop:similarity-general-sandwich}(1), $L_\Sigma = (\sqrt t)_*L_{\Sigma/t}$ and
$L_{tI} = (\sqrt t)_*L_I$, and by (1) the spreads do not change under $(\sqrt t)_*$; so let $t = 1$.
With $Q$ as in \ref{prop:similarity-general-sandwich}(5) and $L_I$ rotation invariant,
$\theta L_I = \theta\,Q_*L_I \le Q_*L_* \le L_\Sigma \le L_I$. For $Q' \in \Orth(d)$,
$Q'_*L_\Sigma \le Q'_*L_I = L_I \le \theta^{-1}L_\Sigma$, so $\operatorname{rs}(L_\Sigma) \le 1/\theta$, and
$\operatorname{es}(\psi_\Sigma) \le \operatorname{rs}(L_\Sigma)$ by (1).

\emph{Clause (3).} As in (2), let $t = 1$. By \ref{prop:similarity-general-sandwich}(5), under either
hypothesis, $\theta\psi_I(\omega) = \theta\psi_I(Q^{\tp}\omega) \le \psi_*(Q^{\tp}\omega) \le \psi_\Sigma(\omega) \le
\psi_I(\omega)$, and for $Q' \in \Orth(d)$, $\psi_\Sigma(Q'\omega) \le \psi_I(Q'\omega) = \psi_I(\omega) \le
\theta^{-1}\psi_\Sigma(\omega)$.

\emph{Clause (4).} If $L \le L'$, then $\int f\,dL \le \int f\,dL'$ for $f \ge 0$, so
$\langle v,\mathcal M_s(L)v\rangle \le \langle v,\mathcal M_s(L')v\rangle$ for every $v$: $\mathcal M_s$ is monotone in the
Loewner order. By \ref{prop:similarity-general-sandwich}(2), $\varepsilon \mapsto \mathcal M_s(L_{\Sigma_\varepsilon})$ is
nondecreasing and nonnegative, so every quadratic form $\langle v,\mathcal M_s(L_{\Sigma_\varepsilon})v\rangle$ converges as
$\varepsilon \downarrow 0$, and by polarization so does the matrix. By (1) and
\ref{prop:similarity-general-sandwich}(1), $Q'\,\mathcal M_s(L_{\Sigma_\varepsilon})Q'^{\tp} = \mathcal M_s(L_{\Sigma_\varepsilon})$ whenever
$Q'e_1 = \pm e_1$, so $\mathcal M_s^*$ commutes with these $Q'$. The reflection $I - 2e_1e_1^{\tp}$ makes the entries
$(\mathcal M_s^*)_{1k}$, $k \ne 1$, vanish, and on $e_1^\perp$ the matrix commutes with the reflections and the
permutations of $e_2,\dots,e_d$, which makes it a multiple of the identity there, as in the proof of
\ref{lem:affine-gauge}. Similarly $\mathcal M_s(L_I)$ commutes with $\Orth(d)$, so it is $\tfrac1dT_s(L_I)\,I$. Now let
$t = 1$ (the general case follows by $\mathcal M_s((\sqrt t)_*L) = t\,\mathcal M_{s/\sqrt t}(L)$, from (1)). By
\ref{prop:similarity-general-sandwich}(5) and its proof, $Q_*L_{\Sigma_\varepsilon} \le L_\Sigma \le L_I$ for
$\varepsilon \le \lambda_{\min}(\Sigma)$, so
$Q\,\mathcal M_s^*\,Q^{\tp} \le Q\,\mathcal M_s(L_{\Sigma_\varepsilon})\,Q^{\tp} \le \mathcal M_s(L_\Sigma) \le \tfrac1dT_s(L_I)\,I$. Hence
$\lambda_{\min}(\mathcal M_s(L_\Sigma)) \ge \min(a_s,b_s)$ and $\lambda_{\max}(\mathcal M_s(L_\Sigma)) \le \tfrac1dT_s(L_I)$,
which is the bound (trivial when $\theta_s = 0$).
\end{proof}

% not in the draft: Q-0264 item 2, the negative half. A one-parameter family of gauges between
% the stable orbit (gamma = 0, which has no rank-one step) and the shape-blind gauge det^{1/d} I
% (gamma = 1/d). Its monotonicity is the Loewner-Heinz inequality, proved here from the integral
% representation of t^p. Checked numerically in records/affine/check-aprime-general.py (a), (c).
% CHANGED Q-0281: \uses re-pointed to the split child it reads (the densities n_H of lem:elliptical-stable beside the laws of its (2)).
% CHANGED Q-0281 wave 2, group I: proved (AffineHemigroup.SimilarityDegenerate). (1) and (3) print
% A3(iii), (2) Lean core; the Assignment is narrowed to that spending (the rank-one limits are
% computed directly, so (i)-(ii) are not read).
\begin{proposition}[without a nondegenerate rank-one limit the shape is unbounded]
\label{prop:similarity-degenerate-family}
\uses{lem:elliptical-stable,lem:stable-tail,prop:stable-step-cone,prop:similarity-stable-family,prop:similarity-general-sandwich,prop:similarity-general-bounded-shape,def:affine-cascade-family,lem:elliptical-stable-density,thm:cited-A3}
\lean{AffineHemigroup.degenerate_cascade, AffineHemigroup.degenerate_unbounded,
AffineHemigroup.degenerate_rotation_invariant_limit}\leanok
Let $d \ge 2$, $0 < \alpha < 2$, $\tfrac{1}{d+\alpha} \le \gamma < \tfrac1d$ and $p := 1 - d\gamma$, so that
$0 < p \le \tfrac{\alpha}{d+\alpha}$. For $\Sigma \in \PSD(d)$ put $h(\Sigma) := (\det\Sigma)^\gamma\,\Sigma^p$,
which is $0$ when $\Sigma$ is singular. For $b \ge 0$ put
$\mu_\Sigma := N(0,2b\Sigma) * S_\alpha(h(\Sigma))$, with $S_\alpha$ and $n_H$ as in
\ref{lem:elliptical-stable} and $S_\alpha(0) := \delta_0$; write $L_\Sigma := n_{h(\Sigma)}(x)\,dx$ and
$\psi_\Sigma$ for the \Levy{} measure and the exponent of the non-Gaussian part, $\Sigma \in \PD(d)$, and
use the notation of \ref{prop:similarity-general-bounded-shape}.
\begin{enumerate}
\item \emph{(A cascade.)} $h$ is continuous on $\PSD(d)$, $h(cQ\Sigma Q^{\tp}) = c\,Q\,h(\Sigma)\,Q^{\tp}$ for
$c > 0$ and $Q \in \Orth(d)$, and $n_{h(\Sigma)} \le n_{h(\Sigma')}$ for every $\Sigma \le \Sigma'$ in $\PD(d)$.
For every $\Sigma \le \Sigma'$ in $\PSD(d)$ there is an infinitely divisible probability measure
$\rho_{\Sigma,\Sigma'}$ with $\mu_{\Sigma'} = \mu_\Sigma * \rho_{\Sigma,\Sigma'}$, and the family
$(\rho_{\Sigma,\Sigma'})$ satisfies (K1), (K2) and (K4) of \ref{def:affine-cascade-family}, and (K3) for
every similarity $A = cQ$.
\item \emph{(Unbounded shape.)} For $\Sigma \in \PD(d)$, $\operatorname{ecc}(h(\Sigma)) = \operatorname{ecc}(\Sigma)^p$,
$\operatorname{es}(\psi_\Sigma) = \operatorname{ecc}(\Sigma)^{p\alpha/2}$, $\operatorname{rs}(L_\Sigma) \ge \operatorname{ecc}(\Sigma)^{p\alpha/2}$, and for
every $s > 0$
\[
  \operatorname{cond}\mathcal M_s(L_\Sigma) \ge \frac{G_1}{d\,G_0}\operatorname{ecc}(\Sigma)^{p\alpha/2} ,
\]
\[
  G_0 := \int_{\RR^d}(z_d^2 \wedge 1)\,|z|^{-d-\alpha}\,dz ,
  \quad
  G_1 := \int_{\RR^d}(|z|^2 \wedge 1)\,|z|^{-d-\alpha}\,dz ,
\]
both finite and positive. The rank-one limit of \ref{prop:similarity-general-sandwich}(4) is $L_* = 0$.
\item \emph{(A nonzero rotation-invariant rank-one limit is not enough.)} Let $0 < \beta < \alpha$, let
$S'_\beta(H)$ and $n'_H$ be the law and the \Levy{} density of \ref{lem:elliptical-stable} for the index
$\beta$, with $S'_\beta(0) := \delta_0$, and put $\mu'_\Sigma := \mu_\Sigma * S'_\beta(\tr(\Sigma)\,I)$, with \Levy{} measure
$L'_\Sigma = (n_{h(\Sigma)} + n'_{\tr(\Sigma)I})(x)\,dx$ and exponent $\psi'_\Sigma$. The family $(\mu'_\Sigma)$
has all the properties of (1); its rank-one limit is $L'_* = n'_I(x)\,dx$, nonzero and rotation
invariant; for every $s > 0$, $\theta_s > 0$ in \ref{prop:similarity-general-bounded-shape}(4); and for
every $\Sigma \in \PD(d)$
\[
  \operatorname{es}(\psi'_\Sigma) \ge \operatorname{ecc}(\Sigma)^{p\alpha/2} ,
  \qquad
  \sup_{s>0}\operatorname{cond}\mathcal M_s(L'_\Sigma) \ge \frac{G_1}{d\,G_0}\operatorname{ecc}(\Sigma)^{p\alpha/2} .
\]
\end{enumerate}
\statusA\quad\ledger{A3}\quad\emph{(\textbf{Assignment.} \ledger{A3}(iii) through
\ref{lem:elliptical-stable}(2) and \ref{prop:stable-step-cone}(1), for the laws and the stable
increments of (1) and (3). The rest cites nothing: the rank-one limits are computed directly, the
Loewner--Heinz inequality $\Sigma \le \Sigma' \Rightarrow
\Sigma^p \le \Sigma'^p$, $0 < p < 1$, is proved in the proof from the integral representation of
$\tau^p$, and the moment bound is two substitutions. \textbf{What it is worth.} It settles conjecture
A$'$ in the form the postdoc stated it, ``under similarity covariance with a cascade along the whole
Loewner order the non-Gaussian part stays within a bounded neighbourhood of isotropy'': that is
false, already for elliptical stable parts with infinitely divisible increments, and with or without
a Gaussian part ($b = 0$ is allowed). The hypothesis that \ref{prop:similarity-stable-family}(2) and
\ref{prop:similarity-orbit-family}(2) add, that $h$ extends to rank one with values in $\PD(d)$,
therefore cannot simply be dropped: here $h$ extends continuously by $0$ and the shape is unbounded.
It is not necessary either: the gauge $(\det\Sigma)^{1/d}I$, the endpoint $\gamma = 1/d$, also has
$h(e_1e_1^{\tp}) = 0$, and its spread is $1$. What survives without it
is a trade-off: the proof of \ref{prop:similarity-stable-family}(2) gives
$\operatorname{ecc}(h(\Sigma)) \le (h_0/\lambda_{\max}(h(\Sigma)))^\alpha$ from the step $\Sigma \le \lambda_{\max}(\Sigma)I$
alone, so the part may be eccentric only where it is small, and in the plane at
$\gamma = \tfrac1{2+\alpha}$ this family attains that bound with equality ($h_0 = 1$, and for
$\Sigma = \Sigma_\varepsilon$, $\lambda_{\max}(h) = \varepsilon^\gamma$ and $\operatorname{ecc}(h) = \varepsilon^{-p}$ with $p = \gamma\alpha$). The
family interpolates between the stable orbit, $\gamma = 0$, which has no rank-one step
(\ref{prop:stable-step-cone}(3)), and the shape-blind gauge $(\det\Sigma)^{1/d}I$, $\gamma = \tfrac1d$. The
threshold $\gamma \ge \tfrac1{d+\alpha}$ is sharp for rank-one steps from $I$: for a smaller $\gamma$ the display in
the proof of (1) gives $n_{h(I + \tau e_1e_1^{\tp})}(e_2) = (1+\tau)^{((d+\alpha)\gamma - 1)/2}n_{h(I)}(e_2) <
n_{h(I)}(e_2)$, and \ref{lem:stable-tail} excludes a positive increment. Clause (3) shows that a
nonzero rank-one limit, even a rotation-invariant one, does not give A$'$ in any of the three
functionals: what \ref{prop:similarity-general-bounded-shape} needs is comparability with $L_I$, and
here $n'_I/(n_I + n'_I) \to 0$ at the origin. The scale-by-scale moment bound holds and is not
uniform. \textbf{Checked numerically} (\texttt{records/affine/check-aprime-general.py} (a), (c), output in
\texttt{check-aprime-general-output.txt}): the ordering on 4000 random Loewner steps for each of
$d \in \{2,3\}$, $\alpha \in \{0.5, 1, 1.5\}$ at the threshold and midway, its failure at
$\gamma = 0.9/(d+\alpha)$, $\operatorname{ecc}(h(\Sigma_\varepsilon)) = \varepsilon^{-p}$, the moment bound in $d = 2$, and the
exponent spread of (3). \textbf{Not claimed}: anything at \Levy{} index $0$, where an elliptical part
is shape-blind whatever its gauge does at rank one (\ref{prop:similarity-orbit-family}(3)--(4)); the
least $\gamma$ for other profiles.)}
\end{proposition}

\begin{proof}
\leanok
% CHANGED Q-0281 wave 2, group I: proof rewritten to the Lean route (inversion reverses the order by
% the variational form of <x, A^{-1}x>; the continuity of Sigma^p from the integral representation;
% in (3) the index-beta part of the moments bounded by its trace; the existence of the rank-one
% moment limit by monotonicity).
\emph{The Loewner--Heinz inequality} (\sscite{bhatia1997matrix}, Thm.~V.1.9, p.~115). For $0 < p < 1$ put $\kappa_p := \int_0^\infty u^{p-1}(1+u)^{-1}du$,
finite and positive; the substitution $s = \tau u$ gives
$\tau^p = \kappa_p^{-1}\int_0^\infty\tau(\tau + s)^{-1}s^{p-1}\,ds$ for $\tau \ge 0$, so by the spectral
decomposition
\[
  \langle x,\Sigma^px\rangle = \kappa_p^{-1}\int_0^\infty\bigl(|x|^2 - s\langle x,(\Sigma + sI)^{-1}x\rangle\bigr)s^{p-1}\,ds
  \qquad (\Sigma \in \PSD(d),\ x \in \RR^d).
\]
Inversion reverses the order on $\PD(d)$: for $A \in \PD(d)$ and all $x, y$,
$\langle y - A^{-1}x, A(y - A^{-1}x)\rangle \ge 0$ gives $2\langle x,y\rangle - \langle y,Ay\rangle \le \langle x,A^{-1}x\rangle$, with
equality at $y = A^{-1}x$; so if $A \le B$, then at $y = B^{-1}x$,
$\langle x,B^{-1}x\rangle = 2\langle x,y\rangle - \langle y,By\rangle \le 2\langle x,y\rangle - \langle y,Ay\rangle \le \langle x,A^{-1}x\rangle$. So
$\Sigma \le \Sigma'$ gives $(\Sigma' + sI)^{-1} \le (\Sigma + sI)^{-1}$ for $s > 0$, then $\Sigma^p \le \Sigma'^p$,
and for $\Sigma \in \PD(d)$, $\Sigma'^{-p} \le \Sigma^{-p}$. The same representation gives the continuity of
$\Sigma \mapsto \Sigma^p$ on $\PSD(d)$: the integrand is continuous in $\Sigma$ and, by the spectral decomposition,
at most $|x|^2s^{p-1}\min(1, \tr\Sigma/s)$, which is integrable and bounded uniformly on $\{\tr\Sigma < T\}$, so
$\langle x,\Sigma^px\rangle$ is continuous by dominated convergence, and the entries of $\Sigma^p$ by polarization.

\emph{Clause (1).} $\det h(\Sigma) = (\det\Sigma)^{d\gamma + p} = \det\Sigma$ and
$h(\Sigma)^{-1} = (\det\Sigma)^{-\gamma}\Sigma^{-p}$, so by \ref{eq:stable-levy-density}
\[
  n_{h(\Sigma)}(x) = c_{d,\alpha}^{-1}\,(\det\Sigma)^{((d+\alpha)\gamma - 1)/2}\,\langle x,\Sigma^{-p}x\rangle^{-(d+\alpha)/2}
  \qquad (\Sigma \in \PD(d),\ x \ne 0).
\]
For $\Sigma \le \Sigma'$ in $\PD(d)$, $\det\Sigma \le \det\Sigma'$ (the eigenvalues of
$\Sigma^{-1/2}\Sigma'\Sigma^{-1/2}$ are at least $1$) with a nonnegative exponent, as
$(d+\alpha)\gamma \ge 1$, and $\langle x,\Sigma'^{-p}x\rangle \le \langle x,\Sigma^{-p}x\rangle$ with a negative
exponent; so both factors increase and $n_{h(\Sigma)} \le n_{h(\Sigma')}$. Equivariance:
$\det(cQ\Sigma Q^{\tp})^\gamma = c^{d\gamma}(\det\Sigma)^\gamma$ and $(cQ\Sigma Q^{\tp})^p = c^pQ\Sigma^pQ^{\tp}$, with
$d\gamma + p = 1$. Continuity: $h$ is the product of $(\det\Sigma)^\gamma$, continuous on $\PSD(d)$ as $\gamma > 0$,
and $\Sigma^p$, continuous by the Loewner--Heinz paragraph; at a singular $\Sigma_0$ it is
$0 = h(\Sigma_0)$. For $\Sigma \le \Sigma'$ in
$\PSD(d)$ let $\rho^{\mathrm s}$ be: the law of \ref{prop:stable-step-cone}(1), infinitely divisible with
$S_\alpha(h(\Sigma')) = S_\alpha(h(\Sigma)) * \rho^{\mathrm s}$, when $\Sigma$ is positive definite (then so is
$\Sigma'$), and $S_\alpha(h(\Sigma'))$ when $\Sigma$ is singular. Put
$\rho_{\Sigma,\Sigma'} := N(0,2b(\Sigma' - \Sigma)) * \rho^{\mathrm s}$. Its transform is
$\exp[-b\langle\omega,(\Sigma' - \Sigma)\omega\rangle - \langle\omega,h(\Sigma')\omega\rangle^{\alpha/2} +
\langle\omega,h(\Sigma)\omega\rangle^{\alpha/2}]$ in every case, and the axioms follow from it as in the
proof of \ref{prop:similarity-stable-family}(3): the exponents telescope, they are continuous in
$(\Sigma,\Sigma')$ because $h$ is, $\mu_I = N(0,2bI) * S_\alpha(I) \ne \delta_0$, and for $A = cQ$ the
equivariance of $h$ gives (K3).

\emph{Clause (2).} The eigenvalues of $h(\Sigma)$ are $(\det\Sigma)^\gamma\sigma_i^p$, $\sigma_i$ those of $\Sigma$,
so $\operatorname{ecc}(h(\Sigma)) = \operatorname{ecc}(\Sigma)^p$, and \ref{prop:similarity-general-bounded-shape}(1) gives
$\operatorname{es}(\psi_\Sigma) = \operatorname{ecc}(h(\Sigma))^{\alpha/2}$ and $\operatorname{rs}(L_\Sigma) \ge \operatorname{es}(\psi_\Sigma)$. For the moments let
$H \in \PD(d)$ with unit eigenvectors $v_\pm$ for its eigenvalues $\eta_+ := \lambda_{\max}(H)$ and
$\eta_- := \lambda_{\min}(H)$, and substitute $x = H^{1/2}y$, under which
$n_H(x)\,dx = c_{d,\alpha}^{-1}|y|^{-d-\alpha}\,dy$ and $\langle v_\pm,x\rangle = \eta_\pm^{1/2}\langle v_\pm,y\rangle$.
Since $(\tau \wedge s^2)\,a/\tau \le a \wedge s^2$ for $0 \le a \le \tau$,
\[
  \langle v_-,\mathcal M_s(n_H)v_-\rangle \le \int(\langle v_-,x\rangle^2 \wedge s^2)\,n_H(x)\,dx
  = c_{d,\alpha}^{-1}\int(\eta_-\langle v_-,y\rangle^2 \wedge s^2)\,|y|^{-d-\alpha}\,dy
  = c_{d,\alpha}^{-1}\eta_-^{\alpha/2}s^{2-\alpha}G_0 ,
\]
by a rotation taking $v_-$ to $e_d$ and the dilation $y = s\eta_-^{-1/2}z$. Since $|x|^2 \le \eta_+|y|^2$ and
$\tau \mapsto (\tau \wedge s^2)/\tau$ is nonincreasing,
\[
  \langle v_+,\mathcal M_s(n_H)v_+\rangle \ge c_{d,\alpha}^{-1}\int(\eta_+|y|^2 \wedge s^2)\,\frac{\langle v_+,y\rangle^2}{|y|^2}\,|y|^{-d-\alpha}\,dy
  = \frac{c_{d,\alpha}^{-1}}{d}\,\eta_+^{\alpha/2}s^{2-\alpha}G_1 ,
\]
the directional factor averaging to $\tfrac1d$ over the sphere and the rest being radial. The ratio of the
two Rayleigh quotients is at most $\operatorname{cond}\mathcal M_s(n_H)$, which gives the bound with
$H = h(\Sigma)$. $G_0$ and $G_1$ are finite because the integrands are at most $|z|^{2-d-\alpha}$ near
$0$ and $|z|^{-d-\alpha}$ near infinity, and positive. Finally, for $x = (x_1, x')$ with $x' \ne 0$, the
display of clause (1) with $\det\Sigma_\varepsilon = \varepsilon^{d-1} \le 1$ gives
$n_{h(\Sigma_\varepsilon)}(x) \le c_{d,\alpha}^{-1}(x_1^2 + \varepsilon^{-p}|x'|^2)^{-(d+\alpha)/2} \to 0$; the set
$\{x' = 0\}$ is null, and $n_{h(\Sigma_\varepsilon)} \le n_{h(I)}$, integrable on every set bounded away from
the origin, so $L_{\Sigma_\varepsilon}(B) \to 0$ by dominated convergence.

\emph{Clause (3).} The gauge $\tr(\Sigma)I$ is equivariant and Loewner-monotone, and
$n'_{\tau I}(x) = c_{d,\beta}^{-1}\tau^{\beta/2}|x|^{-d-\beta}$ increases with $\tau$, so
\ref{prop:stable-step-cone}(1) (for the index $\beta$) gives the second factor's increments, and the
proof of (1) applies to the convolution. Along $\Sigma_\varepsilon$, $\tr\Sigma_\varepsilon \downarrow 1$, so
$L'_{\Sigma_\varepsilon} \ge n'_I(x)\,dx$, and with (2), $L'_* = n'_I(x)\,dx$, as
$n'_{\tau I} = \tau^{\beta/2}n'_I$. The moment matrices $\mathcal M_s(L'_{\Sigma_\varepsilon})$ decrease as $\varepsilon \downarrow 0$
($L'$ is monotone and $\Sigma_\varepsilon$ decreases), so $\mathcal M_s^*$ exists, and
$\mathcal M_s^* \ge \mathcal M_s(n'_I) = \tfrac1dT_s(n'_I)\,I$ by rotation invariance, so $\theta_s \ge
T_s(n'_I)/T_s(L'_I) > 0$. For the exponent spread let $\eta_+$ and $\eta_-$ be the largest and the least
eigenvalue of $h(\Sigma)$, $v_\pm$ unit eigenvectors for them, and $\tau := \tr\Sigma$; as $r \to \infty$,
\[
  \frac{\psi'_\Sigma(rv_+)}{\psi'_\Sigma(rv_-)} = \frac{\eta_+^{\alpha/2}r^\alpha + \tau^{\beta/2}r^\beta}{\eta_-^{\alpha/2}r^\alpha + \tau^{\beta/2}r^\beta}
  \to \Bigl(\frac{\eta_+}{\eta_-}\Bigr)^{\alpha/2} = \operatorname{ecc}(\Sigma)^{p\alpha/2} ,
\]
since $\beta < \alpha$; taking an orthogonal $R$ with $Rv_- = v_+$ and $\omega = rv_-$, every $K$ with
$\psi'_\Sigma(R\omega) \le K\psi'_\Sigma(\omega)$ for all $R$ and $\omega$ is at least this limit. For the moments, $\mathcal M_s(L'_\Sigma) = \mathcal M_s(n_{h(\Sigma)}) + \mathcal M_s(n'_{\tau I})$,
the bounds of (2), and $\langle v_-,\mathcal M_s(n'_{\tau I})v_-\rangle \le T_s(n'_{\tau I}) =
c_{d,\beta}^{-1}\tau^{\beta/2}s^{2-\beta}G'_1$, with $G'_1$ the constant $G_1$ for the index $\beta$ (the substitution
$x = \tau^{1/2}y$ and the dilation of (2)), give
\[
  \frac{\langle v_+,\mathcal M_s(L'_\Sigma)v_+\rangle}{\langle v_-,\mathcal M_s(L'_\Sigma)v_-\rangle}
  \ge \frac{d^{-1}c_{d,\alpha}^{-1}\eta_+^{\alpha/2}G_1}{c_{d,\alpha}^{-1}\eta_-^{\alpha/2}G_0 + c_{d,\beta}^{-1}\tau^{\beta/2}G'_1s^{\alpha-\beta}} ,
\]
which tends to $\frac{G_1}{dG_0}\operatorname{ecc}(\Sigma)^{p\alpha/2}$ as $s \downarrow 0$.

Every step is machine-checked (\texttt{AffineHemigroup.SimilarityDegenerate}); the cited interface is
\ledger{A3}(iii), through \ref{lem:elliptical-stable}(2) and \ref{prop:stable-step-cone}(1), for the stable
increments of (1) and (3).
\end{proof}

% not in the draft: Q-0264 items 2 and 3. The non-elliptical example: a Levy measure laid along
% lines, the one-dimensional profile in direction u dilated by <u, Sigma u>^{1/2}. It separates the
% rotational spread from the other two functionals and refutes exact shape-blindness at index 0 for general parts. Checked
% numerically in records/affine/check-aprime-general.py (b).
% CHANGED Q-0281: blind review R4-1: the tag names line_family_levy and
% line_family_levy_not_shape_blind, the measure-level claims of (1) and (3) from the hypotheses on k
% alone (the cascade and the laws' clause keep the laws nu_Sigma as hypotheses, where a Levy triplet
% would otherwise have carried "L_Sigma is a Levy measure" as an assumption).
% CHANGED Q-0281 wave 2, group I: proved (AffineHemigroup.SimilarityLine). The measures, the
% functionals and the measure half of (3) print Lean core; the cascade A3(iii); the laws of (3)
% A3(ii)-(iii). The Assignment is narrowed to that spending: clause (2) reads no law.
\begin{proposition}[a part laid along lines]
\label{prop:similarity-line-family}
\uses{prop:similarity-general-sandwich,prop:similarity-general-bounded-shape,def:affine-cascade-family,thm:cited-A3}
\lean{AffineHemigroup.line_family_levy, AffineHemigroup.line_family_cascade,
AffineHemigroup.line_family_functionals, AffineHemigroup.line_family_levy_not_shape_blind,
AffineHemigroup.line_family_not_shape_blind}\leanok
Let $d \ge 2$, let $k : (0,\infty) \to [0,\infty)$ be nonincreasing and not identically $0$ with
$\int_0^1rk(r)\,dr + \int_1^\infty k(r)\,dr/r < \infty$, let $\ell(dr) := k(|r|)\,dr/|r|$ on
$\RR \setminus \{0\}$, and let $\overline\sigma$ be the normalised surface measure on $S^{d-1}$. For
$\Sigma \in \PSD(d)$ and $q_\Sigma(u) := \langle u,\Sigma u\rangle$ let $L_\Sigma$ be the measure on
$\RR^d \setminus \{0\}$ with
\[
  \int f\,dL_\Sigma := \int_{S^{d-1}}\overline\sigma(du)\int_{\RR}f\bigl(r\,q_\Sigma(u)^{1/2}u\bigr)\,\ell(dr)
\]
for Borel $f \ge 0$ on $\RR^d$ with $f(0) = 0$, and for $b \ge 0$ let $\mu_\Sigma := N(0,2b\Sigma) * \nu_\Sigma$,
$\nu_\Sigma$ the law with triplet $(0, L_\Sigma, 0)$.
\begin{enumerate}
\item \emph{(A cascade.)} $L_\Sigma$ is a symmetric \Levy{} measure; for $\Sigma \in \PD(d)$ it is the image
of $\overline\sigma(du)\,2k(s\,q_\Sigma(u)^{-1/2})\,ds/s$ under $(u,s) \mapsto su$, a polar profile nonincreasing in
$s$ and at most $2k(0{+})$; $L_{cQ\Sigma Q^{\tp}} = (\sqrt c\,Q)_*L_\Sigma$; and $L_\Sigma \le L_{\Sigma'}$ for
$\Sigma \le \Sigma'$. Hence for every $\Sigma \le \Sigma'$ in $\PSD(d)$ there is an infinitely divisible
probability measure $\rho_{\Sigma,\Sigma'}$ with $\mu_{\Sigma'} = \mu_\Sigma * \rho_{\Sigma,\Sigma'}$, and the family
$(\rho_{\Sigma,\Sigma'})$ satisfies (K1), (K2), (K4) of \ref{def:affine-cascade-family} and (K3) for every
similarity.
\item \emph{(Three functionals.)} For every $\Sigma \in \PD(d)$ and $s > 0$,
$\operatorname{es}(\psi_\Sigma) \le 4d(d+2)$ and $\operatorname{cond}\mathcal M_s(L_\Sigma) \le d+2$, while
$\operatorname{rs}(L_{\Sigma_\varepsilon}) \to \infty$ as $\varepsilon \downarrow 0$. The rank-one limit $L_*$ is at least the
image of $\overline\sigma(du)\,2k(s/|u_1|)\,ds/s$ (the set $u_1 = 0$ is $\overline\sigma$-null), so it
charges every open cone.
\item \emph{(Not shape-blind at index $0$.)} If $\Sigma \in \PD(d)$ is not a multiple of $I$, then $L_\Sigma$,
and so $\nu_\Sigma$, is not rotation invariant. This holds in particular when $k(0{+}) < \infty$, so that
every polar profile of every $L_\Sigma$ is bounded at the origin.
\end{enumerate}
\statusA\quad\ledger{A3}\quad\emph{(\textbf{Assignment.} \ledger{A3}(iii) for the laws $\nu_\Sigma$ and the
increments, whose triplets are $(0, L_\Sigma, 0)$ and $(2b(\Sigma' - \Sigma), L_{\Sigma'} - L_\Sigma, 0)$, and (ii), the
uniqueness of the triplet, in (3), where $Q_*\nu_\Sigma = \nu_\Sigma$ gives $Q_*L_\Sigma = L_\Sigma$ ((iii) makes
$\nu_\Sigma$ infinitely divisible, as (ii) asks). The claims about the measures $L_\Sigma$ in (1) and (3), and
clause (2), cite nothing. \textbf{What it is worth.} It is the family beyond
elliptical parts that Q-0264 asks for, and it is monotone across turning steps: the \Levy{} measure is
laid along lines, and in direction $u$ it is the one-dimensional profile dilated by $q_\Sigma(u)^{1/2}$, so
its exponent is $\psi_\Sigma(\omega) = \int\overline\sigma(du)\,\psi_\ell(q_\Sigma(u)^{1/2}\langle\omega,u\rangle)$. Clause (2)
shows that the rotational spread of \ref{prop:similarity-general-bounded-shape} is not equivalent to
the other two functionals: the
exponent and the moments stay within a bounded neighbourhood of isotropy, and the rotational spread,
the ratio of polar profiles across directions, does not, although the rank-one limit is nonzero and
charges every direction. So a rank-one limit charging every direction is not enough for A$'$ in the
strongest functional; comparability with $L_I$ is what \ref{prop:similarity-general-bounded-shape}(2)
uses, and $k(s/|u_1|)/k(s) \to 0$ near the equator. Clause (3) answers Q-0264's item~3 in the negative:
a general part with polar profiles bounded at the origin need not be shape-blind, even under infinitely
divisible increments. What survives at index $0$ is \ref{prop:similarity-general-index-zero}: for this
family the index at the origin is $2k(0{+})\,\overline\sigma$ for every $\Sigma$, and the shape is carried at
positive radii, by the dilation $q_\Sigma(u)^{1/2}$ of each line. \textbf{Checked numerically}
(\texttt{records/affine/check-aprime-general.py} (b)): for $d = 2$ and $k = \mathbf 1_{(0,1)}$, the largest
condition number over 300 random $\Sigma$ and 13 scales is $2.998$, against the bound $4$. \textbf{Not
claimed}: the value of the rotational spread for a given $\Sigma$, and the asymptotics of $\psi_\Sigma$ at
infinity, hypothesis (D) of \ref{prop:similarity-general-index-zero}(3), which would need a Dini
condition on $k$ at the origin.)}
\end{proposition}

\begin{proof}
\leanok
% CHANGED Q-0281 wave 2, group I: proof rewritten to the Lean route (the fourth moments of
% sigma-bar from rotation invariance; clause (2) at Sigma itself, without the sandwich and without
% any law; fixed caps for the rotational spread; open cones through sectors; clause (3) by an
% iteration at almost-every points in place of right limits; the difference of the Levy measures
% built shell by shell).
% CHANGED Q-0281 wave 2: blind review W2I-5: the facts about psi_ell stand before clause (1), which
% uses them in (K2) and (K4); the bound on I(q) carries its factor log(b/a) = log 2.
\emph{The profile on the line.} For Borel $g \ge 0$ on $\RR$, splitting at $0$,
$\int g\,d\ell = \int_0^\infty\bigl(g(r) + g(-r)\bigr)\,k(r)\,dr/r$, and for $q > 0$ the substitution $s = rq^{1/2}$ gives
\[
  \int g(rq^{1/2})\,\ell(dr) = \int_0^\infty\bigl(g(s) + g(-s)\bigr)\,k(sq^{-1/2})\,\frac{ds}{s} .
\]
In particular $\int(r^2 \wedge 1)\,\ell(dr) = 2\int_0^1rk(r)\,dr + 2\int_1^\infty k(r)\,dr/r < \infty$; and $k(r) \to 0$ as
$r \to \infty$, since $k$ is nonincreasing and $\int_1^\infty k(r)\,dr/r < \infty$.

\emph{Facts about $\psi_\ell$.} Put $\psi_\ell(\tau) := \int(1 - \cos\tau r)\,\ell(dr)$. For $\tau > 0$ the substitution $v = \tau r$ gives
$\psi_\ell(\tau) = 2\int_0^\infty(1 - \cos v)\,k(v/\tau)\,dv/v$, nondecreasing in $\tau$; $\psi_\ell$ is even, finite, as
$1 - \cos\tau r \le (2 + \tau^2)(r^2 \wedge 1)$, and continuous by dominated convergence; and $\psi_\ell(\tau) > 0$ for
$\tau \ne 0$, since for a point $r_0$ with $k(r_0) > 0$ the integrand $2(1 - \cos\tau r)k(r)/r$ is bounded below by a
positive constant on $(\tfrac\beta2,\beta)$, $\beta := \min(r_0, \pi/|\tau|)$. Also $1 - \cos 2\tau \le 4(1 - \cos\tau)$ gives
$\psi_\ell(2\tau) \le 4\psi_\ell(\tau)$, so for $0 < a \le 1$ and the $n$ with $2^{-n} < a \le 2^{1-n}$,
$\psi_\ell(\tau) \le 4^n\psi_\ell(2^{-n}\tau) \le 4a^{-2}\psi_\ell(a\tau)$: $\psi_\ell(a\tau) \ge \tfrac{a^2}4\psi_\ell(\tau)$ for $\tau \ge 0$.

\emph{The moments of $\overline\sigma$.} $\overline\sigma$ is invariant under every orthogonal map. Changing the sign of
one coordinate shows $\int u_iu_j\,d\overline\sigma = 0$ and $\int u_1^2u_iu_j\,d\overline\sigma = 0$ for $i \ne j$; exchanging
two coordinates and $\sum_ju_j^2 = 1$ give $\int u_i^2\,d\overline\sigma = \tfrac1d$, so
$\int\langle x,u\rangle^2\,\overline\sigma(du) = |x|^2/d$. Put $a := \int u_1^4\,d\overline\sigma$ and $b := \int u_1^2u_2^2\,d\overline\sigma$, which
is also $\int u_1^2u_j^2\,d\overline\sigma$ for every $j \ne 1$. A reflection $R$ with $Re_1 = (e_1 + e_2)/\sqrt2$ is
symmetric, so $\langle e_1,Ru\rangle = (u_1 + u_2)/\sqrt2$ and
$a = \int\langle e_1,Ru\rangle^4\,\overline\sigma(du) = \tfrac14\int(u_1 + u_2)^4\,d\overline\sigma = \tfrac14(2a + 6b)$, the odd terms
vanishing; so $a = 3b$, and $a + (d-1)b = \int u_1^2\,d\overline\sigma = \tfrac1d$ gives $b = \frac1{d(d+2)}$. Hence
$\int u_1^2\langle x,u\rangle^2\,\overline\sigma(du) = \frac{|x|^2 + 2x_1^2}{d(d+2)}$, and, a reflection taking a unit vector $v$
to $e_1$,
\[
  \int\langle v,u\rangle^2\langle x,u\rangle^2\,\overline\sigma(du) = \frac{|x|^2 + 2\langle v,x\rangle^2}{d(d+2)} \ge \frac{|x|^2}{d(d+2)} .
\]

\emph{Clause (1).} Write $q := q_\Sigma(u)$. For $q > 0$ the formula above with $g(r) = f(ru)$ gives
$\int f(rq^{1/2}u)\,\ell(dr) = \int_0^\infty\bigl(f(su) + f(-su)\bigr)k(sq^{-1/2})\,ds/s$; as $\overline\sigma$ and $q_\Sigma$ are
invariant under $u \mapsto -u$, the term $f(-su)$ contributes as much as $f(su)$, which gives the polar form for
$\Sigma \in \PD(d)$; its profile is nonincreasing in $s$ and at most $2k(0{+})$. $L_\Sigma$ is symmetric because $\ell$
is, and $\int(|x|^2 \wedge 1)\,dL_\Sigma = \int\overline\sigma(du)\int(qr^2 \wedge 1)\,\ell(dr) \le
\max(1,\lambda_{\max}(\Sigma))\int(r^2 \wedge 1)\,\ell(dr) < \infty$. Covariance: $q_{cQ\Sigma Q^{\tp}}(u) = c\,q_\Sigma(Q^{\tp}u)$, and
the substitution $u = Qv$, under which $\overline\sigma$ is invariant, gives
$\int f\,dL_{cQ\Sigma Q^{\tp}} = \int f(\sqrt c\,Qy)\,L_\Sigma(dy)$. Monotonicity: $\Sigma \le \Sigma'$ gives $q_\Sigma \le q_{\Sigma'}$,
directions with $q_\Sigma(u) = 0$ contribute nothing to $L_\Sigma$, and on the others the formula along a direction
with $k$ nonincreasing gives $L_\Sigma \le L_{\Sigma'}$. The difference is a measure: on each shell
$\{1/(n+2) < |x| \le 1/(n+1)\}$ and on $\{|x| > 1\}$ both measures are finite, and the differences there add up
to a measure $D$ with $D + L_\Sigma = L_{\Sigma'}$; replacing $D$ by $\tfrac12(D + (-I)_*D)$ makes it symmetric, both
measures being symmetric, and $D \le L_{\Sigma'}$ makes it a \Levy{} measure. \ledger{A3}(iii) gives an infinitely
divisible $\rho_{\Sigma,\Sigma'}$ with triplet $(2b(\Sigma' - \Sigma), D, 0)$. The \Levy{} measure being symmetric, the odd
part $i(\sin\langle\omega,x\rangle - \langle\omega,x\rangle\mathbf 1_{\{|x| \le 1\}})$ of the integrand of \ledger{A3}(i)
integrates to $0$, so its transform is $\exp[-b\langle\omega,(\Sigma' - \Sigma)\omega\rangle - \psi_{\Sigma'}(\omega) +
\psi_\Sigma(\omega)]$, with
\[
  \psi_\Sigma(\omega) := \int(1 - \cos\langle\omega,x\rangle)\,L_\Sigma(dx) = \int\overline\sigma(du)\,\psi_\ell\bigl(q_\Sigma(u)^{1/2}\langle\omega,u\rangle\bigr),
\]
and $\mu_\Sigma$ has the transform $e^{-b\langle\omega,\Sigma\omega\rangle - \psi_\Sigma(\omega)}$. These exponents telescope, which
is $\mu_{\Sigma'} = \mu_\Sigma * \rho_{\Sigma,\Sigma'}$ and (K1). (K2): on $\{\tr\Sigma < \tr\Sigma_0 + 1\}$,
$q_\Sigma(u) \le \lambda_{\max}(\Sigma) \le \tr\Sigma$ on the sphere bounds the integrand by
$\psi_\ell\bigl((\tr\Sigma_0 + 1)^{1/2}|\omega|\bigr)$, $\psi_\ell$ being even and nondecreasing on $[0,\infty)$, so
$\psi_\Sigma(\omega)$ is continuous in $\Sigma \in \PSD(d)$ by dominated convergence, $\psi_\ell$ being continuous. (K3):
covariance gives $\psi_{A\cdot\Sigma}(\omega) = \psi_\Sigma(A^{\tp}\omega)$ for every similarity $A$. (K4):
$\psi_I(e_1) = \int\overline\sigma(du)\,\psi_\ell(u_1) > 0$, because $\psi_\ell(\tau) > 0$ for $\tau \ne 0$ and
$\{u_1 = 0\}$ is $\overline\sigma$-null; so $\rho_{0,I} \ne \delta_0$.

\emph{Clause (2): the exponent.} Let $\Sigma \in \PD(d)$, $t := \lambda_{\max}(\Sigma)$ and $v$ a unit eigenvector for it. On
the sphere $q_\Sigma(u) \le t$, and $q_\Sigma(u) \ge t\langle v,u\rangle^2$: writing $u = \langle v,u\rangle v + w$ with
$w \perp v$, the cross term $\langle v,\Sigma w\rangle = t\langle v,w\rangle$ vanishes. For $R \in \Orth(d)$,
$|q_\Sigma(u)^{1/2}\langle R\omega,u\rangle| \le t^{1/2}|\omega|$, so $\psi_\Sigma(R\omega) \le \psi_\ell(t^{1/2}|\omega|)$. For
$\omega \ne 0$, $w := \omega/|\omega|$ and $a := |\langle v,u\rangle|\,|\langle w,u\rangle| \le 1$, we have
$q_\Sigma(u)^{1/2}|\langle\omega,u\rangle| \ge a\,t^{1/2}|\omega|$, so by the facts about $\psi_\ell$ and the moments of
$\overline\sigma$
\[
  \psi_\Sigma(\omega) \ge \int\overline\sigma(du)\,\psi_\ell\bigl(a\,t^{1/2}|\omega|\bigr)
  \ge \tfrac14\psi_\ell(t^{1/2}|\omega|)\int\langle v,u\rangle^2\langle w,u\rangle^2\,\overline\sigma(du)
  \ge \frac{\psi_\ell(t^{1/2}|\omega|)}{4d(d+2)} .
\]
Hence $\psi_\Sigma(R\omega) \le 4d(d+2)\,\psi_\Sigma(\omega)$ for every $R$ and $\omega$, which is the bound.

\emph{Clause (2): the moments.} For $x = rq^{1/2}u$, $|x|^2 = qr^2$ and $\langle y,x\rangle^2/|x|^2 = \langle y,u\rangle^2$, so
$\langle y,\mathcal M_s(L_\Sigma)y\rangle = \int\overline\sigma(du)\,\langle y,u\rangle^2G_s(q_\Sigma(u))$ with
$G_s(\tau) := \int(\tau r^2 \wedge s^2)\,\ell(dr)$, nondecreasing in $\tau$, with $G_s(a\tau) \ge aG_s(\tau)$ for
$0 \le a \le 1$, and finite and positive at $\tau = t$. With $t$ and $v$ as above,
$\langle v,u\rangle^2G_s(t) \le G_s(q_\Sigma(u)) \le G_s(t)$ on the sphere, and the moments of $\overline\sigma$ give
\[
  \frac{G_s(t)}{d(d+2)}\,|y|^2 \le \langle y,\mathcal M_s(L_\Sigma)y\rangle \le \frac{G_s(t)}{d}\,|y|^2 ,
\]
so $\operatorname{cond}\mathcal M_s(L_\Sigma) \le d+2$.

\emph{Clause (2): the rotational spread.} For a symmetric open set $U$ of directions, $0 < a < b$ and
$E_U := \{x : a < |x| < b,\ x/|x| \in U\}$, the formula along a direction gives
$L_\Sigma(E_U) = 2\int_U\overline\sigma(du)\,I(q_\Sigma(u))$, $I(q) := \int_a^bk(sq^{-1/2})\,ds/s$, nondecreasing in $q$. Fix $r_0$
with $k(r_0) > 0$ and take $a := r_0/4$, $b := r_0/2$; then $I(\tfrac14) > 0$, as $k(2s) \ge k(r_0)$ for $s < b$, and
$I(q) \le \log 2\cdot k(aq^{-1/2})$, $k$ being nonincreasing and $\log(b/a) = \log 2$; this bound tends to $0$ as
$q \downarrow 0$. Given $M > 0$ choose $q_2 \in (0,\tfrac14]$ with
$M\,I(q_2) < I(\tfrac14)$, let $U := \{u : u_1^2 < q_2/2,\ u_2^2 > \tfrac14\}$, a symmetric open set containing $e_2$, so that
$\overline\sigma(U) > 0$, and let $R$ exchange $e_1$ and $e_2$. For $0 < \varepsilon < q_2/2$ and unit $u$,
$q_{\Sigma_\varepsilon}(u) = \varepsilon + (1 - \varepsilon)u_1^2$, which is less than $q_2$ on $U$ and at least $u_1^2 > \tfrac14$ on
$R^{-1}U$. So $L_{\Sigma_\varepsilon}(E_U) \le 2\overline\sigma(U)I(q_2)$ while
$(R_*L_{\Sigma_\varepsilon})(E_U) = L_{\Sigma_\varepsilon}(E_{R^{-1}U}) \ge 2\overline\sigma(U)I(\tfrac14)$, and every $K$ with
$R_*L_{\Sigma_\varepsilon} \le KL_{\Sigma_\varepsilon}$ has $K > M$. So $\operatorname{rs}(L_{\Sigma_\varepsilon}) \ge M$ for
$0 < \varepsilon < q_2/2$.

\emph{Clause (2): the rank-one limit.} Let $P$ be the image of $\overline\sigma(du)\,2k(s/|u_1|)\,ds/s$ under
$(u,s) \mapsto su$. For $0 < \varepsilon \le 1$, $q_{\Sigma_\varepsilon}(u) \ge u_1^2$, so $k(s/|u_1|) \le
k(s\,q_{\Sigma_\varepsilon}(u)^{-1/2})$ off the null set $\{u_1 = 0\}$, and $P \le L_{\Sigma_\varepsilon}$ by the polar form. Hence
$P(B) \le \lim_{\varepsilon\downarrow0}L_{\Sigma_\varepsilon}(B) = L_*(B)$ for every Borel $B$ bounded away from the origin,
and $P \le L_*$, $P$ not charging the origin. Let $C$ be a nonempty open cone. It contains a unit $u^0$ with
$u^0_1 \ne 0$; put $c := |u^0_1|/2$, $V := C \cap \{u : |u_1| > c\}$, an open set containing $u^0$, so that
$\overline\sigma(V) > 0$, and $E := \{x : cr_0/2 < |x| < cr_0,\ x/|x| \in V\} \subseteq C$. On $V \cup (-V)$,
$q_{\Sigma_\varepsilon}(u) \ge u_1^2 > c^2$, so the integral $\int_{cr_0/2}^{cr_0}k(sq_{\Sigma_\varepsilon}(u)^{-1/2})\,ds/s$ is at least
$J := \int_{cr_0/2}^{cr_0}k(s/c)\,ds/s \ge k(r_0)\log 2 > 0$. By the symmetry of $L_{\Sigma_\varepsilon}$ and the formula for
symmetric sectors, $2L_{\Sigma_\varepsilon}(E) \ge L_{\Sigma_\varepsilon}(E \cup (-E)) \ge 2\overline\sigma(V)J$ for every
$\varepsilon \in (0,1]$, so $L_*(C) \ge L_*(E) \ge \overline\sigma(V)J > 0$.

\emph{Clause (3).} Let $\lambda_- < \lambda_+$ be the least and the largest eigenvalues of $\Sigma$, which differ as $\Sigma$ is
not a multiple of $I$, $u_\pm$ unit eigenvectors, $\lambda_- < m'' < m' < \lambda_+$, and $R$ the reflection exchanging $u_+$
and $u_-$. The open set $U := \{u : q_\Sigma(u) > m',\ q_\Sigma(Ru) < m''\}$ is symmetric and contains $u_+$, so
$\overline\sigma(U) > 0$, and $q_\Sigma < m''$ on $R^{-1}U$. If $L_\Sigma$ were rotation invariant, then for $0 < a < b$, with
$I_q := \int_a^bk(sq^{-1/2})\,ds/s$ and $E_U$ as above,
\[
  2\overline\sigma(U)I_{m'} \le L_\Sigma(E_U) = L_\Sigma(R^{-1}E_U) = L_\Sigma(E_{R^{-1}U}) \le 2\overline\sigma(R^{-1}U)I_{m''} = 2\overline\sigma(U)I_{m''} ,
\]
with $I_{m''} < \infty$; the integrands satisfy the reverse inequality pointwise, so $k(s/\sqrt{m'}) = k(s/\sqrt{m''})$ for
almost every $s > 0$. Put $\lambda := (m'/m'')^{1/2} > 1$ and take $r_0$ with $k(r_0) > 0$. If $k(t) \ge k(r_0)$, choose such
an $s$ with $s/\sqrt{m'} \in (\lambda^{-1/2}t, t)$; then $s/\sqrt{m''} = \lambda s/\sqrt{m'} > \lambda^{1/2}t$, and
$k(\lambda^{1/2}t) \ge k(s/\sqrt{m''}) = k(s/\sqrt{m'}) \ge k(t) \ge k(r_0)$. By induction $k(\lambda^{n/2}r_0) \ge k(r_0)$ for every
$n$, so $k \ge k(r_0) > 0$ on $(0,\infty)$, against $\int_1^\infty k(r)\,dr/r < \infty$. Finally, $\nu_\Sigma$ is infinitely
divisible, being the law \ledger{A3}(iii) gives for its triplet, and if $Q_*\nu_\Sigma = \nu_\Sigma$ for every
$Q \in \Orth(d)$, the triplet $(0, Q_*L_\Sigma, 0)$ of $Q_*\nu_\Sigma$ is that of $\nu_\Sigma$ by \ledger{A3}(ii), and $L_\Sigma$
would be rotation invariant.

Every step is machine-checked (\texttt{AffineHemigroup.SimilarityLine}); the cited interfaces are
\ledger{A3}(iii), for the increments and for the infinite divisibility of $\nu_\Sigma$, and \ledger{A3}(ii) in (3).
\end{proof}

% not in the draft: Q-0264 items 3 and 4. What survives of exact shape-blindness at index 0 for a
% general part: shape-blindness at the origin. Under infinitely divisible increments the sandwich
% along rays, sI <= Sigma <= tI, gives it at the level of Levy measures with no blow-up. Under
% positive increments without a Gaussian part the exponents are sandwiched instead, and the blow-up
% of prop:log-orbit-rays, whose measure side is one dominated convergence, finishes it in every
% direction where the difference of exponents has a limit.
% CHANGED Q-0281: \uses re-pointed to the split child it reads (eq:polar-decomposition).
% CHANGED Q-0281 wave 2: blind review W2H-4: the assignment names A3(i)-(ii), which the rewritten
% proof of (3) reads directly for the triplet of the increment; a declaration of spending only.
% CHANGED round-2 review (Gemini), item 6 (records/affine/reviews/VERIFY-round2-notation-paper.md): the
% polar profile of the datum is k_d, as in def:isotropic-origin-datum (was k_I, which read like the
% upright k_I of the delay profile). Notation only; no claim, \uses or \lean changes.
\begin{proposition}[index $0$ for a general part: shape-blind at the origin]
\label{prop:similarity-general-index-zero}
\uses{prop:similarity-general-sandwich,lem:isotropic-exponent-log,def:isotropic-origin-datum,prop:similarity-orbit-family,lem:polar-decomposition,thm:cited-A3}
\lean{AffineHemigroup.sim_index_zero_levy, AffineHemigroup.sim_index_zero_exp,
AffineHemigroup.sim_index_zero_blowup}\leanok
Let the family be as in \ref{prop:similarity-general-sandwich}, and suppose $L_I$ is the \Levy{} measure of
an isotropic origin datum $(0, k_d)$ in dimension $d$ (\ref{def:isotropic-origin-datum}) with
$c := k_d(0{+}) < \infty$. Let $\overline\sigma$ be the normalised surface measure on $S^{d-1}$, and for a Borel
set $A \subseteq S^{d-1}$ and $0 < \varepsilon < 1$ put $K_A(\varepsilon) := \{x : \varepsilon < |x| \le 1,\ x/|x| \in A\}$.
\begin{enumerate}
\item \emph{(Infinitely divisible increments: the index at the origin.)} If every Loewner step in $\PD(d)$ has
an infinitely divisible increment, then for every $\Sigma \in \PD(d)$ and every Borel $A \subseteq S^{d-1}$
\[
  \lim_{\varepsilon\downarrow0}\frac{L_\Sigma(K_A(\varepsilon))}{\log(1/\varepsilon)} = c\,\overline\sigma(A) .
\]
\item \emph{(The exponent, to within a bounded term.)} If every Loewner step has an infinitely divisible
increment, or if $C \equiv 0$ and every Loewner step has a positive finite increment, then for every
$\Sigma \in \PD(d)$ every limit point of $\psi_\Sigma(\omega) - \psi_I(\omega)$ as $|\omega| \to \infty$ lies in
$[\tfrac c2\log\lambda_{\min}(\Sigma), \tfrac c2\log\lambda_{\max}(\Sigma)]$.
\item \emph{(The blow-up.)} Under either hypothesis of (2), let $\Sigma \in \PD(d)$ and suppose
\emph{(D)}: for $\overline\sigma$-almost every $u \in S^{d-1}$ the limit
$\Delta_\Sigma(u) := \lim_{r\to\infty}\bigl(\psi_\Sigma(ru) - \psi_I(ru)\bigr)$ exists. Then $\Delta_\Sigma$ is
$\overline\sigma$-almost everywhere equal to a constant.
\end{enumerate}
\statusA\quad\ledger{A3}\quad\emph{(\textbf{Assignment.} \ledger{A3} through
\ref{prop:similarity-general-sandwich}(1)--(2), and \ledger{A3}(iii) once in (3) under the first hypothesis,
for the non-Gaussian factor of an increment, and (i)--(ii) in (3) under the first hypothesis, for the
triplet of the increment. Otherwise it cites \ref{lem:isotropic-exponent-log}(2), which
cites nothing, Fourier inversion on the Schwartz space and dominated convergence. \textbf{What it is
worth.} It is Q-0264's items~3 and~4 in the positive direction. The index-$0$ clause of conjecture A$'$,
exact shape-blindness, is false for a general part (\ref{prop:similarity-line-family}(3)); what holds is
shape-blindness \emph{at the origin}. Under infinitely divisible increments, with or without a Gaussian
part, the \Levy{} measure of every part has the isotropic member's logarithmic index $c$ at the origin,
spread uniformly over directions (clause (1)): the sandwich along rays $\lambda_{\min}I \le \Sigma \le
\lambda_{\max}I$ suffices, and no blow-up is needed. Under positive increments alone and without a Gaussian
part, the exponents are sandwiched in the same way (clause (2)), and the blow-up of
\ref{prop:log-orbit-rays} works for a general part under (D) (clause (3)): its measure side is one
dominated convergence, $\int\phi(x/\varepsilon)\,\rho(dx) \to \phi(0)\rho(\{0\})$, so no dilation-invariance argument
is needed. For an elliptical part $\nu_\Sigma = \nu_{h(\Sigma)}$ of a profile bounded at the origin, (D) holds
with $\Delta_\Sigma(u) = \tfrac c2\log(\langle u,h(\Sigma)u\rangle/h_0)$ by \ref{lem:isotropic-exponent-log}(2), and a
quadratic form constant on the sphere is a multiple of $|u|^2$; so (3) gives
\ref{prop:similarity-orbit-family}(4) again for such profiles. \textbf{Not claimed}: (D) for a general part
whose polar profiles are bounded at the origin (it asks for the second term of the exponent's expansion,
which needs more than boundedness, a Dini condition for example); the directional form of (1) under
positive increments only, which would need (3)'s conclusion to be read back as a statement on \Levy{}
measures, a Tauberian step; anything when there is a Gaussian part and the increments are only positive,
since the Gaussian term dominates the exponent at infinity.)}
\end{proposition}

\begin{proof}
\leanok
\emph{Clause (1).} Let $\lambda_- := \lambda_{\min}(\Sigma)$ and $\lambda_+ := \lambda_{\max}(\Sigma)$, so
$\lambda_-I \le \Sigma \le \lambda_+I$, and by \ref{prop:similarity-general-sandwich}(2)
$L_{\lambda_-I} \le L_\Sigma \le L_{\lambda_+I}$. For $t > 0$, $L_{tI} = (\sqrt t)_*L_I$ by
\ref{prop:similarity-general-sandwich}(1), and by \ref{eq:polar-decomposition}
\[
  L_{tI}(K_A(\varepsilon)) = L_I\bigl(t^{-1/2}K_A(\varepsilon)\bigr) = \overline\sigma(A)\int_{\varepsilon t^{-1/2}}^{t^{-1/2}}k_d(r)\,\frac{dr}{r}
  = \overline\sigma(A)\int_\varepsilon^1k_d(rt^{-1/2})\,\frac{dr}{r} .
\]
As $k_d \le c$, the last integral is at most $c\log(1/\varepsilon)$. Given $\delta > 0$, take $r_0 \in (0,1)$ with
$k_d(rt^{-1/2}) \ge c - \delta$ for $r < r_0$; for $\varepsilon < r_0$ the integral is at least $(c - \delta)\log(r_0/\varepsilon)$.
Dividing by $\log(1/\varepsilon)$, the quotient tends to $c\,\overline\sigma(A)$ for $t = \lambda_-$ and for $t = \lambda_+$, and
$L_\Sigma(K_A(\varepsilon))$ lies between the two.

\emph{Clause (2).} By \ref{prop:similarity-general-sandwich}(2) or (3), $\psi_{\lambda_-I} \le \psi_\Sigma \le
\psi_{\lambda_+I}$, and $\psi_{tI}(\omega) = \psi_I(\sqrt t\,\omega)$ by \ref{prop:similarity-general-sandwich}(1).
$\psi_I$ is the exponent $F_d$ of the datum $(0,k_d)$ (\ref{def:isotropic-origin-datum}), so
\ref{lem:isotropic-exponent-log}(2) with $a = \sqrt t$ gives $\psi_{tI}(\omega) - \psi_I(\omega) \to \tfrac c2\log t$ as
$|\omega| \to \infty$.

\emph{Clause (3).} Put $t := \lambda_{\max}(\Sigma)$, so $\Sigma \le tI$. Under the second hypothesis let $\rho$ be the
positive increment of this step; by \ref{prop:similarity-general-sandwich}(3) it is a probability measure
with $\widehat\rho = e^{-(\psi_{tI} - \psi_\Sigma)}$.
% CHANGED Q-0281 wave 2: proof rewritten to the Lean route (the difference of Levy measures read off
% the increment's triplet and symmetrized; the polar form for "full Lebesgue measure"; the test
% functions taken in the Schwartz space, so that the density of C_c^infty there is not needed).
Under the first, the infinitely divisible increment of the step has a triplet whose \Levy{} measure $\nu$
satisfies $L_{tI} = L_\Sigma + \nu$ (\ledger{A3}(i)--(ii)); its symmetrization $\nu' := \tfrac12(\nu + (-I)_*\nu)$ is a
symmetric \Levy{} measure with $\psi_{\nu'} = \psi_\nu = \psi_{tI} - \psi_\Sigma$, and $\rho$ is the law with
triplet $(0, \nu', 0)$ (\ledger{A3}(iii)), with the same transform. In both cases $0 < \widehat\rho \le 1$. For
$\omega \ne 0$ with $\omega/|\omega|$ in the set where (D) holds, which is a set of full Lebesgue measure (in polar
coordinates, Lebesgue measure on $\RR^d \setminus \{0\}$ is the product of the surface measure and $r^{d-1}\,dr$),
\[
  \widehat\rho(\omega/\varepsilon) = \exp\bigl[-\bigl(\psi_I(\sqrt t\,\omega/\varepsilon) - \psi_I(\omega/\varepsilon)\bigr)
  + \bigl(\psi_\Sigma - \psi_I\bigr)(\omega/\varepsilon)\bigr]
  \to t^{-c/2}e^{\Delta_\Sigma(\omega/|\omega|)} =: H(\omega)
  \qquad (\varepsilon \downarrow 0),
\]
by \ref{lem:isotropic-exponent-log}(2) and (D). $H$ is the almost-everywhere limit of the continuous
functions $\omega \mapsto \widehat\rho(n\omega)$, so it is measurable, and $|H| \le 1$. Let $\phi$ be a Schwartz
function and $\widehat\phi(\omega) := \int e^{i\langle\omega,x\rangle}\phi(x)\,dx$. Fourier inversion and Fubini give
$\int\phi(x/\varepsilon)\,\rho(dx) = (2\pi)^{-d}\int\widehat\phi(\omega)\,\widehat\rho(\omega/\varepsilon)\,d\omega$. As
$\varepsilon \downarrow 0$ the left-hand side tends to $\phi(0)\rho(\{0\})$ by dominated convergence, since $\phi$ is
bounded and $\phi(x/\varepsilon) \to 0$ for $x \ne 0$, a Schwartz function tending to $0$ at infinity; the
right-hand side tends to $(2\pi)^{-d}\int\widehat\phi\,H$, since $|\widehat\rho| \le 1$ and $\widehat\phi$ is integrable.
With $(2\pi)^{-d}\int\widehat\phi = \phi(0)$ this is $\int\widehat\phi\,(H - \rho(\{0\})) = 0$ for every Schwartz
function $\phi$. The Fourier transform maps the Schwartz space onto itself, so $\int g\,(H - \rho(\{0\})) = 0$
for every smooth compactly supported $g$, and $H - \rho(\{0\})$, being bounded and measurable, vanishes almost
everywhere (the fundamental lemma of the calculus of variations). So $e^{\Delta_\Sigma(\omega/|\omega|)} =
t^{c/2}\rho(\{0\})$ for almost every $\omega$, hence $\rho(\{0\}) > 0$, and $\Delta_\Sigma = \tfrac c2\log t + \log\rho(\{0\})$
at $\overline\sigma$-almost every point of the sphere, by the polar form of Lebesgue measure again.
\end{proof}

% not in the draft: the plan's item D4, its last part (notes/PLAN-article-structure.md, "What the null
% array changes"; Q-0265). The A' nodes above are stated for families of kernels from scale zero with
% no cascade, and ask that every Loewner step have an infinitely divisible increment. That hypothesis
% does not follow from their other hypotheses, so it is not removed from them: removing it would need
% a cascade hypothesis in its place, which would narrow them. For the axiomatic families it is a
% theorem, and this node says so.
% Title until Q-0284: "similarity cascades: the infinite divisibility the A' statements assume is a
% theorem of the axioms", and clause (1)'s heading "the setting of the A' nodes" (the names dropped,
% the author, 2026-10-04).
% CHANGED Q-0281: \uses re-pointed to the split child it reads (thm:cascade-infinitely-divisible (3)(a) beside (2)).
% CHANGED Q-0281 wave 2, group G: proved (AffineHemigroup.SimilarityCascade; A6 and A3(i)-(ii),
% clause (3) also A3(iii) through prop:similarity-orbit-family; guarded). The proof of record is the
% Lean's route and is unchanged. Statement unchanged.
% CHANGED Q-0469 (wording only; records/affine/GATE-affine.md, Q-0317 row "App. H 25"): clause (2)'s
% remaining hypotheses, one sentence of about seventy words after a colon, are a list; the sentence
% before it ends at "holds". The hypotheses and the claim are unchanged.
% CHANGED round-2 review (Gemini), item 6 (records/affine/reviews/VERIFY-round2-notation-paper.md): the
% polar profile of the datum is k_d, as in def:isotropic-origin-datum (was k_I, which read like the
% upright k_I of the delay profile). Notation only; no claim, \uses or \lean changes.
\begin{proposition}[similarity cascades: infinitely divisible increments are a theorem of the axioms]
\label{prop:similarity-cascade-families}
\uses{thm:cascade-infinitely-divisible,def:gw-cascade-family,prop:similarity-general-sandwich,prop:similarity-general-bounded-shape,prop:similarity-general-index-zero,prop:similarity-orbit-family,lem:path-condition-orders,thm:cited-A3,thm:cited-A6}
\lean{AffineHemigroup.sim_cascade_infinitelyDivisible, AffineHemigroup.sim_cascade_general,
AffineHemigroup.sim_cascade_elliptical}\leanok
Let $d \ge 2$, let $G$ be a closed subgroup of $\GL(d)$ containing the similarities $cQ$, $c > 0$,
$Q \in \Orth(d)$, and let $(\mu_{\Sigma,\Sigma'})$ be a $(G, \text{Loewner})$-cascade family of kernels
(\ref{def:gw-cascade-family}(2)), with kernels from scale zero $\mu_\Sigma := \mu_{0,\Sigma}$.
\begin{enumerate}
\item \emph{(Infinitely divisible kernels and increments.)} Every $\mu_{\Sigma,\Sigma'}$ is infinitely divisible. For
every $\Sigma \in \PD(d)$, $\mu_\Sigma$ is a symmetric infinitely divisible law with generating triplet
$(C(\Sigma), L_\Sigma, 0)$, where $C(\Sigma) \in \PSD(d)$ and $L_\Sigma$ is a symmetric \Levy{} measure;
$(\sqrt c\,Q)_*\mu_\Sigma = \mu_{cQ\Sigma Q^{\tp}}$ for $c > 0$ and $Q \in \Orth(d)$; and every Loewner
step $\Sigma \le \Sigma'$ in $\PD(d)$ has the infinitely divisible increment $\mu_{\Sigma,\Sigma'}$,
$\mu_{\Sigma'} = \mu_\Sigma * \mu_{\Sigma,\Sigma'}$.
\item \emph{(The general part.)} Consequently $(\mu_\Sigma)_{\Sigma \in \PD(d)}$ is a family as in
\ref{prop:similarity-general-sandwich}, and the hypothesis ``every Loewner step in $\PD(d)$ has an
infinitely divisible increment'' holds. So \ref{prop:similarity-general-sandwich}(1)--(5),
\ref{prop:similarity-general-bounded-shape}(1)--(4) and \ref{prop:similarity-general-index-zero}(1)--(3)
hold for it under their remaining hypotheses alone, which are:
\begin{itemize}
\item the comparability $L_* \ge \theta L_I$ in \ref{prop:similarity-general-bounded-shape}(2);
\item $\psi_* \ge \theta\psi_I$ in (3) there;
\item $L_I \ne 0$ in (4) there;
\item in \ref{prop:similarity-general-index-zero}, that $L_I$ is the \Levy{} measure of an isotropic
origin datum $(0, k_d)$ with $k_d(0{+}) < \infty$, with (D) in its clause (3).
\end{itemize}
\item \emph{(Elliptical parts.)} If moreover $\mu_\Sigma = N(0, C(\Sigma)) * \nu_{h(\Sigma)}$ for every
$\Sigma \in \PD(d)$, with $k$, $(\nu_H)$, $h$ and $C$ as in \ref{prop:similarity-orbit-family}, then that
proposition's clauses (2) and (3) hold under their remaining hypotheses alone (in (2): $h$ equivariant and
extending continuously to $\PSD(d) \setminus \{0\}$ with values in $\PD(d)$; in (3): $k(0{+}) < \infty$
and $h$ equivariant), with or without a Gaussian part.
\end{enumerate}
\statusA\quad\ledger{A6}\quad\ledger{A3}\quad\emph{(\textbf{Assignment.} \ledger{A6} through
\ref{thm:cascade-infinitely-divisible}(2) and (3)(a), for the infinite divisibility of every increment
and of every kernel from scale zero; \ledger{A3}(i)--(ii) once, to give $\mu_\Sigma$ a triplet and to
read its drift $0$ and the symmetry of $L_\Sigma$ off the symmetry of the law. Clauses (2) and (3) apply
the named nodes and spend what they spend. \textbf{What it is worth.} It is the plan's consequence of
the null array for the similarity row: under the axioms of the scheme at (similarities, Loewner), the
hypothesis ``infinitely divisible increments'' of
\ref{prop:similarity-orbit-family}, \ref{prop:similarity-general-sandwich},
\ref{prop:similarity-general-bounded-shape} and \ref{prop:similarity-general-index-zero} is a theorem,
and so is the standing hypothesis of the last three that every kernel from scale zero is a symmetric
infinitely divisible law with drift $0$. Conjecture A$'$ as restated (the author, 2026-10-03) therefore
holds for every axiomatic family of this corner whose rank-one limit is comparable to its isotropic
member, with or without a Gaussian part, and the handoff's open case (a), A$'$ with a Gaussian part
under positive increments only, is no question about axiomatic families. \textbf{Why the existing
nodes keep their hypothesis.} They are stated for a family of kernels from scale zero, with no
increments between other pairs and no continuity in $\Sigma$, and their hypothesis does not follow from
what else they assume. Replacing it by the cascade axioms would add hypotheses (a cascade between all
pairs, continuity) and so could only narrow them; this node adds the cascade case and leaves them as
they are. \textbf{The single step of} \texttt{records/affine/check-gaussian-matern-step.py}. The
increment computed there is positive and, its \Levy{} measure being signed, not infinitely divisible
(\ref{prop:similarity-orbit-family}'s annotation); the kernels from scale zero there have transforms
without zeros, so an increment between them is unique. By (1), those kernels are therefore the kernels
from scale zero of no $(G, \text{Loewner})$-cascade family: the step is an example of why the axioms ask
for continuity along the order, and nothing more. \textbf{Not claimed}: the converse, that a family of
the A$'$ nodes with infinitely divisible increments extends to a cascade family; anything for
similarity families that cascade along a smaller order, which is the plan's item D5: along the rays,
\ref{prop:similarity-ray-families}, and along the stable order, \ref{prop:similarity-stable-order-kernels},
whose clause (2) also gives every (similarities, Loewner)-cascade family a shape bound relative to
$\Sigma$ with no hypothesis on the rank-one limit.)}
\end{proposition}

\begin{proof}
\leanok
\emph{(1).} $(G, \text{Loewner})$ is a wedge datum: $G$ is closed and contains the dilations $cI$, and
the Loewner order is invariant under every $A \in \GL(d)$, as $A(\Sigma' - \Sigma)A^{\tp}$ is positive
semidefinite with $\Sigma' - \Sigma$, and contains the rays, as $c'\Sigma - c\Sigma = (c' - c)\Sigma$.
By \ref{thm:cascade-infinitely-divisible}(3)(a) it satisfies (P), so by
\ref{thm:cascade-infinitely-divisible}(2) every $\mu_{\Sigma,\Sigma'}$ and every
$\mu_\Sigma = \mu_{0,\Sigma}$ is infinitely divisible. Let $(C, L, \beta)$ be the generating triplet of
$\mu_\Sigma$ (\ledger{A3}(i)), with the truncation $\mathbf 1_D$, $D$ the closed unit ball. The matrix
$-I$ is a similarity and $(-I)\cdot\Sigma = \Sigma$, so (K3) gives $(-I)_*\mu_\Sigma = \mu_\Sigma$. The
transform of $(-I)_*\mu_\Sigma$ is $\widehat\mu_\Sigma(-\omega)$, and substituting $y = -x$ in the
\Levy{} integral, $D$ being symmetric, shows that it has the representation of \ledger{A3}(i) with the
triplet $(C, (-I)_*L, -\beta)$. By the uniqueness \ledger{A3}(ii), $\beta = -\beta = 0$ and
$(-I)_*L = L$. Put $C(\Sigma) := C$ and $L_\Sigma := L$. For $c > 0$ and $Q \in \Orth(d)$,
$A := \sqrt c\,Q \in G$ and $A \cdot \Sigma = cQ\Sigma Q^{\tp}$, so (K3) at the pair $(0, \Sigma)$ gives
$(\sqrt c\,Q)_*\mu_\Sigma = \mu_{cQ\Sigma Q^{\tp}}$. For $\Sigma \le \Sigma'$, (K1) at
$0 \le \Sigma \le \Sigma'$ gives $\mu_{\Sigma'} = \mu_{\Sigma,\Sigma'} * \mu_\Sigma$.

\emph{(2).} (1) is the setting of \ref{prop:similarity-general-sandwich}: symmetric infinitely divisible
laws with triplets $(C(\Sigma), L_\Sigma, 0)$, covariant under similarities in the form stated there.
Clauses (1) and (3) of that node have no hypothesis this node must supply; its clauses (2), (4) and (5), and the
named clauses of the other two nodes, ask that every Loewner step in $\PD(d)$ have an infinitely divisible
increment, which (1) supplies, together with the hypotheses listed in the statement.

\emph{(3).} The family is of the form of \ref{prop:similarity-orbit-family}, and by (1) every Loewner
step in $\PD(d)$ has an infinitely divisible increment, which is the hypothesis of that node's clauses
(2) and (3) besides those listed.
\end{proof}

% not in the draft: the plan's item D5, the cell (similarities, K_alpha-order) (Q-0267). The table's
% "K_alpha-order" is read as the order of prop:stable-step-cone(2), extended to PSD(d) by the origin and
% the rays. What writing it down found: the A' nodes do not apply to this cell, since the order has no
% step from a rank-deficient matrix to a definite one, so it supplies neither the Loewner increments nor a
% monotone approach to the rank-one kernel; two multiples of I
% sandwich every shape instead, and the bound on the kernel's shape is relative to the shape of Sigma,
% which it must be, the stable orbit belonging to the cell.
% CHANGED Q-0281: typed (the statement skeleton). Clause (1), the stable order on PSD(d), is split
% to lem:stable-order-psd, placed after this node (ADR-0025; records/formalization/PRICING.md: [T]
% matrix analysis inside an [A] node, and what prop:diagonal-stable-order-families, thm:wedge-axis and
% thm:cone-orders(4)'s reading read), and \uses gains it; the text is unchanged and the \lean tag
% names the child's declaration too.
% CHANGED Q-0281: \uses re-pointed to the split child it reads (prop:stable-step-cone for (1)(b)-(c), the order (2) and (3) only; lem:elliptical-stable (1) beside (2); thm:cascade-infinitely-divisible (3)(b)-(c) beside (2)).
% CHANGED Q-0281: blind review R5-2: the annotation's A3(iii) in (5) is assigned to
% lem:elliptical-stable(2) and this node's clause (2) (at beta = alpha, A3 through
% prop:orbit-step-cone), the route the proof of (5) takes; it named prop:stable-step-cone(1), which
% the proof reads only for the matrix analysis now in lem:stable-step-order. No claim changes.
% CHANGED Q-0281 wave 2, group G: proved (AffineHemigroup.SimilarityStableOrder; clauses (2) and (5)
% A3(iii), clauses (3)-(4) A6 and A3(i)-(ii), the isotropic members of (3) Lean core; guarded). The
% proof of (5) reads the failure of GL(d)-covariance at A = (I + e_1e_1^T)^{1/2} and omega = e_1, the
% Lean's witness, where it read diag(2, 1, ..., 1) at e_2; the rest is the Lean's route. Statement
% unchanged.
% CHANGED writing-rules pass (2026-10-14; records/affine/PASS-writing-rules-affine.md, owed
% item 1): title only, "(similarities, stable order): the families and the shape of their kernels" -> the cell of the overview table said in
% words (writing rule 1: no label of the plan in a title). Label, statement and
% \lean unchanged.
% CHANGED round-1 review (Gemini), required change 7: the statement made self-contained, wording only.
% The preamble defines psi_L, M_s, es and cond (as in prop:similarity-general-sandwich and
% prop:similarity-general-bounded-shape) instead of naming them there, and drops T_s, which no clause
% uses; clause (3) states the similarity covariance of prop:similarity-general-sandwich instead of "a
% family as in" it (the symmetric triplet was already in the clause); clause (4) reads psi_Sigma as
% psi_{L_Sigma}. The content is unchanged, and so are \uses, \lean and the proof.
% CHANGED round-2 review (Gemini), item 9, the author's decision D8 (2026-11-03;
% records/affine/PLAN-review-response-paper.md R2.0, records/affine/reviews/VERIFY-round2-structure-paper.md
% § 9, option (b)): split (ADR-0025, adr/0003). This node keeps clauses (1) and (2), the stable order on
% PSD(d) and the general linear row along it, which thm:wedge-axis spends; clauses (3)-(5), the
% similarity row, are prop:similarity-stable-order-kernels, below, numbered (1)-(3) there. The title
% says the two clauses kept; the preamble keeps what they read (the order of prop:stable-step-cone(2)
% and ecc) and gives S(d), psi_L, M_s, es and cond to the new node. The text of (1) and (2) is
% unchanged. \lean keeps the three declarations of (1) and (2); \uses keeps what (1) and (2) and their
% proof read. The annotation keeps the order's reading and the converse of (2); what it said of the
% similarity cell went with that cell. The proof keeps its parts headed (1) and (2), unchanged.
% CHANGED blind referee of the response (records/affine/REFEREE-round2-response-paper.md, item 12):
% the annotation loses its dated clause on the split ("clauses (3)--(5) of this node until the
% round-2 review split them off"), which the comment above and the changelog hold. Nothing else changes.
% CHANGED round-3 review (report E, section 3.1; the author's decision D10, PLAN-review-response-paper.md
% R3.0): the annotation's "Not claimed: a converse of (2) ... open" is replaced by a pointer to the new
% node prop:wedge-axis-profiles, which proves the converse (proved in print; its Lean is owed). The
% statement keeps its "contains" form, so that this node's Lean still covers all of it (ADR-0025); the
% proof says "the step set of prop:orbit-step-cone(2)" where it said "the cone". No claim, \uses or \lean
% changes.
% CHANGED D14 (2026-11-04; the author's decision, PLAN-review-response-paper.md R3.0): the annotation's
% pointer to prop:wedge-axis-profiles(1) drops "right-continuous", which that clause no longer assumes.
% Annotation only.
\begin{proposition}[the stable order on $\PSD(d)$ and the general linear row along it]
\label{prop:similarity-stable-order-families}
\uses{def:gw-cascade-family,lem:gw-kernels-from-origin,thm:cascade-infinitely-divisible,lem:elliptical-stable,lem:stable-step-order,prop:orbit-step-cone,prop:tempered-step-cone,prop:gw-ray-families,lem:stable-order-psd,lem:elliptical-stable-density,lem:path-condition-orders,thm:cited-A3}
\lean{AffineHemigroup.stablePlusLE_order, AffineHemigroup.stable_order_gl_row,
AffineHemigroup.stable_order_gl_row_examples}\leanok
Let $d \ge 2$ and $0 < \alpha < 2$. Write $\preceq_\alpha$ for the order of
\ref{prop:stable-step-cone}(2) on $\PD(d)$, and $\operatorname{ecc}(\Sigma) :=
\lambda_{\max}(\Sigma)/\lambda_{\min}(\Sigma)$ for $\Sigma \in \PD(d)$.
\begin{enumerate}
\item \emph{(The stable order on $\PSD(d)$.)} For $\Sigma, \Sigma' \in \PSD(d)$ put
$\Sigma \le_\alpha \Sigma'$ when $\Sigma = 0$, or $\Sigma' = c\,\Sigma$ for some $c \ge 1$, or
$\Sigma, \Sigma' \in \PD(d)$ and $\Sigma \preceq_\alpha \Sigma'$. This is a partial order on $\PSD(d)$,
invariant under $\GL(d)$, containing the rays and contained in the Loewner order, and every pair
$\Sigma \le_\alpha \Sigma'$ is joined by a chain as in \ref{thm:cascade-infinitely-divisible}(2); so
$(G, \le_\alpha)$ is a wedge datum with the path condition (P) for every closed subgroup $G$ of $\GL(d)$
containing the dilations. If $\alpha \le \beta < 2$, then $\Sigma \le_\alpha \Sigma'$ implies
$\Sigma \le_\beta \Sigma'$. For $\Sigma \in \PD(d)$ with least and largest eigenvalues
$\lambda_1 \le \lambda_d$, and $\varepsilon, t \ge 0$,
\[
  \varepsilon I \le_\alpha \Sigma \Leftrightarrow \varepsilon \le \varepsilon_\alpha(\Sigma) :=
  \Bigl(\frac{\lambda_1^{d+\alpha}}{\det\Sigma}\Bigr)^{1/\alpha} ,
  \qquad
  \Sigma \le_\alpha tI \Leftrightarrow t \ge t_\alpha(\Sigma) :=
  \Bigl(\frac{\lambda_d^{d+\alpha}}{\det\Sigma}\Bigr)^{1/\alpha} ,
\]
and $t_\alpha(\Sigma)/\varepsilon_\alpha(\Sigma) = \operatorname{ecc}(\Sigma)^{(d+\alpha)/\alpha}$.
\item \emph{(The general linear row.)} Let $k$ be a profile as in \ref{prop:orbit-step-cone} whose
dilation ratio has $R(\lambda) \ge \lambda^\alpha$ for every $\lambda \ge 1$, and $\mu_I$ the law of that
proposition's clause (1). Then $\Sigma \mapsto (\Sigma^{1/2})_*\mu_I$, $\Sigma \in \PSD(d)$, is the family of
kernels from scale zero of a $(\GL(d), \le_\alpha)$-cascade family of kernels. This applies to the stable
law $S_\beta(I)$ of every index $\beta \in [\alpha, 2)$ (the profile $k(r) = c_{d,\beta}^{-1}r^{-\beta}$ of
\ref{lem:elliptical-stable}), and to the profiles of \ref{prop:tempered-step-cone} of every index
$\beta \in [\alpha, 2)$.
\end{enumerate}
\statusA\quad\ledger{A3}\quad\emph{(\textbf{Assignment.} \ledger{A3}(iii) through
\ref{prop:orbit-step-cone} in (2), for the laws of the profiles and their increments. Clause (1) spends
nothing: it is \ref{lem:stable-order-psd}, matrix analysis. \textbf{Reading.} The plan's
``$\mathfrak k_\alpha$-order'' is $\preceq_\alpha$, whose relative steps $M$ are the exponentials of the symmetric
matrices of the cone $\mathfrak k_\alpha$ (\ref{thm:cascade-infinitely-divisible}(3)(c)); it is the whole order generated by
the cone field $\mathfrak k_\alpha$ (\ref{prop:stable-order-generated}(1), the plan's item D3), and nothing here depends
on that. The order $\le_\alpha$ of (1) is $\preceq^+_{C_\alpha}$ of \ref{thm:cone-orders}(4), and the general
linear row along it is \ref{thm:wedge-axis}(3). On $\PSD(d)$ the order is extended by the origin and the rays, which the scheme asks for, and by nothing
else: a nonzero rank-deficient matrix is below no definite one, as for the stable orbit, whose
rank-deficient kernels have \Levy{} measures carried by a proper subspace. \textbf{What it is worth.}
Clause (1) is the order \ref{thm:wedge-axis} reads on $\PSD(d)$, with the two thresholds that sandwich
every definite matrix between multiples of $I$, and clause (2) is the general linear cell along it. The
similarity row along the same order is \ref{prop:similarity-stable-order-kernels}. \textbf{The converse of
(2)} holds for every profile, with any Gaussian part, and is
\ref{prop:wedge-axis-profiles}(1): the increments of a $(\GL(d), \le_\alpha)$-cascade family are infinitely
divisible (\ref{thm:cascade-infinitely-divisible}(2)), so they order the \Levy{} densities
(\ref{prop:similarity-orbit-family}(1)), and \ref{prop:orbit-step-cone}(2) along the steps of the proof of
\ref{lem:dilation-ratio}(4) reads off $\alpha \le \beta_-(k)$. It
is a separate node because the Lean of this one proves sufficiency only. \textbf{Not claimed} here: the
converse; and, for a single step not embedded in a cascade, that a positive increment orders the densities,
which is the open converse of \ref{prop:orbit-step-cone}.)}
\end{proposition}

\begin{proof}
\leanok
\emph{(1), the order.} On $\PD(d)$ the relation is $\preceq_\alpha$: a definite matrix is not $0$, and
the second case is contained in the third, every ray $\Sigma \le c\,\Sigma$, $c \ge 1$, being a step of
$\preceq_\alpha$ by \ref{lem:stable-step-order}(3) with $\lambda = \sqrt c$. If $\Sigma$ is definite, every
$\Sigma'$ with $\Sigma \le_\alpha \Sigma'$ is definite. If $\Sigma$ is singular and not $0$, those $\Sigma'$
are its multiples $c\Sigma$, $c \ge 1$. The properties a wedge datum asks of its order follow, whatever
$G$ is.
\begin{itemize}
\item Reflexivity is the case $c = 1$.
\item Antisymmetry: the relation is contained in the Loewner order, as $\Sigma' - 0 \ge 0$,
$c\Sigma - \Sigma = (c-1)\Sigma \ge 0$, and $\preceq_\alpha$ is contained in it
(\ref{lem:stable-step-order}(1)).
\item Transitivity: from $\Sigma = 0$ there is nothing to prove, from a nonzero singular $\Sigma$ the
relation runs along its multiples, and on $\PD(d)$ it is the transitivity of $\preceq_\alpha$
(\ref{lem:stable-step-order}(2)), the rays being steps of $\preceq_\alpha$.
\item Invariance: $A\cdot 0 = 0$, $A\cdot(c\Sigma) = c\,A\cdot\Sigma$, and $\GL(d)$ preserves $\PD(d)$ and
$\preceq_\alpha$ (\ref{lem:stable-step-order}(2)).
\item The rays: for $0 \le c \le c'$, $c\Sigma \le_\alpha c'\Sigma$ by the first case if $c = 0$ and by the
second, with the factor $c'/c$, otherwise.
\end{itemize}
For $\alpha \le \beta < 2$ only the third case changes: with $M := \Sigma^{-1/2}\Sigma'\Sigma^{-1/2}$,
$\Sigma \preceq_\alpha \Sigma'$ gives $\lambda_{\min}(M) \ge 1$ (\ref{lem:stable-step-order}(1)), so
$\lambda_{\min}(M)^{d+\beta} \ge \lambda_{\min}(M)^{d+\alpha} \ge \det M$.

\emph{(1), the chains.} From $0$ to $\Sigma'$ take $\gamma(s) := s\Sigma'$, and from $\Sigma \ne 0$ to
$c\Sigma$ take $\gamma(s) := (1 + s(c-1))\Sigma$; along both, $\gamma(s) \le_\alpha \gamma(t)$ for $s \le t$
by the rays. For $\Sigma \preceq_\alpha \Sigma'$ in $\PD(d)$ write $M = U\diag(m_1, \dots, m_d)U^{\tp}$ with
$U \in \Orth(d)$, put $M^s := U\diag(m_1^s, \dots, m_d^s)U^{\tp}$, and take the chain of
\ref{lem:path-condition-orders}(c), $\gamma(s) := \Sigma^{1/2}M^s\Sigma^{1/2}$. It is continuous and
definite, with $\gamma(0) = \Sigma$ and $\gamma(1) = \Sigma'$. Its points are ordered by the criterion (b) of
\ref{lem:stable-step-order}(1). For $x = \Sigma^{1/2}y \ne 0$ put $z := U^{\tp}y$; then
\[
  \langle x, \gamma(s)^{-1}x\rangle = \sum_{i} m_i^{-s}z_i^2 =: Q_s , \qquad
  \det\gamma(s) = \det\Sigma\,(\det M)^s ,
\]
so that $n_{\gamma(s)}(x) = c_{d,\alpha}^{-1}(\det\Sigma)^{-1/2}(\det M)^{-s/2}Q_s^{-(d+\alpha)/2}$. Let
$s \le t$ and $q := \lambda_{\min}(M)$. Every $m_i$ is at least $q$, so $m_i^{-t} \le q^{-(t-s)}m_i^{-s}$ and
$Q_t \le q^{-(t-s)}Q_s$. And $\det M \le q^{d+\alpha}$ gives $(\det M)^{(t-s)/2} \le
q^{(t-s)(d+\alpha)/2}$. Together,
$(\det M)^{-s/2}Q_s^{-(d+\alpha)/2} \le (\det M)^{-t/2}Q_t^{-(d+\alpha)/2}$, which is
$n_{\gamma(s)}(x) \le n_{\gamma(t)}(x)$. So $\gamma(s) \preceq_\alpha \gamma(t)$.

\emph{(1), the thresholds.} The value $\varepsilon = 0$ is the first case, and $\Sigma \le_\alpha 0$ is
impossible for $\Sigma \ne 0$. For $\varepsilon > 0$ and $t > 0$ the two relations are steps of
$\preceq_\alpha$, with relative steps $M = \Sigma/\varepsilon$ and $M = t\Sigma^{-1}$, so that
\[
  \lambda_{\min}(\Sigma/\varepsilon) = \lambda_1/\varepsilon , \quad \det(\Sigma/\varepsilon) = \det\Sigma/\varepsilon^d ,
  \qquad
  \lambda_{\min}(t\Sigma^{-1}) = t/\lambda_d , \quad \det(t\Sigma^{-1}) = t^d/\det\Sigma .
\]
Condition (c) of \ref{lem:stable-step-order}(1) then reads
\[
  \lambda_1^{d+\alpha} \ge \varepsilon^\alpha\det\Sigma
  \qquad\text{and}\qquad
  t^\alpha \ge \lambda_d^{d+\alpha}/\det\Sigma ,
\]
which are the two equivalences. The ratio is $((\lambda_d/\lambda_1)^{d+\alpha})^{1/\alpha}$.

\emph{(2).} By \ref{prop:orbit-step-cone}(1), $\mu_I$ is rotation invariant with
$\widehat\mu_I = e^{-\psi_I} > 0$, and as in the proof of \ref{prop:gw-ray-families}(3),
$\mu_\Sigma := (\Sigma^{1/2})_*\mu_I$ has $\widehat\mu_\Sigma(\omega) = F(\langle\omega,\Sigma\omega\rangle)$ with
$F(u) := \widehat\mu_I(\sqrt u\,e_1)$ continuous and positive. So the transforms have no zeros, and (a)--(c)
of \ref{lem:gw-kernels-from-origin}(2) hold for $G = \GL(d)$: (a) as $F(0) = 1$, (b) as
$\widehat{A_*\mu_\Sigma}(\omega) = F(\langle\omega,A\Sigma A^{\tp}\omega\rangle)$, and (c) as $F$ is
continuous; $\mu_I \ne \delta_0$ because
% CHANGED Q-0281 wave 2: blind review W2G-2
$\psi_I$ is not identically $0$, $k$ being positive near $0$. By
\ref{lem:gw-kernels-from-origin}(3) it remains to give each pair $\Sigma \le_\alpha \Sigma'$ an increment.
From $0$: $\mu_{\Sigma'}$. Along a ray, $\Sigma' = c\Sigma$ with $c > 1$: by \ref{prop:orbit-step-cone}(2)--(3)
the ray $I \le cI$ has an increment, $\mu_{cI} = \mu_I * \rho$, and applying $(\Sigma^{1/2})_*$, with
$\Sigma^{1/2}\sqrt c = (c\Sigma)^{1/2}$, gives $\mu_{c\Sigma} = \mu_\Sigma * (\Sigma^{1/2})_*\rho$. On
$\PD(d)$, with $M := \Sigma^{-1/2}\Sigma'\Sigma^{-1/2}$ and $q := \lambda_{\min}(M) \ge 1$, the inequality
$\det M \le q^{d+\alpha}$, condition (c) of \ref{lem:stable-step-order}(1), and $R(q^{1/2})^2 \ge q^\alpha$
give $\det M \le q^dR(q^{1/2})^2$, the step set of
\ref{prop:orbit-step-cone}(2), which supplies the increment. For the stable profile
$k(r) = c_{d,\beta}^{-1}r^{-\beta}$, $k(|x|)|x|^{-d}$ is $n_I(x)$ of \ref{eq:stable-levy-density} at index
$\beta$, so $\mu_I = S_\beta(I)$ by \ref{lem:elliptical-stable}(1)--(2), and $R(\lambda) = \lambda^\beta \ge
\lambda^\alpha$ by \ref{prop:orbit-step-cone}(3); for the profiles of \ref{prop:tempered-step-cone}(1) of index
$\beta$, $R(\lambda) = \lambda^\beta$ by that clause.
\end{proof}

% CHANGED round-2 review (Gemini), item 9, the author's decision D8 (2026-11-03;
% records/affine/PLAN-review-response-paper.md R2.0, records/affine/reviews/VERIFY-round2-structure-paper.md
% § 9, option (b)): split from prop:similarity-stable-order-families (ADR-0025, adr/0003), whose clauses
% (3), (4) and (5) are this node's (1), (2) and (3), their text unchanged but for the clause numbers and
% the pointers to the parent's (1) and (2). The title is the parent's former one. The preamble restates
% S(d), psi_L, M_s, es and cond as the parent did, and reads the order and its thresholds from the
% parent's (1). \lean is the five declarations of the three clauses, which name the parent's clauses in
% their docstrings (Formalization/AffineHemigroup/SimilarityStableOrder.lean: (3) -> (1), (4) -> (2),
% (5) -> (3)). \uses is what the three clauses and their proof read, with the parent for the order. The
% annotation is the parent's, the part about the similarity cell; the proof is the parent's parts
% headed (3), (4) and (5), renumbered, with no step changed. Printed in the paper's section 4.9.
% CHANGED blind referee of the response (records/affine/REFEREE-round2-response-paper.md, items 11
% and 12): \uses made to match what the clauses and the proof cite. It gains prop:stable-step-cone,
% whose (1)(b) and (3) the proof of (3) reads, and drops lem:stable-step-order and
% lem:elliptical-stable-density, which nothing in the node cites. The annotation loses its dated
% sentence on the split, which this comment above and the changelog hold. Statement and proof unchanged.
\begin{proposition}[similarity covariance along the stable order: the families and the shape of their kernels]
\label{prop:similarity-stable-order-kernels}
\uses{def:gw-cascade-family,lem:gw-kernels-from-origin,lem:affine-gauge,thm:cascade-infinitely-divisible,lem:elliptical-stable,lem:stable-tail,prop:stable-step-cone,prop:similarity-general-sandwich,prop:similarity-general-bounded-shape,prop:similarity-cascade-families,prop:similarity-ray-families,prop:similarity-stable-order-families,thm:cited-A3,thm:cited-A6}
\lean{AffineHemigroup.stable_order_sim_kernels, AffineHemigroup.stable_order_sim_isotropic,
AffineHemigroup.stable_order_sandwich, AffineHemigroup.loewner_order_sandwich,
AffineHemigroup.stable_order_cell_strict}\leanok
Let $d \ge 2$ and $0 < \alpha < 2$, and let $\le_\alpha$ be the stable order on $\PSD(d)$ and
$\varepsilon_\alpha(\Sigma)$ and $t_\alpha(\Sigma)$ its thresholds, as in
\ref{prop:similarity-stable-order-families}(1). Write $\mathrm S(d)$ for the similarities $cQ$, $c > 0$,
$Q \in \Orth(d)$, and $\operatorname{ecc}(\Sigma) := \lambda_{\max}(\Sigma)/\lambda_{\min}(\Sigma)$ for
$\Sigma \in \PD(d)$. For a measure $L$ on $\RR^d \setminus \{0\}$ and $s > 0$ put
\[
  \psi_L(\omega) := \int(1 - \cos\langle\omega,x\rangle)\,L(dx) ,
  \qquad
  \mathcal M_s(L) := \int(|x|^2 \wedge s^2)\,\frac{xx^{\tp}}{|x|^2}\,L(dx) ;
\]
for a function $\psi \ge 0$ on $\RR^d$ let the \emph{exponent spread} $\operatorname{es}(\psi)$ be the
infimum of the $K \ge 1$ with $\psi(Q\omega) \le K\,\psi(\omega)$ for every $Q \in \Orth(d)$ and
$\omega \in \RR^d$, the infimum of the empty set being $\infty$; and put
$\operatorname{cond}M := \lambda_{\max}(M)/\lambda_{\min}(M) \in [1,\infty]$ for
$M \in \PSD(d) \setminus \{0\}$. These are the functionals of
\ref{prop:similarity-general-bounded-shape}.
\begin{enumerate}
\item \emph{(The kernels of the similarity row.)} Let $G$ be a closed subgroup of $\GL(d)$ with
$\mathrm S(d) \subseteq G$, and $(\mu_{\Sigma,\Sigma'})$ a $(G, \le_\alpha)$-cascade family of kernels.
Every $\mu_{\Sigma,\Sigma'}$ and every kernel from scale zero $\mu_\Sigma$ is infinitely divisible; for
$\Sigma \in \PD(d)$, $\mu_\Sigma$ is a symmetric law with generating triplet $(C(\Sigma), L_\Sigma, 0)$;
the family $(\mu_\Sigma)_{\Sigma \in \PD(d)}$ is covariant under the similarities,
$(\sqrt c\,Q)_*\mu_\Sigma = \mu_{cQ\Sigma Q^{\tp}}$ for $c > 0$, $Q \in \Orth(d)$ and $\Sigma \in \PD(d)$
(so it satisfies the hypotheses of \ref{prop:similarity-general-sandwich}); and every
step $\Sigma \le_\alpha \Sigma'$ in $\PD(d)$ has the infinitely divisible increment $\mu_{\Sigma,\Sigma'}$.
For $G = \mathrm S(d)$ the isotropic members $\mu_I$ of these families are exactly the rotation-invariant
self-decomposable laws other than $\delta_0$.
\item \emph{(The sandwich, and the shape.)} In the setting of (1) let $\Sigma \in \PD(d)$,
$\varepsilon := \varepsilon_\alpha(\Sigma)$ and $t := t_\alpha(\Sigma)$, and write
$\psi_\Sigma := \psi_{L_\Sigma}$. Then $C(I) = c_0I$ with $c_0 \ge 0$, and
\[
  \varepsilon c_0I \le C(\Sigma) \le tc_0I , \qquad
  (\sqrt\varepsilon\,I)_*L_I \le L_\Sigma \le (\sqrt t\,I)_*L_I , \qquad
  \psi_I(\sqrt\varepsilon\,\omega) \le \psi_\Sigma(\omega) \le \psi_I(\sqrt t\,\omega) .
\]
Consequently $\operatorname{cond}C(\Sigma) \le \operatorname{ecc}(\Sigma)^{(d+\alpha)/\alpha}$ if $c_0 > 0$,
$\operatorname{cond}\mathcal M_s(L_\Sigma) \le \operatorname{ecc}(\Sigma)^{(d+\alpha)/\alpha}$ for every $s > 0$ if
$L_I \ne 0$, and $\operatorname{es}(\psi_\Sigma) \le 4\operatorname{ecc}(\Sigma)^{(d+\alpha)/\alpha}$. For a
$(G, \text{Loewner})$-cascade family of kernels the same holds with $\varepsilon := \lambda_{\min}(\Sigma)$,
$t := \lambda_{\max}(\Sigma)$ and $\operatorname{ecc}(\Sigma)$ in place of
$\operatorname{ecc}(\Sigma)^{(d+\alpha)/\alpha}$.
\item \emph{(The cell is not the union of its neighbours.)} Let $b > 0$ and $\kappa > 0$, and for
$\Sigma \in \PSD(d)$ put $S_\alpha(\Sigma) := (\Sigma^{1/2})_*S_\alpha(I)$ and
\[
  \mu_\Sigma := N\bigl(0,\, 2b(\Sigma + \kappa\tr(\Sigma)\,I)\bigr) * S_\alpha(\Sigma) .
\]
Then $(\mu_\Sigma)$ is the family of kernels from scale zero of an $(\mathrm S(d), \le_\alpha)$-cascade
family of kernels, and it is not covariant under $\GL(d)$, hence the kernels of no
$(\GL(d), \le_\alpha)$-family (\ref{lem:gw-kernels-from-origin}(2)). For $P \in \PSD(d)$ with $0 < \rank P < d$ there
is no positive finite measure $\rho$ with $\mu_{I+P} = \mu_I * \rho$; so $(\mu_\Sigma)$ is the family of
kernels from scale zero of no $(\mathrm S(d), \text{Loewner})$-cascade family.
\end{enumerate}
\statusA\quad\ledger{A6}\quad\ledger{A3}\quad\emph{(\textbf{Assignment.} \ledger{A6} through
\ref{thm:cascade-infinitely-divisible}(2) in (1), for the infinite divisibility of the increments and of the
kernels, and so in (2); \ledger{A3}(i)--(ii) in (1), as in \ref{prop:similarity-cascade-families}(1), to
give a kernel its triplet and read its symmetry, and through \ref{prop:similarity-general-sandwich} in (2);
\ledger{A3}(iii) through \ref{lem:elliptical-stable}(2) and
\ref{prop:similarity-stable-order-families}(2) in (3). The negative half of (3) spends nothing beyond the
existence of the stable law; \ref{lem:stable-tail} cites nothing. The order and its thresholds are read
from \ref{prop:similarity-stable-order-families}(1), which spends nothing. \textbf{What writing it down
found.} The A$'$ nodes (\ref{prop:similarity-general-sandwich} to
\ref{prop:similarity-general-index-zero}) ask that every Loewner step have an infinitely divisible
increment and read a monotone rank-one limit; $\le_\alpha$ does not supply them (the kernel at
$e_1e_1^{\tp}$ exists, and is a limit by (K2), but the shapes approaching it are not
$\le_\alpha$-comparable), and a family of the cell need not satisfy them (the stable orbit). Clause (2) replaces the
rank-one sandwich by two multiples of $I$, and the bound it gives is relative: the kernel's anisotropy is
at most a power of the parameter's. A uniform bound is impossible in this cell, since the stable orbit
$S_\alpha(\Sigma)$ belongs to it by \ref{prop:similarity-stable-order-families}(2) and has exponent spread $\operatorname{ecc}(\Sigma)^{\alpha/2}$
(\ref{prop:similarity-general-bounded-shape}(1)). The Loewner reading of (2) holds for every
(similarities, Loewner)-cascade family with no hypothesis on the rank-one limit, and so complements
\ref{prop:similarity-degenerate-family}, whose spread $\operatorname{ecc}(\Sigma)^{p\alpha/2}$ is unbounded
and within it. \textbf{What it is worth.} The cell (similarities, $\mathfrak k_\alpha$) of the plan's table is not
empty, and it strictly contains its two neighbours, the $(\GL(d), \le_\alpha)$ families of
\ref{prop:similarity-stable-order-families}(2) and the
restrictions of the (similarities, Loewner) families (\ref{lem:gw-kernels-from-origin}(1), the order being
contained in the Loewner order): the stable orbit is in the first and not in the second
(\ref{prop:stable-step-cone}(3), as in the proof of (3)); the Gaussian family with covariances
$2b(\Sigma + \kappa\tr(\Sigma)I)$ is in the second, its covariances growing by $P + \kappa\tr(P)I \ge 0$
along every Loewner step $P$, and not in the first (the proof of (3) without the stable factor); and the family of (3) is in neither. Its isotropic members are those of every similarity
column (\ref{prop:similarity-ray-families}(4)). \textbf{Not claimed}: that the exponent
$(d+\alpha)/\alpha$ is sharp (the stable orbit has $\alpha/2$ in the exponent spread); and a
classification of the cell.)}
\end{proposition}

\begin{proof}
\leanok
\emph{(1).} By \ref{prop:similarity-stable-order-families}(1) and \ref{thm:cascade-infinitely-divisible}(2), every $\mu_{\Sigma,\Sigma'}$ and every
$\mu_\Sigma = \mu_{0,\Sigma}$ is infinitely divisible. The rest is the proof of
\ref{prop:similarity-cascade-families}(1), which uses of the Loewner order only that it gives these
increments: $-I \in G$ and $(-I)\cdot\Sigma = \Sigma$, so (K3) makes $\mu_\Sigma$ symmetric, and the uniqueness
\ledger{A3}(ii) makes its drift $0$ and its \Levy{} measure symmetric; (K3) at $(0, \Sigma)$ with
$A = \sqrt c\,Q$ is the covariance $(\sqrt c\,Q)_*\mu_\Sigma = \mu_{cQ\Sigma Q^{\tp}}$; and (K1) at
$0 \le_\alpha \Sigma \le_\alpha \Sigma'$ gives $\mu_{\Sigma'} = \mu_\Sigma * \mu_{\Sigma,\Sigma'}$. The last sentence
is \ref{prop:similarity-ray-families}(4) with $W' = \le_\alpha$, which by \ref{prop:similarity-stable-order-families}(1) is $\mathrm S(d)$-invariant,
contains the rays and is contained in the Loewner order.

\emph{(2).} \ref{prop:similarity-general-sandwich}(1) with $c = 1$ and $\Sigma = I$ gives $C(I) =
Q\,C(I)\,Q^{\tp}$ for every $Q \in \Orth(d)$, so $C(I) = c_0I$ by the argument of the proof of
\ref{lem:affine-gauge}, and $c_0 \ge 0$; with $Q = I$ and $c = s > 0$ it gives $C(sI) = sc_0I$,
$L_{sI} = (\sqrt s\,I)_*L_I$ and $\psi_{sI}(\omega) = \psi_I(\sqrt s\,\omega)$, and $L_I$ and $\psi_I$ are
rotation invariant. By \ref{prop:similarity-stable-order-families}(1), $\varepsilon I \le_\alpha \Sigma \le_\alpha tI$, with $\varepsilon > 0$, and by (1) both steps have infinitely divisible increments, so \ref{prop:similarity-general-sandwich}(2) gives
$C(\varepsilon I) \le C(\Sigma) \le C(tI)$, $L_{\varepsilon I} \le L_\Sigma \le L_{tI}$ and
$\psi_{\varepsilon I} \le \psi_\Sigma \le \psi_{tI}$, which is the display. The eigenvalues of $C(\Sigma)$ lie
in $[\varepsilon c_0, tc_0]$, so $\operatorname{cond}C(\Sigma) \le t/\varepsilon$.

\emph{The moment matrices.} $\mathcal M_s$ is monotone in the \Levy{} measure, $\mathcal M_s((aI)_*L) =
a^2\mathcal M_{s/a}(L)$, and $\mathcal M_r(L_I) = d^{-1}T_r(L_I)\,I$, as in the proof of
\ref{prop:similarity-general-bounded-shape}(1) and (4). Hence
\[
  \varepsilon\,d^{-1}T_{s/\sqrt\varepsilon}(L_I)\,I \le \mathcal M_s(L_\Sigma) \le t\,d^{-1}T_{s/\sqrt t}(L_I)\,I .
\]
$T_r(L_I) = \int(|x|^2 \wedge r^2)\,L_I(dx)$ is nondecreasing in $r$ and positive when $L_I \ne 0$, and
$s/\sqrt\varepsilon \ge s/\sqrt t$; so $\operatorname{cond}\mathcal M_s(L_\Sigma) \le tT_{s/\sqrt t}(L_I)/
(\varepsilon T_{s/\sqrt\varepsilon}(L_I)) \le t/\varepsilon$.

\emph{The exponent spread.} Put $g(u) := \psi_I(\sqrt u\,e_1)$, so $\psi_I(\omega) = g(|\omega|^2)$ by rotation
invariance. For $c \ge 1$ the ray step $I \le_\alpha cI$ has an infinitely divisible increment by (1), and
\ref{prop:similarity-general-sandwich}(2) gives $\psi_I(\omega) \le \psi_{cI}(\omega) = \psi_I(\sqrt c\,\omega)$,
so $g$ is nondecreasing. As $1 - \cos 2\tau = 2(1 - \cos\tau)(1 + \cos\tau) \le 4(1 - \cos\tau)$,
integrating at $\tau = \langle \omega,x\rangle$ against $L_I(dx)$ gives $\psi_I(2\omega) \le 4\psi_I(\omega)$, that is
$g(4u) \le 4g(u)$. For $c \ge 1$ let $n \ge 0$ be least with $4^n \ge c$,
so that $4^n < 4c$; then $g(cu) \le g(4^nu) \le 4^ng(u) \le 4c\,g(u)$. For $R \in \Orth(d)$ and
$\omega \in \RR^d$,
\[
  \psi_\Sigma(R\omega) \le \psi_I(\sqrt t\,R\omega) = g(t|\omega|^2) \le 4\frac{t}{\varepsilon}\,g(\varepsilon|\omega|^2)
  = 4\frac{t}{\varepsilon}\,\psi_I(\sqrt\varepsilon\,\omega) \le 4\frac{t}{\varepsilon}\,\psi_\Sigma(\omega) ,
\]
so $\operatorname{es}(\psi_\Sigma) \le 4t/\varepsilon$, and $t/\varepsilon = \operatorname{ecc}(\Sigma)^{(d+\alpha)/\alpha}$ by \ref{prop:similarity-stable-order-families}(1). \emph{The Loewner order.} $\lambda_{\min}(\Sigma)I \le \Sigma \le \lambda_{\max}(\Sigma)I$, every
Loewner step of a $(G, \text{Loewner})$-family has an infinitely divisible increment
(\ref{prop:similarity-cascade-families}(1)), and the rays are Loewner steps; the same argument applies with
$t/\varepsilon = \operatorname{ecc}(\Sigma)$.

\emph{(3), the family.} Write $C(\Sigma) := \Sigma + \kappa\tr(\Sigma)I$. By \ref{lem:elliptical-stable}(2),
$\widehat{S_\alpha(\Sigma)}(\omega) = e^{-|\Sigma^{1/2}\omega|^\alpha} = e^{-\langle\omega,\Sigma\omega\rangle^{\alpha/2}}$,
so $\widehat\mu_\Sigma(\omega) = \exp[-b\langle\omega,C(\Sigma)\omega\rangle - \langle\omega,\Sigma\omega\rangle^{\alpha/2}]$,
positive and continuous in $\Sigma$, which is (c) of \ref{lem:gw-kernels-from-origin}(2) with no zeros. (a):
$C(0) = 0$, so $\mu_0 = \delta_0$, and $\widehat\mu_I$ is not identically $1$. (b): for $A = \sqrt c\,Q$,
$A\,C(\Sigma)A^{\tp} = cQ\Sigma Q^{\tp} + \kappa c\tr(\Sigma)I = C(A\cdot\Sigma)$, as $\tr(A\cdot\Sigma) = c\tr\Sigma$,
and $\langle A^{\tp}\omega,\Sigma A^{\tp}\omega\rangle = \langle\omega,(A\cdot\Sigma)\omega\rangle$; so
$\widehat\mu_\Sigma(A^{\tp}\omega) = \widehat\mu_{A\cdot\Sigma}(\omega)$. (d): let $\Sigma \le_\alpha \Sigma'$. By \ref{prop:similarity-stable-order-families}(2) at $\beta = \alpha$ there is a probability measure $\rho^{\mathrm s}$ with $S_\alpha(\Sigma') = S_\alpha(\Sigma) *
\rho^{\mathrm s}$; by \ref{prop:similarity-stable-order-families}(1), $P :=  \Sigma' - \Sigma \ge 0$, so $C(\Sigma') - C(\Sigma) = P + \kappa\tr(P)I \ge 0$, and
$\rho := N(0, 2b(C(\Sigma') - C(\Sigma))) * \rho^{\mathrm s}$ has $\mu_{\Sigma'} = \mu_\Sigma * \rho$.
\ref{lem:gw-kernels-from-origin}(3) gives the family. It is not $\GL(d)$-covariant: put
$M := I + e_1e_1^{\tp}$ and $A := M^{1/2} \in \GL(d)$, so that $A\cdot I = M$ and $\tr M = d + 1$. The
transforms of $A_*\mu_I$ and of $\mu_{A\cdot I} = \mu_M$ at $\omega$ are
$\exp[-b(1 + d\kappa)\langle\omega,M\omega\rangle - \langle\omega,M\omega\rangle^{\alpha/2}]$ and
$\exp[-b\langle\omega,(M + \kappa(d+1)I)\omega\rangle - \langle\omega,M\omega\rangle^{\alpha/2}]$; at $\omega = e_1$,
where $\langle\omega,M\omega\rangle = 2$, the exponents differ by $b\kappa(d - 1) \ne 0$.

\emph{(3), no rank-deficient step.} Suppose $\mu_{I+P} = \mu_I * \rho$ with $\rho$ a positive finite measure.
As $\widehat\mu_I > 0$, $\widehat\rho(\omega) = \exp[-\langle\omega,G\omega\rangle - \psi_{I+P}(\omega) + \psi_I(\omega)]$ with
$G := b(P + \kappa\tr(P)I)$ and $\psi_H(\omega) = \langle\omega,H\omega\rangle^{\alpha/2}$, which is
\ref{eq:stable-tail-transform} with $H = I$ and $H' = I + P$. \ref{lem:stable-tail} gives $n_I \le n_{I+P}$,
which is condition (b) of \ref{prop:stable-step-cone}(1) for the pair $(I, I + P)$, and that node's clause (3)
excludes it, $P$ having rank strictly between $0$ and $d$. An $(\mathrm S(d), \text{Loewner})$-family with these
kernels from scale zero would give $\mu_{I+P} = \mu_{I,I+P} * \mu_I$ by (K1) at $0 \le I \le I + P$.
\end{proof}

% CHANGED Q-0281: split from prop:similarity-stable-order-families (ADR-0025;
% records/formalization/PRICING.md): clause (1), copied verbatim with the parent's clause number. It
% is matrix analysis and spends nothing, where the parent's other clauses spend A6 and A3; the parent
% \uses this node, and its text and numbering are unchanged.
% CHANGED Q-0281 wave 1: the \statusT note updated to the new proof route (blind review W1D-1); the
% statement is unchanged.
\begin{lemma}[the stable order on $\PSD(d)$]
\label{lem:stable-order-psd}
\uses{def:gw-cascade-family,lem:stable-step-order,lem:path-condition-orders}
\lean{AffineHemigroup.stablePlusLE_order}\leanok
Let $d \ge 2$ and $0 < \alpha < 2$. Write $\preceq_\alpha$ for the order of
\ref{lem:stable-step-order}(2) on $\PD(d)$ (the order of \ref{prop:stable-step-cone}(2)) and
$\operatorname{ecc}(\Sigma) :=
\lambda_{\max}(\Sigma)/\lambda_{\min}(\Sigma)$ for $\Sigma \in \PD(d)$.
\begin{enumerate}
\item[(1)] \emph{(The stable order on $\PSD(d)$.)} For $\Sigma, \Sigma' \in \PSD(d)$ put
$\Sigma \le_\alpha \Sigma'$ when $\Sigma = 0$, or $\Sigma' = c\,\Sigma$ for some $c \ge 1$, or
$\Sigma, \Sigma' \in \PD(d)$ and $\Sigma \preceq_\alpha \Sigma'$. This is a partial order on $\PSD(d)$,
invariant under $\GL(d)$, containing the rays and contained in the Loewner order, and every pair
$\Sigma \le_\alpha \Sigma'$ is joined by a chain as in \ref{thm:cascade-infinitely-divisible}(2); so
$(G, \le_\alpha)$ is a wedge datum with the path condition (P) for every closed subgroup $G$ of $\GL(d)$
containing the dilations. If $\alpha \le \beta < 2$, then $\Sigma \le_\alpha \Sigma'$ implies
$\Sigma \le_\beta \Sigma'$. For $\Sigma \in \PD(d)$ with least and largest eigenvalues
$\lambda_1 \le \lambda_d$, and $\varepsilon, t \ge 0$,
\[
  \varepsilon I \le_\alpha \Sigma \Leftrightarrow \varepsilon \le \varepsilon_\alpha(\Sigma) :=
  \Bigl(\frac{\lambda_1^{d+\alpha}}{\det\Sigma}\Bigr)^{1/\alpha} ,
  \qquad
  \Sigma \le_\alpha tI \Leftrightarrow t \ge t_\alpha(\Sigma) :=
  \Bigl(\frac{\lambda_d^{d+\alpha}}{\det\Sigma}\Bigr)^{1/\alpha} ,
\]
and $t_\alpha(\Sigma)/\varepsilon_\alpha(\Sigma) = \operatorname{ecc}(\Sigma)^{(d+\alpha)/\alpha}$.
\end{enumerate}
\statusT\quad\emph{(Cites nothing beyond the spectral theorem: the order of
\ref{prop:stable-step-cone}(2) is read through \ref{lem:stable-step-order}, and the chains are those of
\ref{lem:path-condition-orders}, whose monotonicity on $\PD(d)$ the proof checks through the densities.
Clause (1) of \ref{prop:similarity-stable-order-families}, split off on
2026-10-04 (Q-0281, ADR-0025); the numbering stays the parent's.)}
\end{lemma}

% CHANGED Q-0281 wave 1: proof rewritten to the Lean route. It is the part of the proof of
% prop:similarity-stable-order-families headed (1), with two changes of route: every appeal to
% prop:stable-step-cone is to the split child lem:stable-step-order, and the chain on PD(d) is checked
% to be monotone directly, through the densities, instead of being read from
% lem:path-condition-orders(c), which is not machine-checked; the chain itself is the same.
\begin{proof}
\leanok
\emph{The order.} Reflexivity is the case $c = 1$. On $\PD(d)$ the relation is $\preceq_\alpha$: a
definite matrix is not $0$, and the second case is contained in the third, every ray $\Sigma \le
c\,\Sigma$, $c \ge 1$, being a step of $\preceq_\alpha$ by \ref{lem:stable-step-order}(3) with
$\lambda = \sqrt c$. The relation is contained in the Loewner order, as $\Sigma' - 0 \ge 0$,
$c\Sigma - \Sigma = (c-1)\Sigma \ge 0$, and $\preceq_\alpha$ is contained in it by
\ref{lem:stable-step-order}(1); so it is antisymmetric. For transitivity let
$\Sigma \le_\alpha \Sigma' \le_\alpha \Sigma''$. If $\Sigma = 0$ there is nothing to prove. If
$\Sigma' = c\,\Sigma$, then either $\Sigma'' = c'\Sigma' = c'c\,\Sigma$ with $c'c \ge 1$, or $\Sigma'$ and
$\Sigma''$ are definite with $\Sigma' \preceq_\alpha \Sigma''$; in the second case $\Sigma = c^{-1}\Sigma'$ is
definite, $\Sigma \preceq_\alpha c\,\Sigma$ is a ray, and $\preceq_\alpha$ is transitive
(\ref{lem:stable-step-order}(2)). If $\Sigma \preceq_\alpha \Sigma'$ in $\PD(d)$, then $\Sigma'' = c'\Sigma'$,
a ray of $\preceq_\alpha$, or $\Sigma' \preceq_\alpha \Sigma''$, and again $\preceq_\alpha$ is transitive.
Invariance: $A\cdot 0 = 0$, $A\cdot(c\Sigma) = c\,A\cdot\Sigma$, and $\GL(d)$ preserves $\PD(d)$ and
$\preceq_\alpha$ (\ref{lem:stable-step-order}(2)). The rays: for $0 \le c \le c'$, $c\Sigma \le_\alpha c'\Sigma$
by the first case if $c = 0$ and by the second, with the factor $c'/c$, otherwise. These are the
properties a wedge datum asks of the order, whatever $G$ is. For $\alpha \le \beta < 2$ only the third
case changes: with $M := \Sigma^{-1/2}\Sigma'\Sigma^{-1/2}$, $\Sigma \preceq_\alpha \Sigma'$ gives
$\lambda_{\min}(M) \ge 1$ (the proof of \ref{lem:stable-step-order}(1)), so
$\lambda_{\min}(M)^{d+\beta} \ge \lambda_{\min}(M)^{d+\alpha} \ge \det M$.

\emph{The chains.} From $0$ to $\Sigma'$ take $\gamma(s) := s\Sigma'$, and from $\Sigma \ne 0$ to $c\Sigma$
take $\gamma(s) := (1 + s(c-1))\Sigma$; along both, $\gamma(s) \le_\alpha \gamma(t)$ for $s \le t$ by the
rays. For $\Sigma \preceq_\alpha \Sigma'$ in $\PD(d)$ write $M = U\diag(m_1, \dots, m_d)U^{\tp}$ with
$U \in \Orth(d)$, put $M^s := U\diag(m_1^s, \dots, m_d^s)U^{\tp}$, and take the chain of
\ref{lem:path-condition-orders}(c), $\gamma(s) := \Sigma^{1/2}M^s\Sigma^{1/2}$. It is continuous, as every
$s \mapsto m_i^s$ is, and definite, as every $m_i^s > 0$, with $\gamma(0) = \Sigma$ and $\gamma(1) = \Sigma'$. Its points are ordered by the criterion (b) of
\ref{lem:stable-step-order}(1). For $x = \Sigma^{1/2}y \ne 0$ put $z := U^{\tp}y$; then
\[
  \langle x, \gamma(s)^{-1}x\rangle = \sum_{i} m_i^{-s}z_i^2 =: Q_s , \qquad
  \det\gamma(s) = \det\Sigma\,(\det M)^s ,
\]
so that $n_{\gamma(s)}(x) = c_{d,\alpha}^{-1}(\det\Sigma)^{-1/2}(\det M)^{-s/2}Q_s^{-(d+\alpha)/2}$. Let
$s \le t$ and $m := \lambda_{\min}(M)$. Every $m_i$ is at least $m$, so $m_i^{-t} \le m^{-(t-s)}m_i^{-s}$ and
$Q_t \le m^{-(t-s)}Q_s$; and $\det M \le m^{d+\alpha}$ gives $(\det M)^{(t-s)/2} \le
m^{(t-s)(d+\alpha)/2}$. Together,
$(\det M)^{-s/2}Q_s^{-(d+\alpha)/2} \le (\det M)^{-t/2}Q_t^{-(d+\alpha)/2}$, which is
$n_{\gamma(s)}(x) \le n_{\gamma(t)}(x)$; so $\gamma(s) \preceq_\alpha \gamma(t)$.

\emph{The thresholds.} $\varepsilon = 0$ is the first case, and $\Sigma \le_\alpha 0$ is impossible for
$\Sigma \ne 0$; so at $t = 0$ both sides of the second equivalence fail, as $t_\alpha(\Sigma) > 0$.
For $\varepsilon > 0$ the relation $\varepsilon I \le_\alpha \Sigma$ is
$\varepsilon I \preceq_\alpha \Sigma$, whose relative step $M = \Sigma/\varepsilon$ has
$\lambda_{\min}(M) = \lambda_1/\varepsilon$ and $\det M = \det\Sigma/\varepsilon^d$, so that condition (c) of
\ref{lem:stable-step-order}(1) reads $\lambda_1^{d+\alpha} \ge \varepsilon^\alpha\det\Sigma$. For $t > 0$,
$\Sigma \preceq_\alpha tI$ has $M = t\Sigma^{-1}$, $\lambda_{\min}(M) = t/\lambda_d$ (the least eigenvalue of
$\Sigma^{-1}$ is $1/\lambda_d$) and $\det M = t^d/\det\Sigma$, and (c) reads
$t^\alpha \ge \lambda_d^{d+\alpha}/\det\Sigma$. The ratio is $((\lambda_d/\lambda_1)^{d+\alpha})^{1/\alpha}$.
\end{proof}

% not in the draft: Q-0282 item 3, the cell (diagonal, stable order) of the overview table, which the
% draft rewrite (Q-0279) left marked "not described"; in the style of prop:diagonal-ray-families. What
% writing it down found: on the frame's shapes the stable order's steps are the diagonal exponents of
% the cone K_alpha, so the isotropic member must be B-self-decomposable for each of them; the axis
% steps that force products along the Loewner order are absent (they are rank-deficient), so products
% are not forced, and the condition is strictly stronger than self-decomposability (a strongly
% correlated Gaussian fails it). Off the frame's shapes the order, unlike the rays, links different
% correlation orbits, and the coordinate map Delta that builds the product families does not obviously
% respect it; no description off the frame is attempted.
% CHANGED Q-0281: \uses re-pointed to the split child it reads (condition (c) of prop:stable-step-cone only; lem:parameter-space(5) only; prop:similarity-stable-order-families (1) beside (2)).
% CHANGED writing-rules pass (2026-10-14; records/affine/PASS-writing-rules-affine.md, owed
% item 1): title only, "(diagonal, stable order): the families on the frame's shapes" -> the cell of the overview table said in
% words (writing rule 1: no label of the plan in a title). Label, statement and
% \lean unchanged.
% CHANGED round-2 review (Gemini), item 12, the author's decision D7 (2026-11-03): the exponent cone K_alpha
% is written frak k_alpha (a subset of gl(d), in fraktur as gl(d) is), apart from the step cone
% cal K_alpha. Notation only; no claim, \uses or \lean changes.
\begin{proposition}[diagonal covariance along the stable order: the families on the frame's shapes]
\label{prop:diagonal-stable-order-families}
\uses{def:gw-cascade-family,lem:gw-kernels-from-origin,lem:parameter-space-examples,lem:elliptical-stable,lem:stable-step-order,prop:stable-exponent-cone,prop:orbit-step-cone,prop:diagonal-ray-families,prop:similarity-stable-order-families,lem:stable-order-psd,thm:cited-A3}
\lean{AffineHemigroup.diagonal_stable_datum, AffineHemigroup.diagonal_ray_parameter,
AffineHemigroup.diagonal_stable_frame, AffineHemigroup.diagonal_stable_sources,
AffineHemigroup.diagonal_stable_gaussian_excluded, AffineHemigroup.diagonal_stable_not_product,
AffineHemigroup.diagonal_stable_mixed}\leanok
Let $d \ge 2$ and $0 < \alpha < 2$. Let $\mathrm D_+(d)$, $D_s$, $\Delta(\Sigma)$ and $\mathrm{Corr}(d)$ be as
in \ref{prop:diagonal-ray-families}, let $\le_\alpha$ be the stable order on $\PSD(d)$ of
\ref{prop:similarity-stable-order-families}(1), and use $B$-self-decomposability as in
\ref{prop:stable-exponent-cone}. Write $\mathcal D_\alpha$ for the set of diagonal matrices $\diag(b)$,
$b \in \RR^d$, with $(d+\alpha)\min_ib_i \ge b_1 + \dots + b_d$; these are the diagonal matrices of the
cone $\mathfrak k_\alpha$ of \ref{prop:stable-exponent-cone}.
\begin{enumerate}
\item \emph{(The datum.)} $(\mathrm D_+(d), \le_\alpha)$ is a wedge datum, and its orbits in $\PD(d)$ are
those of \ref{prop:diagonal-ray-families}(1), indexed by $\mathrm{Corr}(d)$. On the orbit of $I$,
$D_s^2 \le_\alpha D_{s'}^2$ if and only if $u := (s'_1/s_1, \dots, s'_d/s_d)$ satisfies
$(\min_iu_i)^{d+\alpha} \ge u_1\cdots u_d$.
\item \emph{(On the frame's shapes.)} The kernels from scale zero of a
$(\mathrm D_+(d), \le_\alpha)$-cascade family of kernels are, on the orbit of $I$,
$\mu_{D_s^2} = (D_s)_*\mu_I$, $s \in (0,\infty)^d$, where $\mu_I \ne \delta_0$ is $B$-self-decomposable for
every $B \in \mathcal D_\alpha$; in particular $\mu_I$ is self-decomposable.
\item \emph{(Two sources of families.)} (a) For self-decomposable laws $\nu_1, \dots, \nu_d$ on $\RR$, not
all $\delta_0$, the map $\Sigma \mapsto (\Delta(\Sigma)^{1/2})_*(\nu_1 \otimes \dots \otimes \nu_d)$ is the
family of kernels from scale zero of a $(\mathrm D_+(d), \le_\alpha)$-cascade family of kernels.
(b) For a profile $k$ as in \ref{prop:orbit-step-cone} with $R(\lambda) \ge \lambda^\alpha$ for every
$\lambda \ge 1$, and $\mu_I$ the law of that proposition's clause (1), so is $\Sigma \mapsto (\Sigma^{1/2})_*\mu_I$;
in particular for $\mu_I = S_\beta(I)$, $\alpha \le \beta < 2$.
\item \emph{(Strictly between its neighbours.)} (a) Let $|r| < 1$ with $r^2 > 4(1+\alpha)/(2+\alpha)^2$, and
let $\Theta_r \in \mathrm{Corr}(d)$ have the entries $(\Theta_r)_{12} = (\Theta_r)_{21} = r$ and agree with $I$ elsewhere.
The centred Gaussian law $N(0, \Theta_r)$ is self-decomposable and is not $B$-self-decomposable for
$B := I + \alpha\,e_1e_1^{\tp} \in \mathcal D_\alpha$. So the $(\mathrm D_+(d), \text{rays})$-family of
\ref{prop:diagonal-ray-families}(2) with $\mu_I = N(0, \Theta_r)$ is the restriction of no
$(\mathrm D_+(d), \le_\alpha)$-family.
(b) $S_\alpha(I)$ is not of the form $\nu_1 \otimes \dots \otimes \nu_d$. So the family of (3)(b) with
$\mu_I = S_\alpha(I)$ is the restriction of no $(\mathrm D_+(d), \text{Loewner})$-family.
(c) Let $\sigma_\alpha$ be the law on $\RR$ with transform $e^{-|\tau|^\alpha}$ (\ref{lem:elliptical-stable}(2)
at $d = 1$). The map $\Sigma \mapsto (\Sigma^{1/2})_*S_\alpha(I) * (\Delta(\Sigma)^{1/2})_*(\sigma_\alpha
\otimes \dots \otimes \sigma_\alpha)$ is the family of kernels from scale zero of a
$(\mathrm D_+(d), \le_\alpha)$-cascade family that is the restriction neither of a
$(\GL(d), \le_\alpha)$-family nor of a $(\mathrm D_+(d), \text{Loewner})$-family.
\end{enumerate}
\statusA\quad\ledger{A3}\quad\emph{(\textbf{Assignment.} \ledger{A3}(iii) for the existence of the stable
laws and of the profile laws and their increments, through \ref{lem:elliptical-stable}(2) and
\ref{prop:similarity-stable-order-families}(2), in (3)(b) and (4)(b)--(c). Clauses (1), (2), (3)(a) and
(4)(a) spend nothing. \ledger{A6} would make every kernel infinitely divisible, the order having the path
condition (\ref{prop:similarity-stable-order-families}(1)); the node does not use it. \textbf{What it is
worth.} The cell (diagonal, stable order) of the overview table is not empty, and on the frame's shapes it
lies strictly between its two neighbours in the diagonal row: every restricted
$(\mathrm D_+(d), \text{Loewner})$-family belongs to it (\ref{lem:gw-kernels-from-origin}(1), as
$\le_\alpha$ is contained in the Loewner order), and every family of it restricts to the ray cell
(\ref{lem:gw-kernels-from-origin}(1) and \ref{prop:diagonal-ray-families}); by (4)(a)--(b) neither inclusion is an equality, and by (4)(c) the cell
is not the union of the restrictions of the $(\mathrm D_+(d), \text{Loewner})$ cell and the
$(\GL(d), \le_\alpha)$ cell. The condition on the isotropic member, $B$-self-decomposability for the diagonal
exponents of $\mathfrak k_\alpha$, sits between self-decomposability, which the ray cell asks for, and the product form,
which the Loewner order's axis steps force (\ref{prop:fixed-frame-families}); those axis steps are
rank-deficient and so are not steps of $\le_\alpha$. \textbf{Not claimed}: the converse of (2), that every
$B$-self-decomposable law for $\mathcal D_\alpha$ is the isotropic member of a family; a description of the
families off the frame's shapes, where $\le_\alpha$, unlike the ray order, relates matrices of different
correlation orbits (\ref{prop:stable-step-cone}(3)); and the group of all invertible diagonal matrices.)}
\end{proposition}

\begin{proof}
\leanok
\emph{(1).} $\mathrm D_+(d)$ is a closed subgroup containing the dilations
(\ref{prop:diagonal-ray-families}(1)), so \ref{prop:similarity-stable-order-families}(1) makes
$(\mathrm D_+(d), \le_\alpha)$ a wedge datum. The orbits depend on the group alone. For the frame:
$D_s^2$ and $D_{s'}^2$ are positive definite, so $D_s^2 \le_\alpha D_{s'}^2$ is $D_s^2 \preceq_\alpha D_{s'}^2$
(the second case of the definition is contained in the third on $\PD(d)$). Its relative step is
$M = D_s^{-1}D_{s'}^2D_s^{-1} = D_u^2$, with $\lambda_{\min}(M) = (\min_iu_i)^2$ and $\det M = (u_1\cdots u_d)^2$,
and condition (c) of \ref{prop:stable-step-cone}(1) is the stated inequality squared.

\emph{(2).} (K3) at the pair $(0, I)$ with $A = D_s$ gives $\mu_{D_s^2} = (D_s)_*\mu_I$, and (K4) gives
$\mu_I \ne \delta_0$. Let $B = \diag(b) \in \mathcal D_\alpha$ and $t > 0$, and put $u_i := e^{tb_i}$, so that
$D_u = e^{tB}$. Then $(\min_iu_i)^{d+\alpha} = e^{t(d+\alpha)\min_ib_i} \ge e^{t(b_1 + \dots + b_d)} =
u_1\cdots u_d$, so $I \le_\alpha D_u^2$ by (1) with $s = \mathbf 1$. (K1) at $0 \le_\alpha I \le_\alpha D_u^2$
gives $(e^{tB})_*\mu_I = \mu_{D_u^2} = \mu_{I,D_u^2} * \mu_I$, and applying $(e^{-tB})_*$,
$\mu_I = (e^{-tB})_*\mu_I * (e^{-tB})_*\mu_{I,D_u^2}$, which is $B$-self-decomposability. $B = I$ lies in
$\mathcal D_\alpha$, as $d + \alpha \ge d$, and $I$-self-decomposability is \ref{eq:self-decomposable} with
$c = e^{-t}$.

\emph{(3).} (a) By \ref{prop:diagonal-ray-families}(3) the map is the family of kernels from scale zero of a
$(\mathrm D_+(d), \text{Loewner})$-cascade family; $\le_\alpha$ is contained in the Loewner order
(\ref{prop:similarity-stable-order-families}(1)), and \ref{lem:gw-kernels-from-origin}(1) restricts it to
$(\mathrm D_+(d), \le_\alpha)$ with the same kernels from scale zero. (b) By
\ref{prop:similarity-stable-order-families}(2) the map is the family of kernels from scale zero of a
$(\GL(d), \le_\alpha)$-cascade family, the stable laws $S_\beta(I)$ included, and
\ref{lem:gw-kernels-from-origin}(1) restricts it to $\mathrm D_+(d) \subseteq \GL(d)$.

\emph{(4)(a).} $\Theta_r$ is positive definite, as $|r| < 1$, with unit diagonal, so $\Theta_r \in \mathrm{Corr}(d)$,
and $N(0, \Theta_r)$ is self-decomposable by \ref{prop:diagonal-ray-families}(4). $B$ has diagonal
$(1+\alpha, 1, \dots, 1)$, and $(d+\alpha)\cdot 1 = (1 + \alpha) + (d - 1)$, so $B \in \mathcal D_\alpha$.
Suppose $N(0, \Theta_r) = (e^{-tB})_*N(0, \Theta_r) * \rho_t$ with probability measures $\rho_t$, $t > 0$. The
transform of $(e^{-tB})_*N(0, \Theta_r)$ is $e^{-\langle\omega,e^{-tB}\Theta_re^{-tB}\omega\rangle/2}$, positive, so
$\widehat\rho_t(\omega) = e^{-\langle\omega,C_t\omega\rangle/2}$ with $C_t := \Theta_r - e^{-tB}\Theta_re^{-tB}$. As
$|\widehat\rho_t| \le 1$, $\langle\omega,C_t\omega\rangle \ge 0$ for every $\omega$ and $t > 0$.
% CHANGED Q-0281 wave 2: proof rewritten to the Lean route (the form of B Theta + Theta B read at
% its negative direction, in place of the determinant of its 2x2 block).
Take $\omega := xe_1 + e_2$ with $x := -r(2+\alpha)/(2(1+\alpha))$. Then
$f(t) := \langle e^{-tB}\omega, \Theta_r e^{-tB}\omega\rangle = x^2e^{-2(1+\alpha)t} + e^{-2t} + 2rxe^{-(2+\alpha)t}$
is largest at $t = 0$ on $[0, \infty)$, so $f'(0) \le 0$. But
$f'(0) = -2(1+\alpha)x^2 - 2 - 2(2+\alpha)rx = (2+\alpha)^2r^2/(2(1+\alpha)) - 2$, which is positive by the
choice of $r$ (which is possible, as $4(1+\alpha) < (2+\alpha)^2$); this is the quadratic form of
$-(B\Theta_r + \Theta_rB)$ at $\omega$, whose $2 \times 2$ block in the rows and columns $1, 2$ has the negative
determinant $4(1+\alpha) - (2+\alpha)^2r^2$. So $N(0, \Theta_r)$ is not $B$-self-decomposable. A
$(\mathrm D_+(d), \le_\alpha)$-family restricting to the ray family of $N(0, \Theta_r)$ would have the same
kernels from scale zero (\ref{lem:gw-kernels-from-origin}(1)), hence the isotropic member $N(0, \Theta_r)$,
against (2).

\emph{(4)(b).} The transform of $S_\alpha(I)$ is $e^{-|\omega|^\alpha}$ (\ref{lem:elliptical-stable}(2)). If
$S_\alpha(I) = \nu_1 \otimes \dots \otimes \nu_d$, then $\widehat\nu_i(\tau) = e^{-|\tau|^\alpha}$, the transform at
$\tau e_i$, and $e^{-|\omega|^\alpha} = \prod_ie^{-|\omega_i|^\alpha}$; at $\omega = e_1 + e_2$ this says
$2^{\alpha/2} = 2$, false for $\alpha < 2$. By \ref{prop:diagonal-ray-families}(3) at $s = \mathbf 1$, the
isotropic member of a $(\mathrm D_+(d), \text{Loewner})$-family is such a product; a restriction has the
same isotropic member (\ref{lem:gw-kernels-from-origin}(1)).

\emph{(4)(c), the family.} $\sigma_\alpha$ is self-decomposable: for $0 < c < 1$,
$e^{-|\tau|^\alpha} = e^{-|c\tau|^\alpha}\,e^{-(1-c^\alpha)|\tau|^\alpha}$, and the second factor is the
transform of the image of $\sigma_\alpha$ under $\tau \mapsto (1 - c^\alpha)^{1/\alpha}\tau$. So (3)(a) with
$\nu_i = \sigma_\alpha$ and (3)(b) with $\beta = \alpha$ give two maps $\Sigma \mapsto \mu^1_\Sigma$ and
$\Sigma \mapsto \mu^2_\Sigma$, each the kernels from scale zero of a $(\mathrm D_+(d), \le_\alpha)$-family. Their
transforms, $\widehat\mu^1_\Sigma(\omega) = e^{-\sum_i\Sigma_{ii}^{\alpha/2}|\omega_i|^\alpha}$ and
$\widehat\mu^2_\Sigma(\omega) = e^{-\langle\omega,\Sigma\omega\rangle^{\alpha/2}}$, have no zeros. Put
$\mu_\Sigma := \mu^1_\Sigma * \mu^2_\Sigma$ and check (a)--(d) of \ref{lem:gw-kernels-from-origin}(2). (a):
$\mu_0 = \delta_0$, and $\widehat\mu_I(e_1) = e^{-2} \ne 1$, so $\mu_I \ne \delta_0$. (b): the image of a
convolution under a linear map is the convolution of the images, and each factor is covariant. (c): the
transform is a product of two continuous ones. (d): for $\Sigma \le_\alpha \Sigma'$,
$\mu^j_{\Sigma'} = \mu^j_\Sigma * \mu^j_{\Sigma,\Sigma'}$ for $j = 1, 2$, so
$\mu_{\Sigma'} = \mu_\Sigma * (\mu^1_{\Sigma,\Sigma'} * \mu^2_{\Sigma,\Sigma'})$. The transforms have no
zeros, so \ref{lem:gw-kernels-from-origin}(3) gives the family.

\emph{(4)(c), neither restriction.} $\widehat\mu_I(\omega) = e^{-|\omega|^\alpha - \sum_i|\omega_i|^\alpha}$. At
$\omega = e_1$ it is $e^{-2}$, and at $\omega = (e_1 + e_2)/\sqrt2$, of the same length, it is
$e^{-1 - 2^{1-\alpha/2}}$, which differs as $2^{1-\alpha/2} > 1$; so $\mu_I$ is not rotation invariant. The
isotropic member of a $(\GL(d), \le_\alpha)$-family is rotation invariant (\ref{lem:parameter-space}(5),
$\Orth(d) \subseteq \GL(d)$), and a restriction keeps it. If $\mu_I$ were a product
$\nu_1 \otimes \dots \otimes \nu_d$, then $\widehat\nu_i(\tau) = \widehat\mu_I(\tau e_i) = e^{-2|\tau|^\alpha}$, and
$e^{-|\omega|^\alpha - \sum_i|\omega_i|^\alpha} = e^{-2\sum_i|\omega_i|^\alpha}$ would give
$e^{-|\omega|^\alpha} = e^{-\sum_i|\omega_i|^\alpha}$, false at $\omega = e_1 + e_2$ as in (4)(b); by
\ref{prop:diagonal-ray-families}(3) it is therefore the isotropic member of no
$(\mathrm D_+(d), \text{Loewner})$-family.
\end{proof}
